\documentclass[11pt,reqno]{amsart}
\usepackage[T1]{fontenc}
\usepackage{lmodern,amsmath,amscd,amssymb,amsthm,mathtools,mathrsfs,microtype,booktabs,longtable}
\usepackage[margin=1.08in]{geometry}
\usepackage{xr-hyper}
\usepackage[hidelinks]{hyperref}
\numberwithin{equation}{section}
\newtheorem{theorem}{Theorem}[section]
\newtheorem{proposition}[theorem]{Proposition}
\newtheorem{lemma}[theorem]{Lemma}
\newtheorem{corollary}[theorem]{Corollary}
\theoremstyle{definition}\newtheorem{definition}[theorem]{Definition}
\newtheorem{example}[theorem]{Example}
\theoremstyle{remark}\newtheorem{remark}[theorem]{Remark}
\newcommand{\C}{\mathbb C}\newcommand{\Q}{\mathbb Q}
\newcommand{\Pj}{\mathbb P}\newcommand{\A}{\mathbb A}
\newcommand{\OO}{\mathcal O}\newcommand{\II}{\mathcal I}
\newcommand{\NN}{\mathcal N}\newcommand{\LL}{\mathcal L}
\newcommand{\ZZ}{\mathcal Z}\newcommand{\RR}{\mathcal R}
\DeclareMathOperator{\Sym}{Sym}\DeclareMathOperator{\Gr}{Gr}
\DeclareMathOperator{\Hom}{Hom}\DeclareMathOperator{\Tot}{Tot}
\DeclareMathOperator{\Fitt}{Fitt}\DeclareMathOperator{\Ann}{Ann}
\DeclareMathOperator{\rank}{rank}\DeclareMathOperator{\Proj}{Proj}
\DeclareMathOperator{\Spec}{Spec}\DeclareMathOperator{\PGL}{PGL}
\DeclareMathOperator{\Pic}{Pic}\DeclareMathOperator{\diag}{diag}
\DeclareMathOperator{\Tor}{Tor}\DeclareMathOperator{\im}{im}
\DeclareMathOperator{\coker}{coker}\DeclareMathOperator{\codim}{codim}
\DeclareMathOperator{\End}{End}
\DeclareMathOperator{\disc}{disc}\DeclareMathOperator{\Disc}{Disc}
\DeclareMathOperator{\Supp}{Supp}\DeclareMathOperator{\Hilb}{Hilb}
\DeclareMathOperator{\rad}{rad}\DeclareMathOperator{\Ass}{Ass}
\DeclareMathOperator{\gr}{gr}\DeclareMathOperator{\Gal}{Gal}
\DeclareMathOperator{\length}{length}
\allowdisplaybreaks[2]
\setlength{\emergencystretch}{3em}

% Embedded companion-paper references; no external aux file is required.
\makeatletter
\@namedef{r@cor:finite-flat-family}{{I.7.5}{I.18}{}{}{}}
\@namedef{r@eq:fixed-spectral-locus}{{I.L.1}{I.59}{}{}{}}
\@namedef{r@eq:general-coefficient-map}{{I.6.6}{I.13}{}{}{}}
\@namedef{r@lem:linear-preserver-v147}{{I.4.1}{I.6}{}{}{}}
\@namedef{r@lem:pencil-irreducibility}{{I.6.2}{I.13}{}{}{}}
\@namedef{r@lem:primitive-pencil}{{I.E.1}{I.43}{}{}{}}
\@namedef{r@prop:universal-envelope-v159}{{I.11.1}{I.25}{}{}{}}
\@namedef{r@sec:closed-pencil-v151}{{I.G}{I.47}{}{}{}}
\@namedef{r@sec:rigidification-v151}{{I.F}{I.46}{}{}{}}
\@namedef{r@thm:artin-local-inverse-v146}{{I.8.3}{I.20}{}{}{}}
\@namedef{r@thm:closed-pencil-v151}{{I.G.1}{I.47}{}{}{}}
\@namedef{r@thm:failure-incidence-v160}{{I.12.1}{I.26}{}{}{}}
\@namedef{r@thm:fixed-spectral-classification}{{I.L.2}{I.60}{}{}{}}
\@namedef{r@thm:intrinsic-envelope-v158}{{I.10.1}{I.24}{}{}{}}
\@namedef{r@thm:recognition-dichotomy-v148}{{I.B.1}{I.34}{}{}{}}
\makeatother

\begin{document}
\title{Intrinsic power geometry and the boundary of quadratic pencils}
\author{Qian Qi}
\date{September 25, 2026}
\subjclass[2020]{14M12, 14B05, 14D20, 13C40}
\hypersetup{pdftitle={Intrinsic power geometry and the boundary of quadratic pencils, revision 160},pdfauthor={Qian Qi}}
\begin{abstract}
Multiplication in a determinant-apolar algebra realizes balanced
products of symmetric minor ideals over every complex base algebra.
Its all-rank power graphs recover relative complete quadrics, and
power contacts together with finite syzygy maps recover the complete
data of arbitrary symmetric pencils. We identify the normalized
Hilbert incidence space, on the simple corank-two open in every
dimension, with the blow-up of a smooth codimension-two centre.
This determines the whole exceptional projective line, its nodal
universal family, its discrepancy, and all multiplication systems,
including residual directions tangent to the determinant conic.
The centre and its Rees algebra are intrinsic multiplication data
of the sharp failure family. For simultaneous ordinary collisions
of arbitrary multiplicities we retain the full nodal-tree calculation,
its exact power ideals and its residual-pencil moduli.
\end{abstract}
\maketitle
\section{Introduction}
\label{sec:paper-ii-v159}
A compactification records more than a limiting pencil when different
smoothing directions produce different resolved curves. We study this
difference through multiplication powers in the finite algebra
\[
 B_h(V)=\Sym(\Sym^2V)/\langle(v^2)^{h+1}:v\in V\rangle,
 \qquad h\geq1.
\]
Here $V$ is a complex vector space of dimension $n$. For a symmetric
matrix $A$, write $J_p(A)$ for the ideal generated by the coefficients
of its $p$th power. Its universal identity is
\[
 J_{hq+s}(A)=I_q(A)^{h-s}I_{q+1}(A)^s,\qquad 0\leq s<h,
\]
where $I_q$ denotes the ideal of $q$-minors
(Theorem~\ref{thm:universal-power-v157}). This is equality of ordinary
ideals, including over nonreduced bases. The proof identifies the
coefficient representation and uses classical invariant-ideal
balancing. We make its classical inputs and its precise apolar
interpretation explicit in Section~\ref{sec:power-conventions-v158}.

\subsection{The full simple-incidence modification}
Let $\mathcal U\subset\Gr(2,\Sym^2V)$ be the open of regular pencils
whose lines avoid rank at most $n-3$ and meet rank at most $n-2$
in at most one reduced point. The locus $\Sigma_2$ of nonempty
intersection is smooth of codimension two. The normalized Hilbert
incidence modification satisfies
\[
       \widehat G|_{\mathcal U}
                 =\operatorname{Bl}_{\Sigma_2}\mathcal U
\]
(Theorem~\ref{thm:incidence-blowup-v160}). This is an identification
of the entire space over this open, not a chosen family in one fibre.
Over an incident line with intersection $p$, the fibre is the full
$\Pj^1$ of lines through the tangent-normal point in the exceptional
$\Pj(N_{Z_{n-2}/P,p})=\Pj^2$. Each parameter gives the strict
transform of the original line with one reduced nodal tail. Locally
the universal family is $xy=t$, with $t$ a boundary parameter, and
the exceptional divisor has discrepancy one. The identification
supplies a Cartier-ideal universal property.

All resolved multiplication systems on the exceptional plane are
complete Veronese systems of degrees
\[
             (j-h(n-2))_+-2(j-h(n-1))_+,
                    \qquad 1\leq j\leq hn
\]
(Proposition~\ref{prop:exceptional-powers-v160}). This includes every
exceptional direction, whether or not its residual determinant is
squarefree. The incidence centre is the image, as a closed subscheme,
of the zero scheme of $J_{h(n-2)+1}$ on the universal pencil line.
Consequently the centre's ordinary powers and their Rees algebra
come from the intrinsic multiplication diagram of the failure family
(Theorem~\ref{thm:failure-incidence-v160} in the companion).
A nonzero first normal deformation selects the actual embedded tail.
The geometry of the closed algebra and the additional direction in
its deformation family are thus distinguished precisely.

The proof uses classical smooth blow-ups and complete quadrics.
The assertion proved here is the identification with the specified
Hilbert graph, the classification of all its fibres over $\mathcal U$,
and its realization by failure multiplication. No novelty is assigned
to the general blow-up universal property or the ambient variety of
complete quadrics.

\subsection{Embedded boundary fibres}
Let a family of regular pencils over $\C[[\tau]]$ have generically
simple spectrum and simultaneous ordinary transverse clusters of
multiplicities $c_a$ in its special fibre. Put $r_a=n-c_a$.
At a cluster, the symmetric Schur complement has the form
\[
 S_a=z_a C_a+\tau D_a+O((z_a,\tau)^2),\qquad
 \det(\alpha C_a+\beta D_a)\text{ squarefree},
\]
with $C_a$ nonsingular. The schematic closure of the lifted generic
pencil in complete quadrics is the blow-up of the cluster points in
the total pencil surface. Its special fibre is a reduced nodal main
curve with one tail at each cluster. The tail is the lifted residual
pencil $[\alpha C_a+\beta D_a]$, marked by its attachment $[C_a]$
(Theorem~\ref{thm:ordinary-fibres-v159}). This statement includes
simultaneous clusters of every allowed multiplicity and is unchanged
by higher-order perturbations with fixed special pencil and first
normal jets.

All multiplication systems on those tails are determined. For
$1\leq p\leq hn$, set
\[
 b_a=(p-hr_a)_+,\qquad s_p=(p-h(n-1))_+,\qquad
 \delta_{a,p}=b_a-c_a s_p.
\]
The exact local power ideal and its resolved tail map are
\[
 J_p(A)=(z_a,\tau)^{\delta_{a,p}}(\det S_a)^{s_p},
 \qquad F_a\longrightarrow\Pj^{\delta_{a,p}}
               \text{ the complete Veronese map}.
\]
A degree-zero map is constant. In particular
$J_{h(n-c_a)+1}=(z_a,\tau)$: one multiplication ideal selects the
exceptional component, and the remaining products give all its linear
systems. This is not the formal pullback of an unrelated flat algebra.

For a fixed semisimple pencil $R_0$, the residual lines through each
$[C_a]$ give a family of distinct embedded limits. If
$\rho:\widehat G\to\Gr(2,\Sym^2V)$ is the normalized Hilbert incidence
modification, then
\[
 \dim\rho^{-1}(R_0)\geq
             \sum_a\left(\binom{c_a+1}{2}-2\right)
\]
(Theorem~\ref{thm:boundary-moduli-v159}). The parameter spaces, their
flat embedded families and their smoothing realizations are explicit.
This is a modular family inside the fibre, not a classification of
all its components. Already a semisimple triple collision gives a
four-dimensional family with the same limiting failure algebra.

\subsection{All ranks and singular pencils}
On the rank-$r$ locus retain the powers through $hr$, including the
last one. The last power is the degree-$2h$ Pluecker image of the
image $r$-plane. The simultaneous graph is precisely
$\mathrm{CQ}(\mathcal U)\to\Gr(r,V)$, without further normalization
(Theorem~\ref{thm:all-rank-graph-v158}). Complete quadrics and their
orbit strata are classical; the theorem realizes their full projection
and scheme structure by powers in one algebra.

For a symmetric pencil of normal rank $r$, the kernel splitting
$\bigoplus_a\OO(-\varepsilon_a)$ is read from finitely many coefficient
syzygy maps. If their nullities are $\kappa_j$, then
\[
 \#\{a:\varepsilon_a=j\}=\kappa_j-2\kappa_{j-1}+\kappa_{j-2}.
\]
The relative boundary contacts give the positive Smith exponents,
with their spectral positions, and satisfy
\[
 \sum_{i=1}^{r-1}(r-i)\deg D_i+
          2\deg(\Lambda\to\Gr(r,V))=r,
 \qquad\deg(\Lambda\to\Gr(r,V))=\sum_a\varepsilon_a.
\]
Together these recover all complex congruence data, including singular
pencils (Theorem~\ref{thm:singular-complete-data-v158}). The complex
Kronecker classification is a classical input to this interpretation.
The ordinary-collision Hilbert theorem is a separate assertion with
its stated transverse hypotheses, not a claim about arbitrary singular
special fibres.

\subsection{Intrinsic origin and logical dependencies}
The companion paper extracts an oriented source projective space from
the first relation of an unmarked sharp failure algebra, before its
last step recovering the pencil coefficient line. The first-jet bundle
and a Cartan quotient on that source construct a scalar-normalized
algebra bundle and its projective multiplication diagram
(Theorem~\ref{thm:intrinsic-envelope-v158} and
Proposition~\ref{prop:universal-envelope-v159}). It is not a preferred
subquotient of the original algebra, and it does not choose an
$n$-plane in an auxiliary coefficient complement of dimension $n^2-1$.
The case $h=1$ gives a distinguished diagram without an exponent choice.

The ideal, graph, all-pencil and boundary proofs here start with $V$
and a pencil and do not use the sharp inverse. Only their interpretation
as functors of the original unmarked failure algebra invokes that
companion. Conversely, the sharp inverse and source descent do not
use the present boundary theorems. The two core proofs are independent;
the intrinsic-envelope application links them with these inputs stated.

The later sections preserve binary normalization and conductor theory,
all-collision incidence charts and reciprocal-fibre homology. The
multiplier formula is explicitly a consequence of the classical
symmetric-determinantal calculation, not an independent novelty claim.
The theorem-level comparison with Ballico's failure-locus paper
\cite{Ballico93} remains limited by unavailable full text; no assertion
of historical nonanticipation follows from that absence.

\section{Power coefficients and their twists}
\label{sec:power-conventions-v158}
For a vector space $V$, put $M(V)=\Sym^2V$, with congruence action
$A\mapsto gAg^{\mathsf t}$, and define
\begin{equation}\label{eq:envelope-presentation-v158}
 B_h(V)=\Sym(M(V))/\langle (v^2)^{h+1}:v\in V\rangle .
\end{equation}
The degree of $M(V)$ is one. The generating space of the ideal is the
Cartan inclusion $\Sym^{2h+2}V\hookrightarrow\Sym^{h+1}(\Sym^2V)$,
normalized by $v^{2h+2}\mapsto(v^2)^{h+1}$. This definition requires
neither a volume form nor a complementary space. It agrees with
$B_{n,h}$ in Theorem~\ref{thm:quadratic-section-v155} when $W=V^*$.

\begin{proposition}[The coefficient map and its twists]
\label{prop:coefficient-conventions-v158}
Let $n=\dim V$ and $N=nh$. The quotient multiplication $\mu_p$ gives
an equivariant injection
\begin{equation}\label{eq:coefficient-duals-v158}
 \mu_p^*:B_h(V)_p^*\hookrightarrow\Sym^p(M(V)^*),
 \qquad (\mu_p^*\varphi)(A)=\varphi(A^p).
\end{equation}
Its image is the full coefficient space defining $J_p$, without a
choice of apolar coordinates. The socle and perfect pairings are
\begin{equation}\label{eq:perfect-twist-v158}
 B_h(V)_N=(\det V)^{2h},\qquad
 B_h(V)_p^*\simeq B_h(V)_{N-p}\otimes(\det V^*)^{2h}.
\end{equation}
For a line $L$, $B_h(V\otimes L)_p=B_h(V)_p\otimes L^{2p}$,
compatibly with all products. On $\Pj(M(V))$, the $p$th power is a
section of $B_h(V)_p\otimes\OO(p)$, and its base ideal is the image
of $B_h(V)_p^*\otimes\OO(-p)\to\OO$.
\end{proposition}
\begin{proof}
A homogeneous polynomial of degree $p$ is identified with its symmetric
$p$-linear polarization, normalized so that its value on $A$ is its
value on $(A,\ldots,A)$. Dualizing the surjection $\mu_p$ gives exactly
\eqref{eq:coefficient-duals-v158}. Multinomial coefficients appearing
in a monomial basis express this normalization and are nonzero; no
dual representation is silently replaced by the original one.
Theorem~\ref{thm:quadratic-section-v155} identifies the quotient with
determinant apolarity. The determinant of $A:V^*\to V$ takes values
in $(\det V)^2$. Choosing a trivialization gives a scalar polynomial
whose inverse-system line has character $(\det V^*)^{2h}$; its dual
socle therefore has character $(\det V)^{2h}$. The Gorenstein pairing
gives the second formula. Each generator in $M(V)$ has scalar weight
two, proving the line-twist assertion directly from
\eqref{eq:envelope-presentation-v158}. The tautological line
$\OO(-1)\to M(V)\otimes\OO$ has $p$th product
$\OO(-p)\to B_h(V)_p\otimes\OO$, which proves the last assertion.
\end{proof}


\subsection{Graph schemes and classical inputs}
The projective graph argument uses the following precise statement.
It applies to the ambient space and also to the two-dimensional total
base of the degeneration computed below.

\begin{lemma}[Simultaneous graph with invertible twists]
\label{lem:simultaneous-graph-v158}
Let $X$ be an integral noetherian complex scheme and let
$W_i\otimes L_i^{-1}\to\OO_X$ be nonzero coherent systems of
sections, with finite-dimensional $W_i$ and invertible $L_i$.
Write $K_i$ for their image ideals, and let
$U=X\setminus\bigcup_i V(K_i)$. The scheme-theoretic graph closure
of $U\to\prod_i\Pj(W_i^*)$ is
$\operatorname{Bl}_{\prod_iK_i}X$ with its specified graph embedding.
This statement identifies the scheme and the maps, not only their
normalizations. Replacing the product ideal by an invertible twist
or a positive power does not change its blow-up over $X$.
\end{lemma}
\begin{proof}
The open $U$ is dense because all $K_i$ are nonzero and $X$ is integral.
After the Segre embedding the simultaneous map is given by all products
of one section from each system. On a trivializing affine open, these
products generate $K=\prod_iK_i$. The homogeneous coordinate algebra
of the graph is the image of the polynomial algebra on these
products in $\OO_X[t]$, namely the Rees algebra
$\bigoplus_{m\geq0}K^m t^m$. It is a subalgebra of the function-field
polynomial ring. Thus its Proj is exactly the schematic closure of
the graph, with no extra component. The Segre equations hold on the
dense graph and hence on this integral closure scheme, so the graph
lands in the product, not merely in its ambient projective space.
Local trivializations glue. An invertible ideal changes the degree-$m$
piece by its $m$th tensor power and does not change relative Proj;
a positive power replaces the Rees algebra by a Veronese subalgebra.
\end{proof}

For the power system, $L_i=\OO(p)$ and $W_i=B_h(V)_p^*$;
for the exterior system of order $q$, $L_i=\OO(q)$ and its sections
are the $q$-minors. The determinant system is invertible on the
integral ambient space. It may be omitted from the blow-up product,
although its zero divisor must still be retained in contact formulas.
Likewise the first-coordinate identity map has unit base ideal.
These observations specify the projective twists in the proof of
Theorem~\ref{thm:power-graph-v157}.

The precise classical scheme-level complete-quadric construction is
\cite[Theorem~6.3]{Vainsencher1982}; the graph of exterior-power maps
and its successive symmetric-rank blow-ups are stated in
\cite[Remark~2.5 and Construction~2.6]{Massarenti2020}.
The relative rank version is recalled in
\cite[Definition~2.5, Remark~2.6 and Theorem~2.14]{CasarottiCornianiMassarenti2023}.
Consequently smoothness, the flag description, and the orbit
stratification in Corollary~\ref{cor:native-strata-v157} are classical
consequences of the graph identification, not new classifications.

There is also a precise invariant-ideal predecessor to the universal
identity. Henriques--Varbaro \cite[\S2.4, Theorems~2.6--2.8]{HenriquesVarbaro2016}
recall Abeasis's classification for the symmetric coordinate ring,
its products of determinantal ideals and their integral closures,
and the passage between invariant generators and products indexed
by diagrams. Our coefficient-space comparison uses this same
multiplicity-free setting. Theorem~\ref{thm:universal-power-v157}
identifies a particular multiplication coefficient map with the
balanced determinantal product as an ordinary ideal; it is not
an independent classification of invariant ideals. The exact
coefficient identification in Proposition~\ref{prop:coefficient-conventions-v158}
explains which map is being specialized on nonreduced bases.
We assert no exhaustive historical nonanticipation of the displayed
identity in all invariant-ideal terminology. Similarly the
multiplier calculation in Corollary~\ref{cor:power-multiplier-v157} is the
substitution of its exponents in the symmetric determinantal formula
\cite[Theorem~4.8]{HenriquesVarbaro2016}, not an independent
multiplier-ideal theorem.


\section{Quadratic sections of reciprocal normalization fibres}
\label{sec:quadratic-sections-v155}
The reciprocal presentation has a distinguished linear coefficient
section on which the entire algebra, rather than only its minimal
relations, can be determined.  Throughout this section a decomposition
$E=\C a\oplus W$ is fixed.  The section is equivariant for
$\operatorname{GL}(W)$; it is not asserted to be independent of an
arbitrary change of that decomposition.

Put $d=\dim W$, $g\geq k\geq2$, and $h=\lfloor g/2\rfloor$.
In the algebra $F_{g,k}(W)$ of \eqref{eq:reciprocal-algebra-v153},
let $\mathfrak l$ be the ideal generated by the coefficients of all
$Q_i$ with $i\ne2$.  Write $Q=Q_2$ and
\begin{equation}\label{eq:quadratic-section-v155}
 B_{d,h}=\Sym((\Sym^2W)^*)/
                    (\operatorname{coeff}Q^{h+1}).
\end{equation}
Here the ordinary degree of a coefficient of $Q$ is one; its
reciprocal weight is two.  These two gradings will always be
specified separately.  A partition $\lambda\subseteq(h^d)$ has
at most $d$ rows, each of length at most $h$.

\begin{theorem}[The entire quadratic coefficient section]
\label{thm:quadratic-section-v155}
There is an isomorphism of $\operatorname{GL}(W)$-equivariant graded
algebras
\begin{equation}\label{eq:exact-section-v155}
                 F_{g,k}(W)/\mathfrak l\simeq B_{d,h}.
\end{equation}
Under the coefficient/differential pairing, $B_{d,h}$ is the apolar
algebra of the $h$th power of the determinant of a generic symmetric
$d\times d$ matrix.  Its ordinary-degree decomposition is
\begin{equation}\label{eq:section-character-v155}
 (B_{d,h})_p\simeq
 \bigoplus_{\substack{\lambda\subseteq(h^d)\\|\lambda|=p}}
                         \mathbb S_{2\lambda}(W^*).
\end{equation}
In particular it is Gorenstein of socle degree $dh$, with
one-dimensional socle representation $(\det W^*)^{2h}$.  Its Hilbert
function and length are
\begin{align}
 H_{d,h}(p)&=\sum_{\substack{\lambda\subseteq(h^d)\\|\lambda|=p}}
       \prod_{1\leq i<j\leq d}
             \frac{2\lambda_i-2\lambda_j+j-i}{j-i},
                     \label{eq:section-hilbert-v155}\\
 L_{d,h}&=\prod_{i=1}^d\frac{h+d-i+1}{d-i+1}
       \prod_{1\leq i<j\leq d}
             \frac{2h+2d-i-j+2}{2d-i-j+2}.
                     \label{eq:section-length-v155}
\end{align}
The Hilbert function is symmetric.  A degree-one element is strong
Lefschetz precisely when its associated symmetric matrix has full
rank.  More generally, an element associated with a rank-$r$ matrix
has nilpotence index $hr+1$ (index one for the zero element).
Consequently
\begin{equation}\label{eq:full-fibre-bounds-v155}
 \length F_{g,k}(W)\geq L_{d,h},\qquad
 \mathfrak m^{dh}\ne0\text{ in }F_{g,k}(W).
\end{equation}
The Gorenstein and Lefschetz assertions concern the section
$B_{d,h}$, not the full fibre $F_{g,k}(W)$.
\end{theorem}
\begin{proof}
After setting the indicated coefficients to zero, the reciprocal
series is $(1+Qz^2)^{-1}$.  The first even integer in
$\{g+1,\ldots,g+k\}$ is $2(h+1)$; it belongs to this interval
because $k\geq2$.  Every odd coefficient vanishes, and the even
ones are signed powers of $Q$.  Thus the specialized ideal is
exactly $(\operatorname{coeff}Q^{h+1})$: all its higher powers have
coefficients in this ideal.  This proves \eqref{eq:exact-section-v155}
as an equality of ideals, not only an equality of radicals.

For clarity about the pairing, write
$Q(x)=\sum_i y_{ii}x_i^2+\sum_{i<j}y_{ij}x_ix_j$ and let $X$ have
independent symmetric entries $X_{ij}$, $i\leq j$.  Send $y_{ij}$
to $\partial/\partial X_{ij}$.  Then $Q(x)$ acts by differentiation
in the rank-one direction $xx^{\mathsf t}$.  Since
$\det(X+zxx^{\mathsf t})^h$ has degree at most $h$ in $z$, the
coefficients of $Q(x)^{h+1}$ annihilate $(\det X)^h$.
The reverse containment is the symmetric-matrix case of
\cite[Remark~11, pp.~8--9]{NagaokaWachi2024}: that apolar ideal is
generated by the full general-linear orbit of the $(h+1)$st power
of a rank-one differentiation direction.  Powers of linear forms
span their symmetric power, so this orbit span is precisely the
coefficient space just displayed.  No assertion that an arbitrary
quotient of an apolar algebra is Gorenstein is used.

We record explicitly how the classical representation formulas
translate here.  In \cite[Example~9 and Proposition~10]{NagaokaWachi2024},
the surviving highest weights are indexed by
$k_1,\ldots,k_d\geq0$ with $\sum k_i\leq h$, and their degrees are
$\sum i k_i$.  Set $\lambda_i=\sum_{j\geq i}k_j$.
These are exactly the partitions in $(h^d)$, of size $\sum i k_i$,
and their polynomial representations are $\mathbb S_{2\lambda}(W^*)$.
This proves \eqref{eq:section-character-v155}; the ordinary Weyl
dimension formula gives \eqref{eq:section-hilbert-v155}.
The same apolar module is the finite-dimensional highest-weight
$\mathfrak{sp}_{2d}$-module of weight $h\omega_d$, with its Levi
character shifted by $\det^h$; see
\cite[Proposition~16 and Lemma~13]{NagaokaWachi2024}.
For the positive roots $e_i-e_j$, $e_i+e_j$, and $2e_i$, its Weyl
dimension factors are respectively $1$,
$(2h+2d-i-j+2)/(2d-i-j+2)$, and
$(h+d-i+1)/(d-i+1)$.  Their product is
\eqref{eq:section-length-v155}.
Apolarity to one homogeneous polynomial of degree $dh$ gives the
perfect pairings and the socle.  Equivalently, the complementary
partition $(h-\lambda_d,\ldots,h-\lambda_1)$ gives the dual summand
with the indicated determinant twist.

The strong Lefschetz assertion is the symmetric-matrix case of
\cite[Theorem~4 and Example~5]{NagaokaWachi2024}.
There is a direct argument for the asserted nilpotence index.
If $\ell_A=\sum_{i\leq j}A_{ij}y_{ij}$, then
$\ell_A^p=0$ in the apolar algebra exactly when
$\partial_z^p\det(X+zA)^h=0$ as a polynomial in $X$.
Congruence reduces a rank-$r$ symmetric matrix to
$\diag(1,\ldots,1,0,\ldots,0)$.
The last determinant power has degree exactly $hr$ in $z$, as
is seen by taking $X$ diagonal with independent entries.
This proves the index.  Finally a surjective local algebra map
surjects on every power of its maximal ideal.  The section has
nonzero degree $dh$, proving both assertions in
\eqref{eq:full-fibre-bounds-v155}.
\end{proof}

\begin{corollary}[Size and Loewy depth of the actual common boundary fibre]
\label{cor:boundary-size-v155}
In the notation of Theorem~\ref{thm:universal-boundary-v152}, put
$d=n^2-1$ and $h=n(n-1)/2$.  The entire normalization fibre at
$K_\infty$ has length at least $L_{d,h}$ and its maximal ideal has
nonzero $dh$th power.  For the ternary pencil boundary,
\begin{equation}\label{eq:ternary-large-section-v155}
 \length F_{6,2}(\C^8)\geq922268360,
       \qquad\mathfrak m^{24}\ne0.
\end{equation}
Its quadratic coefficient section has exactly that length, ordinary
socle degree $24$, and reciprocal-weight socle $48$.
\end{corollary}
\begin{proof}
The common boundary fibre is $F_{g,k_n}(\C^{n^2-1})$, with
$g=n(n-1)\geq k_n\geq2$.  Apply the theorem.
For $n=3$, substitute $(d,h)=(8,3)$ into
\eqref{eq:section-length-v155}.
This computes a whole closed section, not the length of the original
fibre.  The previously obtained $44$ coefficient variables and
$9867$ minimal equations remain counts of different invariants.
\end{proof}


\section{Universal power ideals}
\label{sec:universal-power-v157}
The power-zero construction admits an equality of ideals before any
restriction to a pencil or to a valuation ring. This is the point at
which the reciprocal section acquires a global boundary geometry.
Let $V$ have dimension $n\geq2$, put $M=\Sym^2V$, and use the algebra
$B_{n,h}$ of Theorem~\ref{thm:quadratic-section-v155} with $W=V^*$.
Thus $(B_{n,h})_1=M$. If $A$ is a symmetric matrix with entries in a
commutative $\C$-algebra $T$, denote its element of $M\otimes T$ by
$\ell_A$, and set
\begin{equation}\label{eq:universal-power-definition-v157}
 J_p(A)=\im\bigl((B_{n,h})_p^*\otimes T\longrightarrow T,
                  \quad\varphi\longmapsto\varphi(\ell_A^p)\bigr).
\end{equation}
Write $I_q(A)$ for the ideal of all $q$-minors, with $I_0(A)=T$.
In this section powers of ideals are ordinary powers unless parentheses
explicitly indicate symbolic powers.

\begin{theorem}[Universal determinantal formula]
\label{thm:universal-power-v157}
For every commutative $\C$-algebra $T$, every symmetric $n\times n$
matrix $A$ over $T$, and $0\leq p\leq nh$, write $p=hq+s$ with
$0\leq s<h$. Then
\begin{equation}\label{eq:universal-power-v157}
             J_p(A)=I_q(A)^{h-s}I_{q+1}(A)^s.
\end{equation}
When $s=0$ the second factor is omitted. For $p>nh$, $J_p(A)=0$.
These are equalities of ideals, not of radicals or integral closures,
and commute with every homomorphism of $\C$-algebras. In particular,
for every symmetric map $A:T^n\to T^n$,
\begin{equation}\label{eq:universal-Fitting-v157}
       J_{h(n-j)}(A)=\Fitt_j(\coker A)^h,\qquad 0\leq j\leq n.
\end{equation}
The formulas also hold locally, and hence as ideal sheaves, for
symmetric maps of vector bundles with a line-bundle twist on an
arbitrary $\C$-scheme.
\end{theorem}

We separate the classical determinantal input. For an arbitrary matrix
$X$ over a characteristic-zero polynomial ring, the balancing inclusion
\begin{equation}\label{eq:minor-balancing-v157}
              I_u(X)I_v(X)\subseteq I_{u+1}(X)I_{v-1}(X),
              \qquad 0\leq u\leq v-2,
\end{equation}
is the determinantal symmetrization lemma. See
\cite[Lemma~2.5 and Theorem~2.4]{BrunsConca2003}, which attributes the
characteristic-zero product theorem to De Concini--Eisenbud--Procesi.
Because this is a polynomial ideal inclusion, it remains true after
specialization to a symmetric matrix; no flatness of that specialization
is needed. This classical lemma is not a claim of the present paper.

\begin{lemma}[The bounded-row generating space]
\label{lem:bounded-row-v157}
Let $S=\Sym(\Sym^2V^*)$ and let $A$ be its universal symmetric matrix.
For $p=hq+s$, the degree-$p$ part generating
$I_q(A)^{h-s}I_{q+1}(A)^s$ is
\begin{equation}\label{eq:bounded-row-v157}
 \bigoplus_{\substack{|\lambda|=p,\ \ell(\lambda)\leq n\\
                         \lambda_1\leq h}}\mathbb S_{2\lambda}(V^*)
                         \ \subseteq S_p.
\end{equation}
\end{lemma}
\begin{proof}
The symmetric Cauchy decomposition of $S_p$ is multiplicity free,
with summands $\mathbb S_{2\lambda}(V^*)$, $|\lambda|=p$; see
\cite[\S2.4]{HenriquesVarbaro2016}. We may use the dual general
linear group to regard these as polynomial representations. Every
minor has torus weight at most two in any one coordinate: its row
set and column set each use that coordinate at most once. A product
of $h$ minors therefore has no weight with a coordinate exceeding
$2h$. Its invariant span cannot contain a highest weight
$2\lambda$ with $\lambda_1>h$. This proves one inclusion.

For the converse, let $c_1\geq\cdots\geq c_h\geq0$ be the column
lengths of a partition $\lambda$ occurring on the right, padded with
zeros. The product of the leading principal minors of orders $c_i$
is a nonzero highest-weight vector of weight $2\lambda$ in $S_p$.
It belongs initially to $\prod_{i=1}^h I_{c_i}(A)$.
Whenever two column lengths differ by at least two, apply
\eqref{eq:minor-balancing-v157} to increase the smaller length by one
and decrease the larger by one. The sum of the squares of the lengths
strictly decreases, so the procedure terminates. At termination the
lengths are $q+1$ repeated $s$ times and $q$ repeated $h-s$ times,
since their sum is $p$. Thus that highest-weight vector belongs to
$I_q^{h-s}I_{q+1}^s$. The latter ideal is invariant; its degree-$p$
part contains the full irreducible summand. Multiplicity freeness
proves the equality. The case $q=0$ includes zero-length minors,
which are the constant $1$.
\end{proof}

\begin{proof}[Proof of Theorem~\ref{thm:universal-power-v157}]
First work over $S$ with its universal matrix. The coefficient map of
$\ell_A^p$ is the dual of the multiplication quotient
$\Sym^p M\twoheadrightarrow (B_{n,h})_p$, up to nonzero multinomial
constants. Consequently its degree-$p$ coefficient span is exactly
$(B_{n,h})_p^*\subseteq S_p$. The character formula
\eqref{eq:section-character-v155} identifies that span with
\eqref{eq:bounded-row-v157}. Both ideals in
\eqref{eq:universal-power-v157} are generated in degree $p$, so the
lemma proves their equality in $S$. This step uses the full coefficient
space, not just evaluation on diagonal matrices or closed points.

For a matrix over $T$, substitute its entries in this polynomial
identity. Taking the image of the fixed finite-dimensional coefficient
map commutes with scalar extension, as does taking minors and products
of ideals. This proves the assertion even when $T$ is nonreduced,
non-Noetherian, or not a domain. The Fitting formula follows from its
minor presentation. Changing a vector-bundle frame applies an invertible
congruence to $A$ and an invertible linear transformation to each
coefficient space; changing a local frame of the twisting line multiplies
generators by a unit. Therefore the local ideal identities glue. Finally,
$(B_{n,h})_p=0$ for $p>nh$.
\end{proof}

\begin{example}[A nonreduced test base]
\label{ex:nilpotent-base-v157}
Take $T=\C[\epsilon]/(\epsilon^3)$, $n=h=2$, and
$A=\diag(\epsilon,0)$. Then $J_1(A)=(\epsilon)$ and
$J_2(A)=(\epsilon^2)\ne0$, whereas $J_3(A)=J_4(A)=0$.
Every closed-point matrix is zero. Thus equality on geometric fibres
would not establish the theorem. Nor does
\eqref{eq:universal-Fitting-v157} license taking unique ideal roots:
for example, in $\C[\epsilon]/(\epsilon^2)$ the distinct ideals
$(\epsilon)$ and $0$ have the same square. At $h=1$ the power ideals
are the Fitting ideals themselves over every base.
\end{example}


\section{The reciprocal power graph and complete quadrics}
\label{sec:power-graph-v157}
Let $P=\Pj(M)$, let $P^{\mathrm{nd}}$ be the nonsingular-quadric open,
and define $\mathcal G_h$ as the scheme-theoretic closure of the graph
\begin{equation}\label{eq:power-graph-v157}
 P^{\mathrm{nd}}\longrightarrow
 P\times\prod_{p=1}^{nh-1}\Pj((B_{n,h})_p),
 \qquad [A]\longmapsto([A],[\ell_A],\ldots,[\ell_A^{nh-1}]).
\end{equation}
The first factor records the base point. We use the reduced, integral
graph closure, equivalently its scheme-theoretic image; no additional
normalization is inserted in this definition. Let $\mathrm{CQ}(V)$
be the classical space of complete quadrics: the graph closure of
$[A]\mapsto([\bigwedge^2A],\ldots,[\bigwedge^{n-1}A])$ over $P$.
For $n=2$ the empty product gives $\mathrm{CQ}(V)=P$.

\begin{theorem}[A compactification determined by multiplication]
\label{thm:power-graph-v157}
For every $h\geq1$ there is a canonical $\operatorname{GL}(V)$-equivariant
isomorphism over $P$
\begin{equation}\label{eq:power-graph-CQ-v157}
                      \mathcal G_h\simeq\mathrm{CQ}(V).
\end{equation}
In particular the graph is smooth and projective, is independent of
$h$ as a scheme over $P$, and has a simple-normal-crossing boundary.
Write $\pi:\mathrm{CQ}(V)\to P$, and let $E_r$ be its boundary divisor
of rank break $r$, $1\leq r<n$; for $r<n-1$ it is exceptional over
rank at most $r$, and $E_{n-1}$ is the strict transform of the determinant
divisor. Then every power ideal has the exact total transform
\begin{equation}\label{eq:power-boundary-v157}
 J_p\OO_{\mathrm{CQ}(V)}=
 \OO_{\mathrm{CQ}(V)}\left(-\sum_{r=1}^{n-1}(p-hr)_+E_r\right),
                    \qquad 0\leq p\leq nh,
\end{equation}
where $(a)_+=\max(a,0)$. Removing this common divisor from the universal
$p$th power gives a nowhere-zero projective section on the whole
compactification.
\end{theorem}

The smoothness, boundary and flag interpretation of complete quadrics
are classical, due to the complete-form construction; see
\cite{Vainsencher1982,Massarenti2020}. The assertion here identifies
that particular compactification, with all of its scheme structure,
from the multiplication in one reciprocal normalization-fibre section.
It does not assert that complete quadrics is a newly constructed variety.

\begin{proof}
We first recall the graph calculation with its scheme structure.
If rational maps from an integral scheme are given by generators of
nonzero ideals $K_1,\ldots,K_a$, their simultaneous graph closure is
$\operatorname{Bl}_{K_1\cdots K_a}$ of the base. Indeed, after the Segre
embedding the coordinates are all products of one generator from each
ideal. The closure of this graph is the Proj of the Rees algebra of
that product: on a generator chart its coordinate ring is the subring
of the function field obtained by adjoining the generator ratios.
These are exactly the affine charts of both constructions. Tensoring
an ideal by an invertible ideal, or replacing it by a positive power,
does not change the blow-up, the latter by the Veronese description
of Proj.

Apply this calculation on $P$. The sheaf $I_1\OO_P$ is the unit ideal,
and $I_n\OO_P$ is invertible. For each $2\leq q\leq n-1$, the exponent
of $I_q$ in $\prod_{p=1}^{nh-1}J_p$ is
\[
 \sum_{s=0}^{h-1}(h-s)+\sum_{s=1}^{h-1}s=h^2.
\]
Theorem~\ref{thm:universal-power-v157} therefore gives, up to an
invertible ideal,
\[
       \prod_{p=1}^{nh-1}J_p
                =\left(\prod_{q=2}^{n-1}I_q\right)^{h^2}.
\]
The same graph calculation for the exterior-power maps identifies
$\mathrm{CQ}(V)$ with the blow-up of $\prod_{q=2}^{n-1}I_q$.
This proves \eqref{eq:power-graph-CQ-v157} as an isomorphism of schemes,
not only after normalization. Both isomorphisms are the identity on
the common dense open, hence are canonical and equivariant.

For completeness, the local calculation behind the boundary formula
is as follows. The complete-quadric construction blows up the strict
rank loci successively. On its Gaussian charts, after a congruence
and projective normalization, its matrix has the form
\begin{equation}\label{eq:CQ-chart-v157}
 U\diag(1,\tau_1,\tau_1\tau_2,\ldots,
                         \tau_1\cdots\tau_{n-1})U^{\mathsf t},
\end{equation}
with $U$ lower unitriangular and $E_r=(\tau_r=0)$.
One obtains these charts inductively by making the first pivot one,
taking a Schur complement, blowing up the zero residual matrix, and
repeating. At each step the next nonzero residual quadratic form has
a nonisotropic vector, so congruence translates of these charts cover
the construction. This also explains the smooth charts and the normal
crossings stated in the classical construction.

Cauchy--Binet, applied to \eqref{eq:CQ-chart-v157}, shows that every
$q$-minor is divisible by
$\tau_1^{q-1}\tau_2^{q-2}\cdots\tau_{q-1}$.
The leading principal $q$-minor is exactly that product. Thus, on
these charts and therefore globally,
\begin{equation}\label{eq:minors-on-CQ-v157}
 I_q\OO_{\mathrm{CQ}(V)}=
              \OO\left(-\sum_{r<q}(q-r)E_r\right).
\end{equation}
Using \eqref{eq:universal-power-v157}, its exponent along $E_r$ is
$(h-s)(q-r)_++s(q+1-r)_+=(p-hr)_+$.
This proves \eqref{eq:power-boundary-v157}. The coefficient sections
generate the transformed ideal; division by its local generator makes
them generate the unit ideal, proving the last assertion.
\end{proof}

\begin{corollary}[Native boundary strata]
\label{cor:native-strata-v157}
The power graph has one congruence orbit for every subset
$I\subseteq\{1,\ldots,n-1\}$. It is the stratum
$E_I^\circ=(\bigcap_{r\in I}E_r)\setminus\bigcup_{r\notin I}E_r$,
of codimension $|I|$, and its closure is $\bigcap_{r\in I}E_r$.
Its points are flags with nondegenerate projective quadratic forms on
successive quotients; the quotient dimensions are the successive gaps
in $\{0\}\cup I\cup\{n\}$. The simultaneous normalized power values
recover this complete-quadric datum. They are not the one common limit
for the full cotangent-coordinate group.
\end{corollary}
\begin{proof}
The flag description and boundary stratification are the classical
complete-quadric description. The general linear group acts transitively
on flags of fixed dimensions, and over $\C$ every nondegenerate
quadratic form on a fixed-dimensional quotient is congruent to the
identity. Independent diagonal scalings remove the projective scalars.
The chart \eqref{eq:CQ-chart-v157} realizes each subset of vanishing
$\tau_r$ and permits further coordinates to vanish, giving the stated
codimensions and closures. The graph is a closed subscheme of the
product in \eqref{eq:power-graph-v157}; its projections therefore
retain the point, including its flag and quotient forms. The group
here is the native congruence group of $V$.
\end{proof}

\begin{remark}[Relative construction and base change]
\label{rem:relative-CQ-v157}
For a vector bundle $\mathcal V$ on any $\C$-scheme $S$, the associated
bundle with fibre $\mathrm{CQ}(\C^n)$ is a relative smooth projective
scheme over $S$. Equivariance glues the power maps and the boundary
formula. Formation of this associated bundle and the polynomial ideal
identity commute with arbitrary base change. This does not say that
taking the graph closure of an arbitrary restricted family commutes
with nonflat base change. The relative ambient resolution and the
closure of an individual family's graph are different constructions.
\end{remark}


\section{Power graphs of every rank}
\label{sec:all-rank-graph-v158}
For $1\leq r\leq n$, let $P^{(r)}\subset\Pj(\Sym^2V)$ be the
rank-exactly-$r$ locus. Define $\mathcal G_{r,h}$ as the
scheme-theoretic closure of
\begin{equation}\label{eq:rank-graph-definition-v158}
 P^{(r)}\longrightarrow
 \Pj(\Sym^2V)\times\prod_{p=1}^{hr}\Pj(B_h(V)_p),
 \qquad[A]\longmapsto([A],[A],\ldots,[A^{hr}]).
\end{equation}
Keeping the final power is essential when $r<n$. Write $\mathcal U$
for the tautological rank-$r$ bundle on $\Gr(r,V)$ and
$\mathrm{CQ}(\mathcal U)$ for its relative complete-quadric space.
For rank one this means $\Pj(\Sym^2\mathcal U)=\Gr(1,V)$.

\begin{theorem}[The rank-stratified power graph]
\label{thm:all-rank-graph-v158}
There is a canonical isomorphism
\begin{equation}\label{eq:rank-graph-isom-v158}
 \mathcal G_{r,h}\simeq\mathrm{CQ}(\mathcal U)
                  \longrightarrow\Gr(r,V).
\end{equation}
The final power in \eqref{eq:rank-graph-definition-v158} recovers the
image $r$-plane by the degree-$2h$ Pluecker embedding. In particular
the graph is smooth and projective, independently of $h$. For its
relative boundary divisors $E_i$, $1\leq i<r$, its power base ideals
have total transforms
\begin{equation}\label{eq:rank-boundary-v158}
 J_p\OO_{\mathcal G_{r,h}}=
 \OO_{\mathcal G_{r,h}}\left(-\sum_{i=1}^{r-1}(p-hi)_+E_i\right),
 \qquad 0\leq p\leq hr.
\end{equation}
The projective power maps in this formula have their common divisors
removed. All powers $p>hr$ vanish on the rank-$r$ locus.
\end{theorem}
\begin{proof}
For $U\subset V$, the map $B_h(U)\to B_h(V)$ from the presentation
\eqref{eq:envelope-presentation-v158} is injective: any linear projection
$V\to U$ splitting the inclusion induces a retraction on the quotient
algebras. The injection itself is independent of that projection.
Locally this gives subbundles $B_h(\mathcal U)_p\subset B_h(V)_p$.
If $A$ has image $U$ and rank $r$, it is a nonsingular element of
$\Sym^2U$, and $A^{hr}$ generates the socle
$B_h(U)_{hr}=(\det U)^{2h}$. As $U$ varies these lines give the
homogeneous embedding
\[
 \Gr(r,V)\hookrightarrow\Pj(\mathbb S_{(2h)^r}V)
                         \hookrightarrow\Pj(B_h(V)_{hr}).
\]
Indeed the line at the coordinate $r$-plane is a highest-weight line
of weight $(2h)^r$; its orbit map is the Pluecker embedding followed
by its degree-$2h$ linear system. In particular this is a closed
immersion, not just a recovery of the image plane on closed points.

Over the recovered Grassmannian, the remaining graph is precisely the
relative simultaneous power graph for $B_h(\mathcal U)$. By
Theorem~\ref{thm:power-graph-v157} on each frame chart, it is
$\mathrm{CQ}(\mathcal U)$. The relative graph embedding, the subbundle
injections above, and the final-power Grassmannian embedding together
give a closed immersion into the product in
\eqref{eq:rank-graph-definition-v158}. Its dense open is the graph
of $P^{(r)}$. Both its source and that graph closure are integral,
so equality of their scheme-theoretic images proves
\eqref{eq:rank-graph-isom-v158} without a normalization step.
Applying \eqref{eq:power-boundary-v157} in rank $r$ proves the
transform formula. Local split inclusion of the coefficient bundles
ensures that ambient and relative power ideals agree. The vanishing
assertion follows from the socle degree of $B_h(U)$.
\end{proof}

The varieties of complete singular quadrics are classical: see
\cite[Definition~2.5, Remark~2.6 and Theorem~2.14]{CasarottiCornianiMassarenti2023}.
The point of the theorem is the exact recovery of their image-plane
projection and entire scheme structure from the power maps in one
fixed algebra $B_h(V)$, simultaneously for every $r$. The classical
orbit stratification is background, not an additional novelty claim.


\section{Contacts and minimal indices of arbitrary pencils}
\label{sec:singular-pencil-v158}
Let $R=\langle A_0,A_1\rangle\subset\Sym^2V$ be a genuine pencil,
$\Lambda=\Pj(R)$, and $r$ its normal rank. Its tautological symmetric
map is $A:V^*\otimes\OO_\Lambda\to V\otimes\OO_\Lambda(1)$.
Write
\[
 K=\ker A,\qquad U=K^\perp\subset V\otimes\OO_\Lambda.
\]
The kernel is saturated and hence a subbundle on the smooth curve;
$U$ is the saturation of the generic image. Thus $A$ is a generically
nonsingular section of $\Sym^2U\otimes\OO_\Lambda(1)$.
The map on the rank-$r$ open lifts uniquely to
$\widetilde\Lambda\to\mathcal G_{r,h}$ by properness. Its
Grassmannian projection is the map defined by $U$, and its source is
still $\Lambda$. Let $D_i$ be its contact divisor with $E_i$.

\begin{lemma}[Symmetric diagonalization with a kernel]
\label{lem:symmetric-dvr-v158}
Over $\C[[t]]$, a symmetric matrix of rank $r$ over the fraction field
is congruent to $\diag(t^{e_1},\ldots,t^{e_r},0,\ldots,0)$, with
$e_1\leq\cdots\leq e_r$. For the local matrix of a genuine pencil,
$e_1=0$. The exponents are its nonzero Smith exponents.
\end{lemma}
\begin{proof}
Choose a matrix entry of smallest valuation. If a diagonal entry
has that valuation, it is a pivot. Otherwise an off-diagonal entry
does, and the value of the quadratic form on $e_i+c e_j$ has the
same valuation for some $c\in\C$: its leading coefficient is a
nonzero polynomial in $c$, since $2$ is invertible. Extend this
primitive vector to a basis. The pivot divides every entry in its
row, so symmetric row-and-column elimination clears the row and
column over the ring. All remaining entries have valuation at least
that of the pivot. Repeat until the remaining block is zero.
Units have square roots in $\C[[t]]$ and can be absorbed into the
basis. Invertible congruence preserves determinantal ideals, which
identify the resulting exponents with Smith exponents. A genuine
pencil has a nonzero matrix at every projective parameter, giving a
unit entry and $e_1=0$. This proof supplies the local congruence step
rather than relying on left-right Smith form alone; compare the
classical pencil congruence treatment in \cite{Thompson1991}.
\end{proof}

For $j\geq0$ define the finite multiplication-syzygy map
\begin{equation}\label{eq:toeplitz-map-v158}
 C_j(R):V^*\otimes\Sym^jR^*\longrightarrow
                         V\otimes\Sym^{j+1}R^*,\qquad q\longmapsto A q,
 \qquad\kappa_j=\dim\ker C_j(R).
\end{equation}
Set $\kappa_{-1}=\kappa_{-2}=0$. Under a basis of $R$ this is the block
Toeplitz map with rows $A_0q_0$, $A_0q_i+A_1q_{i-1}$, and $A_1q_j$;
its definition is independent of that basis.

\begin{theorem}[Complete power-and-syzygy data]
\label{thm:singular-complete-data-v158}
For every complex symmetric pencil, including every singular pencil,
the following formulas recover its complete congruence data.
If $\varepsilon_1,\ldots,\varepsilon_c$ are its minimal indices,
where $c=n-r$, then
\begin{align}
 K&\simeq\bigoplus_{a=1}^c\OO_\Lambda(-\varepsilon_a),
 &\kappa_j&=\sum_{a=1}^c(j-\varepsilon_a+1)_+,
                       \label{eq:kernel-splitting-v158}\\
 \#\{a:\varepsilon_a=j\}
   &=\kappa_j-2\kappa_{j-1}+\kappa_{j-2}.
                       \label{eq:minimal-differences-v158}
\end{align}
It is enough to use $0\leq j\leq\lfloor(n-1)/2\rfloor$.
At $x\in\Lambda$, let $0=e_1(x)\leq\cdots\leq e_r(x)$ be the
nonzero local Smith exponents. With $v_0=0$ and
$v_p(x)=\length(\OO_{\Lambda,x}/J_p)$ for $p\leq hr$, one has
\begin{align}
 v_{hq+s}(x)&=h\sum_{i=1}^{q}e_i(x)+s e_{q+1}(x),
                         \label{eq:singular-valuations-v158}\\
 e_i(x)&=\frac{v_{hi}(x)-v_{h(i-1)}(x)}h,
 &\operatorname{mult}_xD_i&=e_{i+1}(x)-e_i(x).
                         \label{eq:singular-contacts-v158}
\end{align}
In the first formula $s=0$ at $q=r$ and the last term is omitted.
The positive exponents at each supported point, together with the
minimal indices, determine the pencil up to congruence and projective
reparametrization of $\Lambda$. The spectral datum includes the
positions of the points, not only their local partitions.

There is also a global conservation law:
\begin{equation}\label{eq:conservation-v158}
 \sum_{i=1}^{r-1}(r-i)\deg D_i
       +2\deg\bigl(\Lambda\longrightarrow\Gr(r,V)\bigr)
       =r,\qquad
 \deg\bigl(\Lambda\longrightarrow\Gr(r,V)\bigr)
       =\sum_a\varepsilon_a.
\end{equation}
The Grassmannian degree uses its Pluecker polarization. The first
summand is the dimension of the regular summand in the Kronecker
congruence form. Every object and number in the theorem is invariant
under the ungraded failure-algebra isomorphisms of
Theorem~\ref{thm:intrinsic-envelope-v158}.
\end{theorem}
\begin{proof}
We use the classical complex symmetric Kronecker congruence theorem
\cite{Thompson1991}. A singular block of minimal index $\varepsilon$
has size $2\varepsilon+1$ and form
\[
 S_\varepsilon(s,t)=
 \begin{pmatrix}0&L_\varepsilon(s,t)^{\mathsf t}\\
                 L_\varepsilon(s,t)&0\end{pmatrix},\qquad
 L_\varepsilon=
 \begin{pmatrix}s&t& &0\\ &s&t& \\ &&\ddots&t\end{pmatrix},
\]
where $L_\varepsilon$ has $\varepsilon$ rows and $\varepsilon+1$
columns; $S_0$ is a zero $1\times1$ block. Its kernel is generated
by $((-t)^\varepsilon,s(-t)^{\varepsilon-1},\ldots,s^\varepsilon)$
in the first block. These components have no common projective zero,
so it gives $\OO(-\varepsilon)$ as a kernel subbundle. A regular
block has zero kernel as a sheaf. Direct sums prove the first formula
in \eqref{eq:kernel-splitting-v158}. Taking sections after twist by
$\OO(j)$ identifies them with \eqref{eq:toeplitz-map-v158}, proving
the second formula. The second finite difference of
$(j-\varepsilon+1)_+$ is one exactly at $j=\varepsilon$ and zero
otherwise. Each singular block uses $2\varepsilon+1\leq n$
dimensions, proving the stated finite bound. The canonical-form
classification is used here as a classical input, not reproved or
claimed as new.

Locally, symmetric diagonalization factors the matrix through the
saturated $r$-plane $U$ with an $r\times r$ symmetric block. Theorem
\ref{thm:universal-power-v157}, or its rank-$r$ version via the split
inclusion $B_h(U)\hookrightarrow B_h(V)$, yields
\eqref{eq:singular-valuations-v158}. Finite differences recover the
$e_i$. In the relative complete-quadric chart the successive ratios
of diagonal entries have valuations $e_{i+1}-e_i$, proving the contact
formula. The torsion of the pencil cokernel has precisely the positive
Smith exponents. The regular blocks of the classical congruence form
are specified by these elementary divisors, and its singular blocks
by the minimal indices; over $\C$ there are no additional signature
parameters. Retaining the supported points on $\Lambda$ makes the
statement compatible with arbitrary projective changes of the pencil
basis.

The exact sequence $0\to U\to V\otimes\OO\to K^*\to0$ gives
$\deg U=-\sum\varepsilon_a$, so the Grassmannian degree is
$\sum\varepsilon_a$. The determinant of the induced symmetric map
$U^*\to U(1)$ is a nonzero section of
$(\det U)^2\otimes\OO(r)$; its zero divisor has degree
$r-2\sum\varepsilon_a$. At each point its order is
$\sum e_i=\sum_{i<r}(r-i)\operatorname{mult}_xD_i$.
This proves \eqref{eq:conservation-v158}. The block-size identity
$n=n_{\mathrm{reg}}+\sum(2\varepsilon_a+1)$ identifies the first
summand with $n_{\mathrm{reg}}$. Finally the intrinsic first relation
extracts $\Lambda_A$ on its source and preserves all the maps involved;
changes of source lift or pencil basis give isomorphic kernel maps
and equal ideals, hence the asserted invariance.
\end{proof}

\begin{example}[Power contacts alone miss a singular invariant]
\label{ex:singular-separation-v158}
The size-six pencils $S_0\oplus S_2$ and $S_1\oplus S_1$ both have
constant rank four and no spectral torsion. All their power ideals
are the unit ideal for $p\leq4h$ and zero for $p>4h$, and every
contact divisor is empty. Their kernel splittings are respectively
$\OO\oplus\OO(-2)$ and $\OO(-1)^2$. Thus their first three
syzygy nullities are $(1,2,4)$ and $(0,2,4)$, distinguishing the two
congruence types. Both have Grassmannian degree two, as required by
\eqref{eq:conservation-v158}; the full splitting, rather than just
its degree, is necessary.
\end{example}

\begin{proposition}[Finite closed tests in families]
\label{prop:syzygy-family-v158}
For $j\geq0$ and $0\leq a\leq n(j+1)$, on the pencil
Grassmannian each condition $\kappa_j\geq a$ is the
closed determinantal condition on the bundle map $C_j$ given by minors
of size $n(j+1)-a+1$. These conditions commute with arbitrary base
change as schemes. Their fibrewise ranks need not be constant, so no
unqualified assertion that the kernel bundles commute with every
base change is made. Along a specialization over an integral curve,
$\kappa_j$ can only increase.
\end{proposition}
\begin{proof}
Use the tautological rank-two subbundle in
\eqref{eq:toeplitz-map-v158}. Its source has rank $n(j+1)$ and its
target rank $n(j+2)$. Nullity at least $a$ is equivalent to rank at
most $n(j+1)-a$, giving the stated minors. The ideals are polynomial
in the bundle map and therefore commute with arbitrary substitution.
Upper semicontinuity of nullity, or specialization of these minors,
gives the last assertion.
\end{proof}


\section{Boundary contacts and spectral schemes}
\label{sec:boundary-contacts-v157}
Let $R\subset M$ be a regular pencil and $\Lambda=\Pj(R)$. Its map
to $P$ lifts uniquely to a morphism
$\widetilde\iota_R:\Lambda\to\mathrm{CQ}(V)$: it lifts on the
nonsingular open, and properness extends it uniquely across the missing
points of the smooth curve. Put
$D_r(R)=\widetilde\iota_R^*E_r$. These are effective divisors because
the generic member is nonsingular.

\begin{theorem}[Contact-divisor reconstruction of spectral data]
\label{thm:contact-Segre-v157}
For every $h\geq1$, every regular pencil, and $0\leq p\leq nh$,
\begin{equation}\label{eq:contact-ideals-v157}
 J_p|_\Lambda=\OO_\Lambda\left(-\sum_{r=1}^{n-1}(p-hr)_+D_r(R)\right).
\end{equation}
At a point $x\in\Lambda$, let $0=e_1\leq\cdots\leq e_n$ be the
symmetric Smith exponents of the local pencil matrix. Then
\begin{equation}\label{eq:contact-exponents-v157}
 \operatorname{mult}_xD_r(R)=e_{r+1}-e_r\quad(1\leq r<n).
\end{equation}
Consequently the supported contact divisors determine every elementary
divisor, hence the full Segre symbol. If
$v_p(x)=\length(\OO_{\Lambda,x}/J_p)$ and $v_0(x)=0$, then
\begin{equation}\label{eq:contact-second-difference-v157}
 \operatorname{mult}_xD_r(R)=
 \frac{v_{h(r+1)}(x)-2v_{hr}(x)+v_{h(r-1)}(x)}{h}.
\end{equation}
\end{theorem}
\begin{proof}
Pull back \eqref{eq:power-boundary-v157}; equality of extended ideals
under composition gives \eqref{eq:contact-ideals-v157}.
The pencil matrix is nonzero at every point, since $R$ is a genuine
two-dimensional subspace; hence its entries generate the unit ideal
locally and $e_1=0$. Over $\C[[t]]$ symmetric elimination gives a
congruence to $\diag(t^{e_1},\ldots,t^{e_n})$, absorbing unit factors
by square roots. Its unique lift has, in
\eqref{eq:CQ-chart-v157}, $\tau_r=t^{e_{r+1}-e_r}$ up to a unit.
This proves \eqref{eq:contact-exponents-v157}; faithful flatness of
completion gives the local divisor statement. Taking valuations in
\eqref{eq:contact-ideals-v157} and second differences at multiples of
$h$ proves \eqref{eq:contact-second-difference-v157}.
Recover $e_i$ by summing the first $i-1$ contact multiplicities. The
positive exponents, with their points on $\Lambda$, are precisely
the elementary-divisor partitions of the spectral cokernel.
\end{proof}

\begin{example}[A boundary point and its contact order are distinct data]
\label{ex:contact-profile-v157}
For a symmetric germ with exponents $(0,1,3)$ and $h=2$, the valuations
of $J_1,\ldots,J_6$ are $(0,0,1,2,5,8)$.
Its two boundary contacts are $1$ and $2$. The germs
$\diag(1,t,t^3)$ and $\diag(1,t^2,t^5)$ reach the same full-flag
complete-quadric point but have contacts $(1,2)$ and $(2,3)$.
The compactification retains the flag; the arc with its divisorial
contacts retains the elementary exponents. These assertions must not
be conflated. For singular pencils the universal ideal theorem still
holds and detects all Fitting schemes and the normal rank; Fitting
ideals alone do not determine Kronecker minimal indices.
\end{example}


\section{A projective incidence compactification for pencils}
\label{sec:pencil-compactification-v157}
The complete-quadric graph concerns members of pencils. There is also
a projective parameter space carrying both the original finite failure
algebra and flat limits of their resolved member curves. This construction
keeps the pencil and does not identify distinct pencils at a common
full-coordinate-change limit.

\begin{proposition}[Resolved pencil curves and the failure family]
\label{prop:pencil-compactification-v157}
Let $n\geq3$, $G=\Gr(2,M)$, and let $G^\circ$ be the open set of
regular pencils whose lines avoid the rank-at-most-$(n-2)$ locus in
$P$. There is a normal projective $\operatorname{GL}(V)$-equivariant
scheme $\widehat G$ with a projective birational morphism
$\rho:\widehat G\to G$, an isomorphism over $G^\circ$, and a flat
family of subschemes
$\mathcal C\subseteq\widehat G\times\mathrm{CQ}(V)$ with the following
properties. Over $G^\circ$ it is exactly the lifted pencil line. Every
fibre over $R$ lies scheme-theoretically in $\pi^{-1}(\Pj(R))$.
All the resolved power maps restrict to $\mathcal C$, for every $h$.
The sharp finite failure algebra on $G$ pulls back to a finite locally
free algebra on $\widehat G$, with its original multiplication.
\end{proposition}
\begin{proof}
The rank-at-most-$(n-2)$ locus has codimension three in $P$, so a general
line avoids it; a general pencil is regular. Thus $G^\circ$ is nonempty
and dense. Over that open, $\pi$ is an isomorphism along every pencil
line and these lifts form a flat family of embedded projective lines.

For $1\leq q<n$, let $H_q$ be the pullback to $\mathrm{CQ}(V)$ of the
hyperplane bundle for its $q$th exterior-power projection, taking $q=1$
to mean $P$. The bundle $L=\bigotimes_{q=1}^{n-1}H_q$ is very ample,
by the graph embedding and Segre embedding. On a lifted line over
$G^\circ$, $\deg H_q=q$, since its $q$-minors have no common zero.
The Hilbert polynomial with respect to $L$ is therefore
$\binom n2 m+1$. The lifted family defines a morphism
\[
 G^\circ\longrightarrow
       \Hilb^{\binom n2 m+1}(\mathrm{CQ}(V)).
\]
Take the reduced closure of its graph in $G$ times this projective
Hilbert scheme, and then its normalization; this is $\widehat G$.
Normalization is finite here, so the result is projective and normal.
The graph is already closed over $G^\circ$ and that open is smooth,
so $\rho$ is an isomorphism there. Equivariance of the family and the
uniqueness of normalization give the stated group action.

Pull back the universal Hilbert family to obtain $\mathcal C$; it is
flat. To verify the asserted scheme-theoretic incidence, pull back
the linear equations for the universal pencil line to $\mathcal C$.
They vanish over the dense open $G^\circ$. On an affine open of the
integral base $\widehat G$, flatness makes the coordinate modules of
$\mathcal C$ torsion free over the base. A section which becomes zero
at the generic point is therefore zero. Applied locally to these
linear equations, this proves the incidence on the whole family.
The power maps are everywhere defined on $\mathrm{CQ}(V)$ by
Theorem~\ref{thm:power-graph-v157}, so their restrictions exist.
Finally pull back the finite locally free algebra of
Corollary~\ref{cor:finite-flat-family} along $\rho$; arbitrary base
change preserves its rank and multiplication.
\end{proof}

This is an incidence compactification over the Grassmannian, with its
native congruence action, not a properness assertion for the quotient
stack by that action. A limiting Hilbert fibre may have extra or
nonreduced components; no classification of all such fibres is used.
The construction is independent of $h$: its target, polarization and
Hilbert polynomial were specified using the $h$-independent complete
quadric graph. The unmarked inverse in the companion reconstruction
paper recovers the underlying pencil, so the incidence construction is
compatible with that inverse rather than a replacement of it.


\section{The simple corank-two incidence modification}
\label{sec:incidence-blowup-v160}
The ordinary-collision calculation describes specified families of
limits. We now identify the entire normalized Hilbert incidence space
on a natural open of the pencil Grassmannian. The result holds in every
dimension. It does not require a semisimple collision or a squarefree
residual pencil.

Let $n\geq3$, $M=\Sym^2V$, $P=\Pj(M)$, and $G=\Gr(2,M)$.
Write $Z_j\subset P$ for the symmetric rank-at-most-$j$ scheme, and
put $Z_0=\varnothing$ when $n=3$. Let $\mathcal U\subset G$ consist
of regular pencils $R$ such that their lines $L_R=\Pj(R)$ avoid
$Z_{n-3}$ and their scheme-theoretic intersections with $Z_{n-2}$
are either empty or a single reduced point. Write $\mathcal P$ for
the universal line over $\mathcal U$. We use projective spaces of
lines, so that the exceptional fibre of a smooth blow-up is the
projectivization of the normal space.

\begin{lemma}[The incidence centre and its normal space]
\label{lem:incidence-centre-v160}
The locus $\mathcal U$ is open and contains $G^\circ$ of
Proposition~\ref{prop:pencil-compactification-v157}. Its incidence
scheme
\[
 \mathcal D=\mathcal P\times_P Z_{n-2}\longrightarrow\mathcal U
\]
is a closed immersion with smooth image $\Sigma_2$ of codimension two.
At $R\in\Sigma_2$, let $p=L_R\cap Z_{n-2}$ and let
$a_R\subset N_{Z_{n-2}/P,p}$ be the nonzero image of $T_pL_R$.
There is a canonical isomorphism
\begin{equation}\label{eq:incidence-normal-v160}
 N_{\Sigma_2/\mathcal U,R}
       \simeq N_{Z_{n-2}/P,p}/a_R.
\end{equation}
Near $(p,R)$, after an etale change of base and coordinates, the
ideal of $\mathcal D$ in $\mathcal P$ is
\begin{equation}\label{eq:incidence-local-v160}
       I_{n-1}(A)=(x,u,v),\qquad I_{\Sigma_2}=(u,v),
\end{equation}
where $x$ is a relative line coordinate and $u,v$ are part of a
regular system of parameters on the base. The generators may be
changed by an invertible matrix; the asserted equality is an equality
of ideals.
\end{lemma}
\begin{proof}
On $P\setminus Z_{n-3}$ the scheme $Z_{n-2}$ is smooth of codimension
three. A nonsingular $(n-2)$-block reduces its equations to the three
entries of a symmetric $2\times2$ Schur complement. In particular
these are regular local equations for its given determinantal scheme
structure.

The universal line evaluation $\mathcal P_G\to P$ is smooth: over a
point of $P$ its fibre is the projective space of lines through that
point. Hence its inverse image of the smooth rank locus is smooth
and has dimension $\dim G-2$. Lines contained in the rank locus, or
meeting the lower rank locus, form closed exceptional sets after the
usual proper incidence projection. On the complementary open the
incidence projection is finite. The condition that its fibre have
length at most one is open, by upper semicontinuity for a finite
module. This proves the asserted openness. On this open the unit map
$\OO_{\mathcal U}\to q_*\OO_{\mathcal D}$ is surjective: on each
geometric fibre its target is either zero or a one-dimensional
algebra with unit, and Nakayama's lemma applies to its coherent
cokernel. Thus $q$ is a closed immersion. Its source is smooth by the
smooth evaluation just noted; its codimension is two.

First-order motion of a line evaluates surjectively onto
$T_pP/T_pL_R$: in the description
$T_RG=\Hom(R,M/R)$ this is restriction to the one-dimensional
subspace representing $p$. To remain incident to the rank locus,
this evaluation must vanish in $T_pP/(T_pZ_{n-2}+T_pL_R)$.
A reduced intersection means that the normal image of $T_pL_R$ is
nonzero. The kernel is precisely $T_R\Sigma_2$, since the incidence
point may move along both the line and the rank locus. The quotient
is therefore \eqref{eq:incidence-normal-v160}.

For the local assertion choose a rank-locus equation with nonzero
relative derivative. Its zero hypersurface is etale over the base
near $(p,R)$; after an etale base change it gives a section. Use that
equation as the coordinate $x$. Subtract multiples of $x$ from the
other two equations. Their restrictions to the section are base
functions $u,v$, and the ideal becomes $(x,u,v)$. The normal-space
calculation proves that $du,dv$ are independent. This proves
\eqref{eq:incidence-local-v160}. The argument is the smooth-centre
coordinate calculation, not an assertion that arbitrary graph
closures commute with base change.
\end{proof}

The open is nonempty along its centre. For example, a diagonal pencil
with two equal linear entries and all other entries having distinct
zeros away from their common zero belongs to $\Sigma_2$.
Although $a_R\ne0$, its residual binary quadratic form may have rank
one. This case is allowed throughout the following theorem.

\begin{theorem}[The full simple-incidence Hilbert modification]
\label{thm:incidence-blowup-v160}
For the normalized Hilbert incidence modification
$\rho:\widehat G\to G$ of
Proposition~\ref{prop:pencil-compactification-v157}, there is a
canonical isomorphism over $\mathcal U$
\begin{equation}\label{eq:incidence-blowup-v160}
 \widehat G\times_G\mathcal U
       \simeq B:=\operatorname{Bl}_{\Sigma_2}\mathcal U.
\end{equation}
It is equivariant under source congruence. In particular this part
of the incidence space is smooth. Its exceptional divisor
$F=\Pj(N_{\Sigma_2/\mathcal U})$ is a $\Pj^1$-bundle, and
\begin{equation}\label{eq:incidence-canonical-v160}
 K_B=\beta^*K_{\mathcal U}+F,
 \qquad \OO_B(F)|_{F_R}=\OO_{\Pj^1}(-1),
\end{equation}
where $\beta:B\to\mathcal U$ is the blow-up.

For every $R\in\Sigma_2$ the entire fibre of $\rho$ is $F_R\simeq
\Pj^1$. Its point $[\bar b]\in\Pj(N_{Z_{n-2}/P,p}/a_R)$ represents
the reduced nodal curve
\begin{equation}\label{eq:incidence-fibre-v160}
 C_R\cup\ell_{[\bar b]},\qquad
 \ell_{[\bar b]}=\Pj(a_R+\C b)
                  \subset\pi^{-1}(p)=\Pj(N_{Z_{n-2}/P,p}).
\end{equation}
Here $C_R$ is the strict transform of $L_R$ and $b$ is any lift of
$\bar b$. The tail attaches to $C_R$ at $[a_R]$. Every line through
$[a_R]$ in this exceptional $\Pj^2$ occurs, and no further components,
embedded points, or nonreduced fibres occur over this open. The two
components have exterior multidegrees
\begin{equation}\label{eq:incidence-multidegrees-v160}
 \deg_{C_R}H_q=q-\mathbf1_{q=n-1},\qquad
 \deg_{\ell}H_q=\mathbf1_{q=n-1}\quad(1\leq q<n).
\end{equation}
Their common Hilbert polynomial is $\binom n2m+1$ for
$\bigotimes_{q=1}^{n-1}H_q$.

The universal curve over $B$ is the simultaneous multiplication graph
on $\mathcal P\times_{\mathcal U}B$. Near a node it is etale locally
$xy=t$, where $t=0$ defines $F$. Thus it is a flat family of nodal
curves, including along directions for which the residual determinant
has a repeated zero. It commutes with base change as this specified
flat family. The blow-up is characterized by the universal property
that $I_{\Sigma_2}\OO_B$ is an invertible ideal; the corresponding
factorization property holds for every $T\to\mathcal U$ on which
that image ideal is an effective Cartier ideal.
\end{theorem}
\begin{proof}
On $P\setminus Z_{n-3}$, the complete-quadric morphism is exactly the
blow-up of the smooth centre $Z_{n-2}$: all earlier rank blow-ups
are isomorphisms there and the determinant divisor requires no
blow-up. We use the scheme-level classical construction
\cite[Remark~2.5 and Construction~2.6]{Massarenti2020}, originating
in \cite[Theorem~6.3]{Vainsencher1982}. Nothing in this step declares
the space of complete quadrics itself new.

Pull the universal line back to $B$. In a base blow-up chart
$u=t$, $v=tw$, the ideal \eqref{eq:incidence-local-v160} becomes
$(x,t)$. Its blow-up has the two charts
\[
        \OO_B[x,y]/(xy-t),\qquad \OO_B[z],\quad x=tz.
\]
They glue by the usual ratio change. Their total spaces are smooth,
and they are flat over $B$. For example
$\OO_B[x,y]/(xy-t)$ is free as an $\OO_B$-module with basis
$1,x,x^2,\ldots,y,y^2,\ldots$; reduction by the monic relation $xy-t$
gives this basis. Its fibres at $t=0$ have two reduced branches and
one ordinary node. The other base chart has the same description.
Outside the incidence section the pulled ideal is the unit ideal.
Blowing up the globally defined pulled ideal therefore gives a
projective flat family $\mathcal C_B$ with these local charts.
All smaller minor ideals are units on this open. The universal power
identity makes the product of the nontrivial power ideals an
$h^2$th power of $I_{n-1}$, up to an invertible determinant ideal.
Thus this same blow-up is the simultaneous multiplication graph
for every $h$, not only the exterior graph.

There is a closed immersion
$\mathcal C_B\hookrightarrow\mathrm{CQ}(V)\times B$. Indeed the
base-changed Rees algebra of the rank-locus ideal surjects onto the
Rees algebra of its image on the universal line; taking Proj gives
a closed immersion into the base-changed complete-quadric graph.
The universal line is already closed in $P\times B$. Composing these
two closed immersions proves the assertion. Equivalently, use the
simultaneous graph lemma on the integral total space. This argument
does not identify a general base change of a blow-up with the blow-up
of the image ideal without the indicated Rees quotient.

In normal coordinates at $p$, the exceptional map is
$[x:t]\mapsto[x a_R+t b]$, with $b$ representing the normal direction
of the base motion. Formula \eqref{eq:incidence-normal-v160} implies
that $a_R,b$ are independent. This is a line embedded in
$\Pj(N_{Z_{n-2}/P,p})$, not a multiple cover. The strict-transform
branch meets it at $[a_R]$. This proves
\eqref{eq:incidence-fibre-v160}, including the converse parametrization
of all lines through that point. The local charts show the absence
of embedded components and multiplicities.

The minor systems of orders $q<n-1$ have no base point on $L_R$.
The $(n-1)$-minors have one simple common zero, by
\eqref{eq:incidence-local-v160}. Their degrees on $C_R$ are therefore
$q$ and $n-2$, respectively. On the exceptional $\Pj^2$ only the
last exterior system has positive degree, namely one. This proves
\eqref{eq:incidence-multidegrees-v160}. The two projective lines
have Euler characteristic $2-1=1$, since their intersection is one
reduced point. Their total degree is $\binom n2$.

The flat embedded family defines a morphism from $B$ to the reduced
Hilbert graph closure over $\mathcal U$. It is proper because $B$ is
proper over $\mathcal U$ and the Hilbert graph is separated over that
base. Its image is exactly the graph closure: it is closed and contains
the dense graph over $\mathcal U\setminus\Sigma_2=G^\circ$, while
$B$ is integral. Over an incident pencil the displayed exceptional
lines are pairwise distinct as embedded curves; over a nonincident
pencil the fibre is a point. The morphism thus has finite geometric
fibres and is finite by \cite[Tag~02LS]{StacksProject160}. It is
birational, and $B$ is normal because its centre is smooth. A finite
birational map from a normal integral scheme identifies it with the
normalization: on affine charts its ring is the integral closure in
the common function field. Normalization commutes with restriction
to the open $\mathcal U$. This proves
\eqref{eq:incidence-blowup-v160}, without assuming that an arbitrary
nondominant normal test scheme lifts to a normalization.

All maps are specified by the universal incidence and agree on the
dense nonincident locus. Their uniqueness proves canonicity and
congruence equivariance. For completeness, in coordinates $v=uw$,
$du\wedge dv=u\,du\wedge dw$, giving the discrepancy one in
\eqref{eq:incidence-canonical-v160}. The exceptional normal bundle
is tautological of degree $-1$. The universal property is the usual
Rees-algebra factorization for an effective Cartier image ideal
\cite[Tag~0806]{StacksProject160}; here the preceding argument
identifies the actual Hilbert incidence modification with that
universal object. Flat pullback preserves the explicitly constructed
nodal family. No assertion about unrestricted graph-closure base
change is used.
\end{proof}

\subsection{Every power on every exceptional direction}
\begin{proposition}[Complete exceptional power systems]
\label{prop:exceptional-powers-v160}
Fix $R\in\Sigma_2$ and $p=L_R\cap Z_{n-2}$. For $h\geq1$ and
$1\leq j\leq hn$, define
\begin{equation}\label{eq:exceptional-delta-v160}
 \delta_j=(j-h(n-2))_+-2(j-h(n-1))_+.
\end{equation}
On the full exceptional plane $\pi^{-1}(p)\simeq\Pj^2$, the resolved
$j$th multiplication map is the complete degree-$\delta_j$ Veronese
map into its linear span, with degree zero interpreted as constant.
Its restriction to every tail in \eqref{eq:incidence-fibre-v160} is
therefore the complete degree-$\delta_j$ system on $\Pj^1$.
In particular $j=h(n-2)+1$ detects the tail as an embedded line.
These assertions include tangent lines to the residual determinant
conic and attachments on that conic.
\end{proposition}
\begin{proof}
A unit $(n-2)$-block leaves a residual two-dimensional space $W$.
Lemma~\ref{lem:block-products-v159} identifies the first nonzero
normal coefficient at degree $j>h(n-2)$ with the power of order
$b=j-h(n-2)$ in $B_h(W)$, multiplied by a nonzero unit-block socle.
For $j\leq h(n-2)$ the order-zero coefficient is constant on the
exceptional fibre. In the three-variable coordinate ring of
$\Sym^2W$, put $\mathfrak a=I_1$ and $f=\det$. The universal ideal
formula gives
\[
 J_b=\mathfrak a^b\quad(b\leq h),\qquad
 J_b=\mathfrak a^{2h-b}f^{b-h}\quad(h\leq b\leq2h).
\]
The generators all have degree $b$. Their span is the entire degree
$b$ part of the displayed ideal. After removal of the common factor
$f^{(b-h)_+}$ it is the full space of ternary forms of degree
$b-2(b-h)_+=\delta_j$. Hence it defines the complete Veronese system
on the exceptional plane. Restriction of all ternary forms to any
line gives all binary forms of that degree, independently of its
intersection multiplicities with the conic $f=0$. This proves all
assertions, including the endpoint $j=hn$ where the map is constant.
\end{proof}


\section{A computed boundary of the pencil incidence space}
\label{sec:computed-boundary-v158}
We give the incidence compactification a precise universal property and
compute a family of its boundary fibres in every dimension. The new
boundary is selected by one power base ideal, not by pulling an
unrelated algebra back to a parameter space.

\begin{proposition}[The normal main-component flattening property]
\label{prop:flattening-property-v158}
Let $G^\circ\subset G=\Gr(2,\Sym^2V)$ and
$\widehat G\to G$ be as in Proposition~\ref{prop:pencil-compactification-v157}.
Among normal integral schemes $T$ with a dominant morphism $f:T\to G$
for which the closure of the lifted pencil family over $f^{-1}(G^\circ)$
is an embedded flat family in $\mathrm{CQ}(V)\times T$ with Hilbert
polynomial $\binom n2 m+1$, there is a unique morphism $T\to\widehat G$
over $G$, and that family is the pullback of its universal family.
Consequently $\widehat G$ is the minimal normal dominant modification
with this specified embedded-flat-closure property.
\end{proposition}
\begin{proof}
The flat embedded family gives a morphism to the indicated Hilbert
scheme. Together with $f$ it agrees, over a dense open of $T$, with
the graph defining the reduced graph closure. Since $T$ is integral,
its morphism factors through this closed graph scheme: each equation
vanishing on the dense open vanishes on $T$. The induced morphism is
dominant. A dominant morphism from a normal integral scheme to an
integral scheme factors uniquely through the normalization: locally,
an element integral over the target ring becomes integral over the
source ring and lies in its fraction field, hence belongs to the
normal source ring. Thus the morphism factors uniquely through
$\widehat G$. The Hilbert-scheme representing property identifies
the families. Conversely its pulled-back universal family is flat;
flatness over an integral base makes it the schematic closure of its
restriction to the dense open. Uniqueness follows also from separatedness.
\end{proof}

The hypothesis of dominance is part of this property; this proposition
does not assert that normalization commutes with arbitrary base change.
The property is a standard consequence of the Hilbert functor and
normalization. Its usefulness here is that the following calculation
identifies a nontrivial fibre of the resulting modification.

Assume $n\geq3$. Choose distinct nonzero numbers $a_3,\ldots,a_n$ and put
$D=\Spec\C[[\tau]]$. On $\Pj^1_D$, consider the symmetric pencil
\begin{equation}\label{eq:collision-family-v158}
 A_\tau(s,t)=\diag\bigl(s,s+\tau t,
                           t-a_3s,\ldots,t-a_ns\bigr).
\end{equation}
For $n=3$ there is just one number $a_3$. Write $X=\Pj^1_D$ and let
$o$ be the point $s=\tau=0$ in the special fibre. Its generic pencil
has $n$ distinct rank-$(n-1)$ spectral points, so defines a point of
$G^\circ$.

\begin{theorem}[A transverse corank-two tail and its power maps]
\label{thm:hilbert-tail-v158}
The closure of the generic lifted curve of
\eqref{eq:collision-family-v158} in $\mathrm{CQ}(V)\times D$ is
\begin{equation}\label{eq:tail-blowup-v158}
                         Z=\operatorname{Bl}_o X.
\end{equation}
It is flat over $D$, and its special fibre is the reduced nodal union
$C\cup F$, where $C\simeq\Pj^1$ is the lifted special pencil and
$F\simeq\Pj^1$ is an exceptional tail. There are no embedded components.
If $H_q$ is the exterior-power hyperplane bundle of order $q$, then
\begin{equation}\label{eq:tail-degrees-v158}
 (\deg_C H_q,\deg_F H_q)=
 \begin{cases}(q,0),&1\leq q\leq n-2,\\
              (n-2,1),&q=n-1.
 \end{cases}
\end{equation}
Thus the special fibre has the same Hilbert polynomial
$\binom n2m+1$ as the generic lifted line.

The tail lies in $E_{n-2}$, its two intersections with $E_{n-1}$
are distinct and reduced, and its attachment to $C$ is neither of
those two points. For any $h\geq1$, let $L_p$ be the hyperplane
bundle of the resolved $p$th power map. Then
\begin{equation}\label{eq:tail-power-degree-v158}
 \deg_F L_p=
 \begin{cases}
 0,&p\leq h(n-2),\\
 \min\{b,2h-b\},&p=h(n-2)+b,\quad 0\leq b\leq2h.
 \end{cases}
\end{equation}
When this degree is positive, the restricted map on $F$ is the
complete Veronese map of that degree into its linear span.
Most directly, in the local coordinate $z=s/t$ at $o$,
\begin{equation}\label{eq:one-power-centre-v158}
 J_{h(n-2)+1}(A_\tau)=(z,\tau).
\end{equation}
Hence a single multiplication power ideal selects the blow-up centre
and its exceptional component.
\end{theorem}
\begin{proof}
The rank is at least $n-2$ on $X$. In a neighbourhood of $o$ the last
$n-2$ diagonal entries are units; all other spectral points have rank
$n-1$ and are pairwise distinct. Consequently the extended minor ideals
$I_q$ are the unit ideal for $q\leq n-2$, while
\[
 I_{n-1}=(z,z+\tau)=(z,\tau),\qquad
 I_n=(u z(z+\tau))
\]
near $o$, with $u$ a unit. Away from $o$, $I_{n-1}$ is a unit ideal;
$I_n$ is everywhere invertible. The simultaneous exterior-power graph
on the integral surface $X$ is therefore the blow-up of $(z,\tau)$.
This is an application of the graph-by-generators lemma, not a claim
that arbitrary base change of a blow-up is a blow-up. Since the
universal family of pencil lines embeds $X$ into $\Pj(\Sym^2V)\times D$,
this graph embeds into $\mathrm{CQ}(V)\times D$ and is the required
scheme-theoretic closure.

On the chart with $\tau=zv$, the special fibre has equation $zv=0$;
on the other chart, $z=\tau w$, it has equation $\tau=0$.
These charts prove flatness over the DVR, multiplicity one of both
components, their transverse intersection, and absence of embedded
components. The map on $C$ is the unique lift of the special pencil.
The $q$-minors of that pencil have degree $q$. They have no common
factor for $q\leq n-2$, and precisely one factor $s$ for $q=n-1$.
This proves their degrees on $C$. On $F$, all lower maps are constant,
and the $(n-1)$st exterior map is linear. This gives
\eqref{eq:tail-degrees-v158}. For $L=\bigotimes H_q$, the degrees
on $C$ and $F$ add to $\binom n2$. The normalization sequence of
$C\cup F$ subtracts one point, so its Euler characteristic is one
and its Hilbert polynomial is as stated. Thus it is the Hilbert
limit used by $\widehat G$; equivalently the generic morphism from
$D$ extends to $\widehat G$ by properness, with this universal fibre.

Use homogeneous coordinates $[\alpha:\beta]=[z:\tau]$ on $F$.
After removing the exceptional factor, the residual quadratic form
on the two-dimensional kernel is
$\diag(\alpha,\alpha+\beta)$. This is a line in
$\mathrm{CQ}(\C^2)=\Pj(\Sym^2\C^2)$, and gives the claimed image
of $F$ in the fibre of $E_{n-2}$. Its determinant $\alpha(\alpha+\beta)$ has two
simple zeros $[0:1]$ and $[-1:1]$, which are its intersections with
$E_{n-1}$. The main component attaches at $[1:0]$, where this residual
form is nondegenerate.

For clarity, the whole list of extended power ideals near $o$ is
\begin{equation}\label{eq:tail-all-ideals-v158}
 J_p=
 \begin{cases}
 \OO,&p\leq h(n-2),\\
 (z,\tau)^b,&p=h(n-2)+b,\ 0\leq b\leq h,\\
 (z,\tau)^{2h-b}[z(z+\tau)]^{b-h},
       &p=h(n-2)+b,\ h\leq b\leq2h,
 \end{cases}
\end{equation}
up to an invertible factor. These are actual ideals, by
Theorem~\ref{thm:universal-power-v157}. In particular $b=1$ gives
\eqref{eq:one-power-centre-v158} for every $h$.

It remains to check the linear system on the tail, not only its
degree. On diagonal matrices the coefficient polynomials of a $p$th
power span exactly the monomials of total degree $p$ in the diagonal
entries with each exponent at most $h$. For containment, use the
torus-weight bound in Lemma~\ref{lem:bounded-row-v157}; for the reverse
inclusion, take the coefficients of the independent diagonal entries
of $X$ in $\det(X+uA)^h$. Thus, for $p=h(n-2)+b$, the leading normal
coefficients use exponent $h$ in each of the $n-2$ invertible entries
and span
\[
  \alpha^j(\alpha+\beta)^{b-j},\qquad
           \max\{0,b-h\}\leq j\leq\min\{b,h\}.
\]
If $b\leq h$ these span every
binary form of degree $b$. If $b>h$, remove the common factor
$[\alpha(\alpha+\beta)]^{b-h}$; the remaining monomials span every binary form of
degree $2h-b$. Lower powers have nonzero constant normal term and
are constant on $F$. This proves the complete-linear-system assertion
and \eqref{eq:tail-power-degree-v158}.
\end{proof}

\begin{remark}[What the failure family contributes]
\label{rem:tail-integration-v158}
Along \eqref{eq:collision-family-v158}, the original finite failure
algebras form the finite locally free family already constructed in
Corollary~\ref{cor:finite-flat-family}. Its pullback to $F$ is constant;
we do not count that pullback as a new interaction. Instead, its first
relation canonically recovers the source and the pencil line, and the
source-normalized envelope has the total-base ideal
\eqref{eq:one-power-centre-v158}. Resolving that multiplication ideal
creates $F$, and the other products prescribe all the linear systems
\eqref{eq:tail-power-degree-v158} on it. A degeneration arc, not only
its closed failure algebra, supplies this total-base ideal. Thus the
calculation distinguishes information in the family from information
in the special point and supplies an explicit modular Hilbert limit.
\end{remark}


\section{Ordinary collision strata and their Hilbert fibres}
\label{sec:ordinary-collisions-v159}
We compute simultaneous collisions of arbitrary multiplicities, and
allow higher-order perturbations of the total family. Besides describing
its reduced components, we determine every resolved power system and
exhibit positive-dimensional families in a fibre of the pencil
incidence modification. The calculation uses the ordinary power ideals
on the two-dimensional total base.

\subsection{A local ideal calculation}
\begin{lemma}[Distinct tangent factors]
\label{lem:tangent-factors-v159}
Let $\ell_1,\ldots,\ell_c$ be pairwise nonproportional linear forms
in two variables. For $0\leq j<c$, the products of $j$ distinct
$\ell_i$ span all binary forms of degree $j$. Consequently, if
$R$ is a regular local complex surface ring with maximal ideal
$\mathfrak n=(z,\tau)$, and a symmetric $c\times c$ matrix $S$ satisfies
\[
 S=zC+\tau D+O(\mathfrak n^2),
 \qquad C\text{ nonsingular},\qquad
 \det(\alpha C+\beta D)\text{ squarefree},
\]
then $I_j(S)=\mathfrak n^j$ for $j<c$, while
$I_c(S)=(\det S)$.
\end{lemma}
\begin{proof}
Choose $j+1$ of the linear forms. Their products obtained by omitting
one factor are linearly independent: evaluation at the zero of the
omitted factor kills all other products and not that one. Their
number is $j+1$, the dimension of the required binary-form space.
For $j=0$ the assertion is immediate.

The endomorphism $C^{-1}D$ is self-adjoint for $C$ and has distinct
eigenvalues. Its eigenspaces are orthogonal and none is isotropic,
since otherwise $C$ would be degenerate. A constant congruence
therefore makes $\alpha C+\beta D$ diagonal with distinct linear
entries. Its $j$-minors span $\Sym^j(\mathfrak n/\mathfrak n^2)$.
The entries of $S$ lie in $\mathfrak n$, so
$I_j(S)\subset\mathfrak n^j$ and
$I_j(S)+\mathfrak n^{j+1}=\mathfrak n^j$. Nakayama's lemma applied
to $\mathfrak n^j/I_j(S)$ proves equality. The argument may be made
after completion and descended by faithful flatness. The determinant
assertion is its definition.
\end{proof}

\begin{lemma}[Block products and the first normal term]
\label{lem:block-products-v159}
For $V=U\oplus W$, multiplication gives an injection
\[
                 B_h(U)\otimes B_h(W)\longrightarrow B_h(V).
\]
If $\dim U=r$, $Q\in\Sym^2U$ is nonsingular, and
$S\in\mathfrak n\otimes\Sym^2W$ has a generically nonsingular linear term $S_1$, the first
nonzero normal term of $(Q+S)^p$, for $hr<p\leq h\dim V$, is
\[
                  \binom p{hr}Q^{hr}S_1^{p-hr}.
\]
For $p\leq hr$ its order-zero term $Q^p$ is nonzero. These statements
remain valid with a nonsingular formal unit block and after a formal
congruence, with the corresponding invertible changes of coefficients.
\end{lemma}
\begin{proof}
The Cartan relations on each summand vanish in $B_h(V)$, so the
algebra map is defined. Its source is Artin Gorenstein. Its socle is
the product of the two socles, and that product is nonzero in $B_h(V)$:
in determinant apolarity this is the nonzero mixed top derivative
of $\det(X_U)^h\det(X_W)^h$, obtained by restricting to block
diagonal matrices. A nonzero ideal in an Artin local Gorenstein
algebra meets its socle: successive multiplication of a nonzero
ideal by the maximal ideal ends in a nonzero subspace annihilated
by that maximal ideal. The kernel is therefore zero.

Expand $(Q+S)^p$ in the injected tensor product. Powers of $Q$
above $hr$ vanish, whereas $Q^{hr}$ is a nonzero socle generator.
For $p>hr$ the term with the fewest factors of $S$ consequently
has precisely $p-hr$ such factors and is the displayed term.
The unit block and a formal change of frame specialize to invertible
changes on the associated graded coefficient space. This proves the
last assertion as well.
\end{proof}

\subsection{The simultaneous collision theorem}
Let $D_0=\Spec\C[[\tau]]$, let $\mathcal R$ be a rank-two pencil
subbundle of $\Sym^2V\otimes\OO_{D_0}$, and put
$X=\Pj(\mathcal R)$. Assume its generic pencil has $n$ distinct
spectral points. Its special pencil is assumed regular. Apart from
simple spectral points, suppose it has distinct points
$x_1,\ldots,x_k$ of coranks $c_a$, where $2\leq c_a\leq n-1$.
Write $r_a=n-c_a$.

At $x_a$, choose a parameter $z_a$ on the special pencil and a unit
$r_a\times r_a$ block. Its symmetric Schur complement is required to
satisfy
\begin{equation}\label{eq:ordinary-cluster-v159}
 S_a=z_a C_a+\tau D_a+O((z_a,\tau)^2),\qquad
 C_a\text{ nonsingular},\qquad
 \det(\alpha C_a+\beta D_a)\text{ has }c_a\text{ simple zeros}.
\end{equation}
We call this an ordinary transverse cluster. Intrinsically,
$C_a,D_a\in\Sym^2W_a$, where
$W_a=(\ker(A(x_a):V^*\to V))^*$. Their pencil is the first derivative
of the form restricted to its radical. Taking a Schur complement does
not alter that first derivative. Replacing the parameter changes the
basis of this residual pencil and preserves its marked member $[C_a]$.
Thus the hypothesis does not depend on the chosen complement or frame.

\begin{theorem}[Simultaneous ordinary collision fibres]
\label{thm:ordinary-fibres-v159}
Under \eqref{eq:ordinary-cluster-v159}, the schematic closure of the
generic lifted pencil in $\mathrm{CQ}(V)\times D_0$ is
\begin{equation}\label{eq:ordinary-blowup-v159}
               Z=\operatorname{Bl}_{\{x_1,\ldots,x_k\}}X.
\end{equation}
It is regular and flat over $D_0$. Its special fibre is the reduced
nodal tree
\[
                          C\cup F_1\cup\cdots\cup F_k,
\]
where $C$ is the lift of the special pencil and each $F_a\simeq\Pj^1$
meets $C$ in one point. There are no embedded components and the tails
are pairwise disjoint. The tail $F_a$ is the complete-quadric lift of
the pointed residual pencil
$([\alpha C_a+\beta D_a],[C_a])$ on $W_a$. It lies in $E_{r_a}$,
meets $E_{n-1}$ in $c_a$ distinct reduced points, meets no other
boundary divisor, and attaches at the nonsingular member $[C_a]$.

Let $H_q$ be the hyperplane bundle of the resolved $q$th exterior map,
$1\leq q\leq n-1$. Then
\begin{equation}\label{eq:ordinary-multidegrees-v159}
 \deg_{F_a}H_q=(q-r_a)_+,\qquad
 \deg_C H_q=q-\sum_a(q-r_a)_+.
\end{equation}
For the product polarization $\bigotimes_{q=1}^{n-1}H_q$ the Hilbert
polynomial is $\binom n2 m+1$. In particular
$\deg F_a=\binom{c_a}{2}$ and
$\deg C=\binom n2-\sum_a\binom{c_a}{2}$.

For every $h\geq1$ and $1\leq p\leq hn$, set
\begin{equation}\label{eq:ordinary-power-numbers-v159}
 b_a=(p-hr_a)_+,\qquad s_p=(p-h(n-1))_+,\qquad
 \delta_{a,p}=b_a-c_a s_p.
\end{equation}
These integers are nonnegative. The resolved $p$th power restricts
to the complete linear system of degree $\delta_{a,p}$ on $F_a$.
It is constant if that degree is zero, and otherwise the complete
Veronese map into its linear span. Near $x_a$, with
$\mathfrak n_a=(z_a,\tau)$ and $f_a=\det S_a$, the exact ideals are
\begin{equation}\label{eq:ordinary-power-ideals-v159}
 J_p(A)=\mathfrak n_a^{\delta_{a,p}}(f_a)^{s_p},\qquad
                   J_{hr_a+1}(A)=\mathfrak n_a.
\end{equation}
Equalities are up to multiplication by a unit after a local frame
is fixed, not up to radical or integral closure. Consequently the
multiplication powers select the graph modification and every tail
linear system. This description depends only on the first normal jet
in \eqref{eq:ordinary-cluster-v159}; with the special pencil fixed,
higher-order terms do not change its embedded special fibre or its
power maps.
\end{theorem}
\begin{proof}
Near $x_a$ the unit block and its Schur complement give
$I_q(A)=\OO$ for $q\leq r_a$ and
$I_{r_a+j}(A)=I_j(S_a)$. Lemma~\ref{lem:tangent-factors-v159}
therefore gives
\[
 I_{r_a+j}(A)=\mathfrak n_a^j\quad(1\leq j<c_a),\qquad
 I_n(A)=(u_a f_a),\quad u_a\in\OO^\times.
\]
At every other point the rank is at least $n-1$, so all minor ideals
of order at most $n-1$ are units. The determinant ideal is invertible.
The simultaneous exterior graph is the blow-up of the product of
its base ideals by Lemma~\ref{lem:simultaneous-graph-v158}.
Near $x_a$ this product is
$\mathfrak n_a^{\binom{c_a}{2}}$. A positive power has the same
blow-up as the ideal itself, and the points are disjoint. These local
identifications glue uniquely over $X$ to \eqref{eq:ordinary-blowup-v159}.
This computes the graph on the total surface; it is not an assertion
that graph closure commutes with arbitrary base change.

The tautological family of genuine pencil lines embeds $X$ into
$\Pj(\Sym^2V)\times D_0$. Hence the simultaneous graph embedding
places $Z$ in $\mathrm{CQ}(V)\times D_0$, with exactly the prescribed
generic curve. The two charts $\tau=z_av$ and $z_a=\tau w$ of each
point blow-up show that $Z$ is regular. In the first chart its special
fibre is $z_av=0$. Both branches have multiplicity one and intersect
transversally. In the second chart only the exceptional branch occurs.
The parameter $\tau$ is a nonzerodivisor, so $Z$ is flat over the DVR.
The charts also rule out embedded components.

On $F_a$, divide the Schur block by the exceptional parameter.
Its value is $\alpha C_a+\beta D_a$. All earlier unit-block data
specialize to their values at $x_a$, independently of the direction.
The complete-quadric chart thus identifies $F_a$ with the lifted
residual pencil in the fibre over $[A(x_a)]$. Its only residual rank
drops have corank one and occur at the $c_a$ simple determinant zeros.
This proves the boundary assertions, including the attachment at
$[\alpha:\beta]=[1:0]$. For $q=r_a+j<n$, the leading $q$-minors
are the $j$-minors of this residual pencil. They have no common zero
and have degree $j$. The lower maps are constant. On $C$, the common
factor in the $q$-minors has order $(q-r_a)_+$ at $x_a$ and no other
zeros. This proves \eqref{eq:ordinary-multidegrees-v159}.
The sum of the component degrees is $\binom n2$. The normalization
sequence for a tree of $k+1$ projective lines subtracts its $k$ nodes,
so its Euler characteristic is one. This proves the Hilbert polynomial
and identifies the curve as the fibre in the Hilbert incidence space.

For the power ideals apply Theorem~\ref{thm:universal-power-v157}.
If $p\leq hr_a$, the ideal is the unit ideal. Otherwise write
$p=hr_a+hj+s$ with $0\leq s<h$. For $j<c_a-1$, the two minor
ideals are $\mathfrak n_a^j$ and $\mathfrak n_a^{j+1}$, so the
power ideal is $\mathfrak n_a^{hj+s}$. For $j=c_a-1$ it is
\[
        \mathfrak n_a^{(c_a-1)(h-s)}(f_a)^s.
\]
At $p=hn$ it is $(f_a)^h$. These cases are exactly
\eqref{eq:ordinary-power-ideals-v159}. The value $p=hr_a+1$ gives
the point ideal for every $h$.

It remains to identify the complete coefficient space on $F_a$.
By Lemma~\ref{lem:block-products-v159}, its first normal coefficients,
for $p>hr_a$, are exactly those of the $b_a$th power of the residual
pencil, multiplied by a nonzero unit-block socle vector. In the
polynomial ring $\C[\alpha,\beta]$, Lemma
\ref{lem:tangent-factors-v159} gives $I_j=(\alpha,\beta)^j$ for
$j<c_a$. The universal power identity gives the homogeneous ideal
\[
 (\alpha,\beta)^{\delta_{a,p}}
                    \det(\alpha C_a+\beta D_a)^{s_p}.
\]
All its original power generators have degree $b_a$; therefore their
span is its entire degree-$b_a$ part. Remove the common determinant
factor. The remaining space is exactly
$H^0(\Pj^1,\OO(\delta_{a,p}))$, proving the assertion about the
linear system. This argument determines the map as well as its degree.
All surviving leading coefficients use only the first normal jet;
the higher-order terms vanish on the exceptional divisor.
\end{proof}

\subsection{A modular family over a fixed pencil}
Let $R_0$ be a regular pencil whose spectral elementary divisors all
have exponent one, with repeated spectral multiplicities
$c_1,\ldots,c_k\geq2$. Thus these are semisimple collisions, not
nontrivial Jordan blocks. At $x_a$ keep the residual space $W_a$ and
the nonsingular tangent member $[C_a]$. Put
\begin{equation}\label{eq:residual-parameter-v159}
 U_a=\{\ell\text{ a line through }[C_a]\text{ in }\Pj(\Sym^2W_a):
             \ell\text{ meets the determinant in }c_a\text{ simple points}\}.
\end{equation}
This is a nonempty open in
$\Pj(\Sym^2W_a/\langle C_a\rangle)$, of dimension
$\binom{c_a+1}{2}-2$.

\begin{theorem}[Residual-pencil moduli in an incidence fibre]
\label{thm:boundary-moduli-v159}
The product $U=\prod_aU_a$ parametrizes an algebraic flat family of
distinct embedded Hilbert limits over the fixed pencil $R_0$.
Each is the fixed main curve with the tails prescribed by its
residual lines, and each is realized by an ordinary transverse
one-parameter smoothing of $R_0$. Hence, for the normal incidence
modification $\rho:\widehat G\to G$, one has
\begin{equation}\label{eq:boundary-fibre-dimension-v159}
                \dim\rho^{-1}(R_0)\geq
                    \sum_a\left(\binom{c_a+1}{2}-2\right).
\end{equation}
Distinctness is detected by the resolved multiplication power
$J_{h(n-c_a)+1}$ on the $a$th tail. The statement describes a specified
family in the fibre, not the entirety of the fibre or its normalization.
\end{theorem}
\begin{proof}
Since $C_a$ is nonsingular, lines through $[C_a]$ are parametrized
by the displayed projective quotient. The condition on determinant
zeros is open and nonempty, as is seen by choosing $D_a$ diagonal
with distinct eigenvalues relative to $C_a$.

The complete-quadric morphism is an isomorphism over forms of rank
at least $c_a-1$. Each line in $U_a$ stays in this open, and thus
its universal line has an algebraic lift. The fibre of
$\mathrm{CQ}(V)\to\Pj(\Sym^2V)$ over $[A(x_a)]$ contains the
residual $\mathrm{CQ}(W_a)$: the nondegenerate first quotient form
is fixed, and its radical carries the remaining complete form.
This follows also by taking the unit block in the Schur charts used
in the preceding proof. Thus the universal lifted lines give an
algebraic family of embedded tails, each passing through the fixed
attachment corresponding to $[C_a]$.

Adjoin the fixed main curve. The scheme-theoretic intersection of
each tail family with that main curve is just this reduced section.
Indeed the main curve has contact order one with $E_{r_a}$, since
its local nonzero Smith exponents are $0$ and $1$, whereas the tail
lies inside that divisor. The tails lie over different points, hence
do not intersect each other. The union exact sequence
\[
 0\longrightarrow\OO_{C\cup\bigcup F_a}
 \longrightarrow\OO_{C\times U}\oplus\bigoplus_a\OO_{F_a}
 \longrightarrow\bigoplus_a\OO_{\{\mathrm{attachment}_a\}\times U}
 \longrightarrow0
\]
proves flatness over $U$: the two terms on the right are flat, so
its kernel is flat. The degrees from
\eqref{eq:ordinary-multidegrees-v159} give the required fixed Hilbert
polynomial. We therefore obtain a morphism $U$ to the Hilbert scheme.
The embedded fibre over $[A(x_a)]$ recovers its residual tail. Its
first nonconstant exterior map, equivalently the power of index
$h(n-c_a)+1$, has degree one and recovers the line $\ell$ itself.
Distinct parameters give distinct embedded curves.

To see that these Hilbert points belong to the graph closure defining
$\widehat G$, choose a nonsingular member of $R_0$. Semisimplicity
makes its spectral eigenspaces orthogonal, so in a congruence frame
the pencil is a direct sum of blocks $(z-x_a)C_a$ and simple blocks.
Choose any representative $D_a$ of each residual line and replace
these blocks by $(z-x_a)C_a+\tau D_a$. Each cluster splits into
$c_a$ simple spectral points and distinct clusters stay separated.
Theorem~\ref{thm:ordinary-fibres-v159} identifies the Hilbert limit
with the constructed curve. This realizes every parameter by a
smoothing, independently on all the blocks.

Let $\overline G$ denote the reduced Hilbert graph before normalization.
The morphism from the smooth, reduced scheme $U$ factors through its
fibre over $R_0$: the defining equations vanish at every closed point
by the smoothing construction, hence vanish on $U$. Its geometric
fibres are zero-dimensional by the distinctness just proved; its
image consequently has dimension $\dim U$. Normalization
$\widehat G\to\overline G$ is finite and surjective, and the preimage
of this image has the same dimension. This gives
\eqref{eq:boundary-fibre-dimension-v159}.
We use this finite preimage, rather than asserting that a nondominant
map from a normal scheme must lift to a normalization. Such an
assertion would not be a valid normalization universal property.
\end{proof}

\begin{example}[Triple collision]
\label{ex:triple-collision-v159}
For one semisimple triple collision, $c=3$, the tail has exterior
multidegree $(0,\ldots,0,1,2)$ and total degree three. At $h=2$ its
successive residual powers $b=0,\ldots,6$ have degrees
\[
                             0,1,2,3,4,2,0.
\]
The residual lines through a fixed nonsingular ternary quadratic
form form an open in $\Pj^4$, and give a four-dimensional family of
distinct Hilbert limits over the same special pencil. Every member
has three distinct spectral points on its tail and one nonspectral
attachment. This is a family of embedded limits, not just a count
of branches of a normalization fibre.
\end{example}

The special finite failure algebra remains $A_{R_0}$ throughout this
family of limits. What varies is the total-family first-relation
geometry: its intrinsic envelope produces the maximal-ideal centre,
its successive products give the tail systems, and their degree-one
member recovers the residual line. Thus the family supplies information
not contained in its closed algebra alone. The construction is
equivariant under source congruence and under reparametrization of the
pencil, but the dimension bound concerns the embedded incidence space,
not a quotient by its stabilizer.


\section{Singularities of the universal power-zero schemes}
\label{sec:power-multiplier-v157}
There is an explicit further consequence for all the universal ideals,
which we record separately from the geometry of the full reciprocal
normalization fibre. Multiplier ideals are taken on the affine smooth
space $M$, and $I_{r+1}^{(a)}$ denotes symbolic power, with exponent
zero interpreted as the unit ideal.

\begin{corollary}[Multiplier ideals and threshold]
\label{cor:power-multiplier-v157}
For $1\leq p\leq nh$, put $w_r=(p-hr)_+$ and
$c_r=\binom{n-r+1}{2}$ for $0\leq r<n$. For every real $b\geq0$,
\begin{align}
 \mathcal J(M,b\cdot J_p)
   &=\bigcap_{r=0}^{n-1}
       I_{r+1}^{(\max\{0,\lfloor b w_r\rfloor+1-c_r\})},
                         \label{eq:power-multiplier-v157}\\
 \operatorname{lct}_M(J_p)
   &=\min_{\substack{0\leq r<n\\p>hr}}
                    \frac{\binom{n-r+1}{2}}{p-hr}.
                         \label{eq:power-lct-v157}
\end{align}
\end{corollary}
\begin{proof}
Blow up the affine origin and then the strict rank loci. This is the
affine complete-quadric resolution, whose charts are
\eqref{eq:CQ-chart-v157} with one additional scalar coordinate.
Let $E_0$ be the divisor over the origin. The same minor calculation
now includes $r=0$, and the total transform of $J_p$ is
$\sum_{r=0}^{n-1}w_rE_r$. The discrepancy along $E_r$ is $c_r-1$:
over the generic rank-$r$ point the earlier blow-ups are isomorphisms,
and the smooth rank-$r$ locus has codimension $c_r$; for $r=n-1$
this is the strict determinant divisor and the discrepancy is zero.
All boundary divisors have normal crossings, so this is a log resolution.

The defining pushforward formula for the multiplier ideal requires
$\operatorname{ord}_{E_r}(f)\geq\lfloor b w_r\rfloor+1-c_r$.
At the generic rank-$r$ point this valuation is the order in the
regular local ring with maximal ideal $I_{r+1}$, so the condition on
polynomials is exactly membership in the indicated symbolic power.
Intersecting the conditions proves the first formula; the first value
at which one becomes nontrivial proves the second. Equivalently this
is the specialization, via Theorem~\ref{thm:universal-power-v157}, of
the classical symmetric-determinantal multiplier-ideal calculation
\cite[\S2.4 and Theorem~4.8]{HenriquesVarbaro2016}.
\end{proof}

The ideal formula, graph identification and contact reconstruction are
statements about the entire universal coefficient space. They do not
supply the conductor of every multivariate common-divisor image or the
Hilbert function of every full $F_{g,k}(W)$. Those objects remain in
the preserved incidence theory; the present results connect their exact
quadratic sections to a native global boundary and its spectral schemes.


\section{Spectral divisor fibres and the recovery of degeneration type}
\label{sec:spectral-divisor-v154}
The preceding completion concerns the native congruence action. To
retain degeneration type before passage to a coarse quotient, we
associate a tower of binary divisor incidences to a regular pencil.
This construction is intrinsic after the unmarked inverse and is
unchanged by a simultaneous change of the pencil and source frames.

Let $\Lambda=\Pj(R)$ be the pencil line and let
$\mathcal T_R$ be the cokernel of the universal symmetric map on
$\Lambda$. A regular pencil means that its determinant is not
identically zero. Thus $\mathcal T_R$ is a torsion sheaf. Define
$D_j(R)$ by
\begin{equation}\label{eq:spectral-Fitting-v154}
 \II_{D_j(R)}=\Fitt_j(\mathcal T_R),\qquad 0\leq j\leq n,
\end{equation}
where $D_n=\varnothing$. These are effective divisors on the smooth
curve $\Lambda$. Write $s_j=\deg D_j$. If $s_j>0$, let $G_j$ be
its binary equation and put
\[
 J_j=G_j H^0(\Lambda,\OO_\Lambda(1))
       \subset H^0(\Lambda,\OO_\Lambda(s_j+1)).
\]
Changing the scalar representative of $G_j$ does not change $J_j$.
Let $\mathcal F_j(R)$ be the whole scheme fibre at $J_j$ of the
degree-one divisor normalization $\nu_1$, with $r=2$ and $D=s_j+1$.
For $s_j=0$ put $\mathcal F_j=\varnothing$.

\begin{theorem}[A fibre-theoretic Segre invariant]\label{thm:spectral-fibres-v154}
There is a natural identification of finite schemes over $\Lambda$
\begin{equation}\label{eq:fibre-divisor-v154}
                        \mathcal F_j(R)\simeq D_j(R).
\end{equation}
At a spectral point $x$, write the positive elementary divisors of
$\mathcal T_R$ as
$\lambda_{x,1}\geq\cdots\geq\lambda_{x,r_x}>0$, and extend this
sequence by zeroes. Then
\begin{equation}\label{eq:fibre-length-v154}
 (\mathcal F_j(R))_x=\Spec\C[\epsilon]/(\epsilon^{b_{x,j}}),
 \qquad b_{x,j}=\sum_{i>j}\lambda_{x,i},
\end{equation}
with the empty-scheme convention when $b_{x,j}=0$. Consequently
\begin{equation}\label{eq:Segre-difference-v154}
                 \lambda_{x,j+1}=b_{x,j}-b_{x,j+1}.
\end{equation}
The tower, with its maps to the same intrinsic spectral line,
therefore determines the full Segre symbol. Conversely, the
elementary-divisor partitions together with their spectral support
on that line determine the tower. The tower retains the projective
configuration of the support as well as the Segre symbol.
It is an invariant of the unmarked finite failure algebra.

The construction is compatible with base change on every family
stratum where the Fitting ideals in \eqref{eq:spectral-Fitting-v154}
define relative effective Cartier divisors of fixed degrees. On such
a stratum $\mathcal F_j$ is finite locally free of degree $s_j$.
These hypotheses are part of the relative statement; no assertion of
flatness across a jump of $s_j$ is intended.
\end{theorem}
\begin{proof}
Over the discrete valuation ring $\OO_{\Lambda,x}$, Smith normal
form gives diagonal entries with exponents
$\lambda_{x,r_x},\ldots,\lambda_{x,1}$ and $n-r_x$ zero
exponents. The minors of size $n-j$ have greatest common divisor
of order $\sum_{i>j}\lambda_{x,i}$. This proves the formula for
$D_j$, including its nonreduced structure. It is the usual
Fitting-ideal description of elementary divisors, used here on the
intrinsic spectral line.

The residual two-space $H^0(\OO_\Lambda(1))$ is gcd-free. Hence
Theorem~\ref{thm:binary-fibres-v153} identifies the degree-one
normalization fibre at $J_j$ with the scheme of linear factors of
$G_j$. On an affine chart write a monic linear factor as $z-a$.
Monic division gives $G_j=(z-a)H+G_j(a)$, so the entire factor
scheme is $G_j(a)=0$, not its reduction. These descriptions agree
on overlaps of the projective line. They give
\eqref{eq:fibre-divisor-v154} and the local algebra
\eqref{eq:fibre-length-v154}. Taking successive differences proves
\eqref{eq:Segre-difference-v154}. The common map to $\Lambda$ is
essential: forgetting it separately in each member of the tower
would discard the alignment of the spectral supports.

Pencil frame changes act on $\Lambda$ by projective transformations;
source congruence changes give isomorphic cokernels and preserve all
Fitting ideals. The intrinsic inverse of
Theorem~\ref{thm:artin-local-inverse-v146} therefore defines this
invariant from the unmarked algebra. The equivalence with Segre data
uses the classical elementary-divisor convention of
\cite[\S5]{FevolaMandelshtamSturmfels}.

For a family with the stated Cartier and degree conditions, the same
monic division argument works over any base algebra, including one
with nilpotents. It represents the divisor's functor of points, so
it commutes with base change. A relative effective divisor of fixed
degree on a projective-line bundle is finite locally free of that
degree. Local choices of $G_j$ differ by units and glue. Formation
of Fitting ideals itself commutes with base change; requiring them
to remain flat Cartier divisors is precisely what permits the
additional finite-flat conclusion.
\end{proof}

\begin{example}[Equal discriminant, different fibre towers]
\label{ex:segre-fibres-v154}
At a root of multiplicity four, the partitions $(3,1)$ and $(2,2)$
have the same zeroth fibre $\C[\epsilon]/(\epsilon^4)$.
Their first fibres are respectively $\C$ and
$\C[\epsilon]/(\epsilon^2)$; all higher fibres are empty there.
Thus the discriminant alone does not detect the distinction, whereas
the divisor-fibre tower does. Both are realized by regular symmetric
pencils: for a block of size $l$, take the reversal matrix $H_l$ and
the symmetric pair $(H_l,H_lJ_l(\alpha))$, where $J_l(\alpha)$ is
a Jordan block. Direct sums realize every indicated partition.
Additional simple blocks may be appended. The example is a statement
on the regular-pencil stack, not a claim that distinct Segre strata
always remain distinct after strictly semistable GIT identification.
\end{example}

\begin{remark}[The two divisor constructions]\label{rem:two-divisors-v154}
The primitive first relation of an actual pencil has exact common
factor $(\det T)^{n-1}$, by Lemma~\ref{lem:primitive-pencil}.
Its common-factor degree does not jump when the pencil changes
Segre type. The divisors $D_j$ above instead lie on the recovered
spectral line and define auxiliary binary first-relation subspaces
$J_j$. Thus \eqref{eq:Segre-difference-v154} relates pencil data to
intrinsically induced divisor ramification, not to a nonexistent
change in the conductor of the original exact-gcd locus. The
identification supplies a moduli consequence of the all-binary-fibre
theorem while respecting this distinction.
\end{remark}

% BEGIN PRESERVED PART 37-common-divisor-strata-v149.tex

\section{Common-divisor strata and intrinsic families}
\label{sec:common-divisor-strata}

The coefficient condition in Theorem~\ref{thm:recognition-dichotomy-v148}
is not imposed in this section.  We instead retain the common divisor
of a space of first relations.  This geometric condition gives a
larger smooth parameter space, includes relations without either
factor symmetry, and has a scheme-theoretic inverse.  The distinction
between a stratum and the condition on its geometric points is
essential when the parameter base is nonreduced.

Write \(S_j(E)=\Sym^j E\), where \(E\) is a complex vector space.
For \(r\ge2\), denote by
\[
 X_{h,r}(E)\subset\Gr(r,S_h(E))
\]
the locus of subspaces with greatest common divisor equal to \(1\),
up to a scalar.  Empty Grassmannians are omitted throughout.  When
\(h=0\) this notation can also be used for \(r=1\).

\begin{lemma}[Universal common divisor]\label{lem:universal-gcd}
Let \(D=g+h\).  The spaces of \(r\) degree-\(D\) forms whose exact
greatest common divisor has degree \(g\) form a locally closed
reduced subscheme \(Z_g\) of \(\Gr(r,S_D(E))\).  Multiplication
induces an isomorphism
\begin{equation}\label{eq:universal-gcd}
 \Pj(S_g(E))\times X_{h,r}(E)\ \xrightarrow{\ \sim\ }\ Z_g,
                 \qquad ([f],J)\longmapsto fJ.
\end{equation}
In particular, the common-divisor line and the residual subbundle of
a family with classifying map in \(Z_g\) are recovered uniquely and
commute with arbitrary complex base change.
\end{lemma}
\begin{proof}
For each \(a\), multiplication from
\(\Pj(S_a(E))\times\Gr(r,S_{D-a}(E))\) has closed image: its domain
is projective.  The union for \(a>g\) is closed.  Remove this union
from the image for \(a=g\), and give the resulting locally closed
set its reduced structure.  Its points are exactly those in the
statement.  The inverse image of this set under multiplication is
\(\Pj(S_g(E))\times X_{h,r}(E)\).  The same closed-image argument
with \(D=h\) shows that \(X_{h,r}(E)\) is open.

The restricted multiplication map is proper, by base change.  It
is injective on geometric points over every algebraically closed
extension of \(\C\), since a greatest common divisor is unique up
to scalar in a polynomial ring.  Its differential is injective as
well.  Indeed, a tangent vector in its kernel satisfies
\[
                   \dot f\, j+f\dot j\in fJ\quad(j\in J).
\]
Consequently \(f\mid\dot f\,j\) for every \(j\in J\).  For each
irreducible factor of \(f\), some element of \(J\) is not divisible
by that factor.  Comparison of its valuation gives
\(f\mid\dot f\).  Equality of degrees makes \(\dot f\) a scalar
multiple of \(f\); then \(\dot J=0\) in
\(\Hom(J,S_h(E)/J)\).

In characteristic zero the geometric-point injectivity and this
differential calculation imply that the map is universally
injective and unramified.  A proper, universally injective,
unramified morphism is a closed immersion
\cite[Lemma~41.7.2]{StacksClosedImmersion}.
Here its image is all of the reduced target.  The defining ideal
of a closed immersion with this support is zero in a reduced
scheme.  This proves the isomorphism, including its assertion for
nonreduced test schemes.  Equivalently, one can check properness
on the base change of the original projective multiplication map;
that base change has the same underlying open preimage.  Its source
is reduced because it is this open subscheme of the original
smooth product.
\end{proof}

\begin{remark}\label{rem:gcd-scheme-condition}
A map from a nonreduced scheme is required to factor through the
locally closed \emph{scheme} \(Z_g\), not merely to have geometric
points in its underlying set.  These are different requirements.
For example, over \(\C[\epsilon]/(\epsilon^2)\) the span of
\(x^2+\epsilon y^2,xy\) specializes to a pair with common divisor
\(x\), but has no factor \(x+\epsilon(ax+by)\).  Comparison of the
\(y^2\) coefficient in the first form proves this.  The stratum in
Lemma~\ref{lem:universal-gcd} excludes this transverse infinitesimal
family.  No use of a fibrewise radical can replace its scheme structure.
\end{remark}

Now let \(\dim V=\dim U=n\ge3\),
\(E=V^*\otimes U\), and \(\delta=\det T\).  Denote by
\(G\subset\operatorname{GL}(E)\) the stabilizer of the line
\(\C\delta\).  By Lemma~\ref{lem:linear-preserver-v147},
\begin{equation}\label{eq:determinant-group}
 G^\circ=(\operatorname{GL}(V)\times\operatorname{GL}(U))/\mathbb G_m,
                 \qquad G/G^\circ=\mathbf Z/2\mathbf Z.
\end{equation}
The second component exchanges the two tensor factors; the diagonal
scalars in the displayed product act trivially on the matrix space.
Fix \(s\ge1\) and \(D=sn+h\).  Inside \(Z_{sn}\), impose the condition
that the common-divisor line is the \(s\)-th power of a determinant
form after a linear change of the \(n^2\) variables.  Call the resulting
locally closed scheme \(Z^{\det}_{s,h,r}\).

\begin{theorem}[Intrinsic determinant-divisor stratum]
\label{thm:determinant-stratum}
There are canonical isomorphisms of schemes and algebraic stacks
\begin{equation}\label{eq:determinant-stratum}
 Z^{\det}_{s,h,r}\simeq
 \operatorname{GL}(E)\times^G X_{h,r}(E),\qquad
 [Z^{\det}_{s,h,r}/\operatorname{GL}(E)]\simeq[X_{h,r}(E)/G].
\end{equation}
They hold over arbitrary complex parameter schemes.  The stack is
smooth.  At a residual relation \(J\), its tangent complex, in degrees
\(-1,0\), is
\begin{equation}\label{eq:ambient-tangent-complex}
 [\,\operatorname{Lie}(G)\longrightarrow
                    \Hom(J,S_h(E)/J)\,],
             \qquad X\longmapsto(j\mapsto Xj\bmod J).
\end{equation}
No invariance of \(J\) under either full factor group is assumed.
The equivalence identifies stabilizers, infinitesimal families, and
specializations whose classifying maps remain in this stratum.
\end{theorem}
\begin{proof}
The power map \(\Pj(S_n(E))\to\Pj(S_{sn}(E))\),
\([f]\mapsto[f^s]\), is a closed immersion in characteristic zero.
It is proper and geometrically injective, and its differential has
kernel given by \(sf^{s-1}\dot f\in\C f^s\), hence by
\(\dot f\in\C f\).  The same closed-immersion criterion as above
applies.  The orbit of \([\delta]\) is the smooth locally closed
homogeneous space \(\operatorname{GL}(E)/G\).  The preimage of its
power orbit under the first projection in
\eqref{eq:universal-gcd} is therefore a locally closed subscheme.
It is exactly \(Z^{\det}_{s,h,r}\).

Pulling back along the \(G\)-torsor
\(\operatorname{GL}(E)\to\operatorname{GL}(E)/G\) trivializes the
common-divisor form, leaving the unrestricted residual relation
\(J\in X_{h,r}(E)\).  Descent gives the first isomorphism in
\eqref{eq:determinant-stratum}; descent of equivariant maps and
morphisms gives the second.  This proof uses morphisms of schemes,
not a comparison only on complex points.  It also supplies local
lifts of all isomorphisms over an fppf cover.

The space \(X_{h,r}(E)\) is open in a Grassmannian and \(G\) is
smooth.  The quotient is algebraic and smooth
\cite[Theorem~97.17.2]{StacksQuotient}; see also
\cite[Lemma~101.33.7]{StacksSmoothQuotient}.
For completeness, write a first-order residual subspace as the
graph of \(\epsilon\phi:J\to S_h(E)/J\).  An infinitesimal change
\(1+\epsilon X\) replaces \(\phi\) by \(\phi+X|_J\bmod J\).
This proves the tangent complex.  Its kernel records infinitesimal
automorphisms and its cokernel records first-order isomorphism
classes.  The same torsor description proves compatibility with
specialization.  No assertion is made about limits outside the
specified common-divisor stratum.
\end{proof}

The next step connects these parameter spaces with ungraded finite
algebras.  It is useful to state the family condition intrinsically,
since an absolute nilradical is unsuitable over a nonreduced base.
Put \(e=\dim E\) and
\(\ell=\binom{e+D}{D}-r\).  Consider finite locally free commutative
unital algebras \(\mathcal A\) of rank \(\ell\) on a complex scheme
\(T_0\), subject to the following conditions.  The map
\(\epsilon(a)=\ell^{-1}\operatorname{tr}(M_a)\) is multiplicative.
For \(\mathcal N=\ker\epsilon\), one has
\(\mathcal N^{D+1}=0\); all quotients of its power filtration are
locally free; \(\mathcal E=\mathcal N/\mathcal N^2\) has rank \(e\);
and multiplication gives isomorphisms
\(\Sym^j\mathcal E\simeq\mathcal N^j/\mathcal N^{j+1}\) for
\(j<D\).  In degree \(D\) its kernel has rank \(r\).  These
conditions specify the maximal-lower-Hilbert-function stratum; they
do not assert that arbitrary smoothings of a local algebra remain
in it.

\begin{proposition}[Relative first-relation algebra]
\label{prop:relative-first-relation}
Every such algebra is locally isomorphic, over its base, to
\[
 \mathcal A(\mathcal K)=
 \Sym\mathcal E/((\mathcal K)+\mathfrak m^{D+1}),\qquad
 \mathcal K=\ker(\Sym^D\mathcal E\longrightarrow\mathcal N^D).
\]
Its trace augmentation, \(\mathcal E\), and \(\mathcal K\) are
intrinsic and compatible with arbitrary base change in this
stratum.  For two such algebras, taking the linear part gives a
locally surjective map of isomorphism sheaves to the sheaf of
isomorphisms of \((\mathcal E,\mathcal K)\).  Its kernel is the
smooth substitution group whose underlying scheme in a local
presentation is
\(\Hom(\mathcal E,\mathcal N^2)\).

After quotienting morphisms by this tangent-identity kernel and
fppf stackifying, the algebras with first-relation map in
\(Z^{\det}_{s,h,r}\) have stack \([X_{h,r}(E)/G]\).
\end{proposition}
\begin{proof}
Locally lift a basis of \(\mathcal E\) to \(\mathcal N\).
Nilpotence and multiplication on the successive quotients make the
induced map from the truncated symmetric algebra surjective.
A relation with a nonzero term of least degree \(j<D\) would give
a nonzero kernel in degree \(j\) of the associated graded algebra,
which is impossible.  Its entire kernel therefore lies in degree
\(D\) and is precisely \(\mathcal K\).  Changing a lift by an
element of \(\mathcal N^2\) does not alter a product of \(D\)
generators, since the difference has degree at least \(D+1\).
This also proves the asserted uniqueness of the first relation.

In this presentation multiplication by a positive-degree element
is strictly filtration-increasing.  Its trace is zero over any
base ring.  Multiplication by a scalar has trace \(\ell\) times
that scalar, proving that the presentation's augmentation is the
specified intrinsic trace augmentation.  All filtration sequences
are locally split sequences of vector bundles, so their kernels
and quotients commute with base change.

An isomorphism must preserve trace and induces a linear map carrying
\(\mathcal K\) to \(\mathcal K'\).  Conversely such a map lifts
locally to an isomorphism by the homogeneous presentations.  All
other lifts are obtained by substituting generators plus arbitrary
elements of \(\mathcal N'^2\); the inverse is constructed degree by
degree in a finite filtration.  This gives the stated kernel and
its smoothness, without replacing its substitution law by addition.
After identifying morphisms with the same linear part, frames of
\(\mathcal E\) give the quotient
\([Z^{\det}_{s,h,r}/\operatorname{GL}(E)]\).
Theorem~\ref{thm:determinant-stratum} completes the proof.
\end{proof}

\begin{remark}\label{rem:algebra-inertia}
The original algebra stack is not this quotient stack: it has the
additional nonlinear inertia just computed.  The phrase
\emph{effective first-relation stack} will refer to the specified
quotient on morphisms followed by stackification.  The construction
retains the full automorphism group before taking that quotient.
It does not identify algebras with the same geometric fibres, nor
does it discard nilpotent parameters.  In particular, its relation
to a stratum of first relations is stronger than a bijection of
isomorphism classes over \(\C\).
\end{remark}

% BEGIN PRESERVED PART 40-canonical-divisor-boundary-v151.tex

\section{The common-divisor boundary and its normalization}
\label{sec:canonical-boundary-v151}

The exact-common-divisor stratum has a canonical completion which need
not be normal.  Its normalization records a divisor, and its failure to
be an isomorphism is controlled by two different phenomena: choices of
that divisor and overlap with the residual common factor.  We establish
these assertions for arbitrary spaces of forms.  No determinant, Cauchy,
or Pluecker equation is imposed in this section.

Let \(E\) have dimension \(e\geq2\), put \(S_j=\Sym^j E\) and
\(s_j=\dim S_j=\binom{e+j-1}{j}\), and fix
\[
             g,h\geq1,\qquad D=g+h,\qquad 2\leq r\leq s_h.
\]
Write \(\mathscr G=\Gr(r,S_D)\).  A \emph{degree-\(g\) divisor
incidence} over a complex scheme \(T_0\) consists of a line subbundle
\(\mathcal L\subset S_g\otimes\OO_{T_0}\) and a rank-\(r\)
subbundle \(\mathcal K\subset S_D\otimes\OO_{T_0}\) such that
\begin{equation}\label{eq:divisor-incidence-v151}
 \mathcal K\longrightarrow
 (S_D\otimes\OO_{T_0})/(\mathcal L\cdot S_h)=0.
\end{equation}
Subbundle always means that the quotient is locally free.  In
particular, the form generating \(\mathcal L\) is nonzero on every
geometric fibre.  Condition~\eqref{eq:divisor-incidence-v151} is the
vanishing of a bundle map, not a condition imposed only on closed points.

\begin{proposition}[Canonical divisor incidence]
\label{prop:incidence-v151}
The incidence functor is represented, with its full scheme structure, by
\[
 \mathscr X_g=\Pj(S_g)\times\Gr(r,S_h).
\]
Its forgetful morphism is multiplication
\begin{equation}\label{eq:normalization-map-v151}
 \nu_g:\mathscr X_g\longrightarrow\mathscr G,
                   \qquad ([f],J)\longmapsto fJ.
\end{equation}
This morphism is finite.  Its scheme-theoretic image \(Y_g\) is an
integral projective scheme.  If \(K\) is a geometric point of \(Y_g\)
and \(F_K=\gcd(K)\), the underlying set of \(\nu_g^{-1}(K)\) is
\begin{equation}\label{eq:divisor-fibre-set-v151}
       \{[f]: f\mid F_K,\ \deg f=g\},
\end{equation}
with residual space \(J=K/f\).  The incidence construction and its
scheme-theoretic fibres commute with arbitrary complex base change.
\end{proposition}
\begin{proof}
On \(\Pj(S_g)\), multiplication from the tautological line tensored
with \(S_h\) to \(S_D\) has constant rank \(s_h\).  Indeed,
multiplication by a nonzero polynomial over a field is injective.
A nonzero maximal minor locally splits this map, so its image and
cokernel are vector bundles, also after a nonreduced base change.
The closed incidence~\eqref{eq:divisor-incidence-v151} is therefore
the relative Grassmannian of rank-\(r\) subbundles of
\(\mathcal L\otimes S_h\).  Tensoring by \(\mathcal L^{-1}\)
identifies it canonically with the stated product.  In more elementary
terms, \(\mathcal K\) has a unique inverse image in
\(\mathcal L\otimes S_h\); its quotient there is locally free
because the quotient of \(S_D/\mathcal K\) onto
\(S_D/(\mathcal L\cdot S_h)\) is a surjection of vector bundles.
This proves the functorial assertion, including its scheme structure.

The source of \(\nu_g\) is projective.  Unique factorization in the
polynomial ring over an algebraically closed field gives
\eqref{eq:divisor-fibre-set-v151}.  A fixed polynomial has only finitely
many divisors of a fixed degree, up to scalar.  The morphism is thus
proper with finite fibres and is finite
\cite[Tag~02LS]{StacksFiniteV151}.
Define \(Y_g\) by the kernel of
\(\OO_{\mathscr G}\to(\nu_g)_*\OO_{\mathscr X_g}\), as in
\cite[Tag~01R5]{StacksImageV151}.  Its integrality follows from the
integrality of \(\mathscr X_g\): on an affine open meeting the image,
the coordinate ring of the image injects into the domain which is the
coordinate ring of its finite inverse image.  Reducedness is therefore
a consequence of this incidence construction, not an extra choice.
Representability of incidence gives the base-change assertion.
\end{proof}

\begin{remark}\label{rem:degree-divisor-v151}
The condition defining \(Y_g\) is the existence of a divisor of
\emph{degree \(g\)}.  In several variables it is not the condition
\(\deg\gcd(K)\geq g\): an irreducible polynomial of larger degree
need not have a degree-\(g\) factor.  Nor do we assert that formation
of the scheme-theoretic image, or normalization itself, commutes with
arbitrary nonflat base change.  The assertion for arbitrary bases in
Proposition~\ref{prop:incidence-v151} concerns the incidence functor
and the base change of its representing morphism.
\end{remark}

We next retain the nilpotents of a fibre.  This will distinguish a
ramified divisor choice from several distinct divisor choices.

\begin{lemma}[Fibre equations and the overlap space]
\label{lem:fibre-equations-v151}
Let \(K=fJ\), set \(H=\gcd(J)\), and put \(c=\gcd(f,H)\) and
\(k=\deg c\).  Then the tangent space of the scheme-theoretic fibre
at \(([f],J)\) is naturally, after the displayed representatives have
been chosen,
\begin{equation}\label{eq:overlap-kernel-v151}
 \ker(d\nu_g)_{([f],J)}\simeq S_k/\C c,
                     \qquad \dim\ker(d\nu_g)=s_k-1.
\end{equation}
More explicitly, choose a projective coefficient chart
\(f(u)=f+\sum u_\alpha m_\alpha\), with one coefficient normalized
to its value in \(f\).  Let \(M(f(u))\) be the matrix of
multiplication \(S_h\to S_D\).  Choose a set \(A_0\) of \(s_h\)
rows for which \(M_{A_0}(f)\) is invertible, and let \(B_0\) denote
the remaining rows.  For a matrix \(C\) whose columns form a basis
of \(K\), the full fibre on this chart is defined by
\begin{equation}\label{eq:fibre-schur-v151}
 C_{B_0}-M_{B_0}(f(u))M_{A_0}(f(u))^{-1}C_{A_0}=0,
             \qquad \det M_{A_0}(f(u))\ne0.
\end{equation}
In particular its completed local ring is the power-series ring in
the \(s_g-1\) variables \(u_\alpha\), modulo the entries in
\eqref{eq:fibre-schur-v151}; no radical is taken.
\end{lemma}
\begin{proof}
For each column of \(C\), its possible quotient by \(f(u)\) is forced
by the \(A_0\) rows to be \(M_{A_0}^{-1}C_{A_0}\).  The remaining
rows are exactly the condition that this quotient works.  These
quotients span a subbundle wherever the equations hold, by the
incidence argument above.  Thus elimination has introduced neither
additional solutions nor discarded nilpotents, proving the fibre
presentation over arbitrary coefficient rings in the chart.

A tangent vector is a pair \((\dot f\bmod\C f,\dot J)\).  It lies
in the fibre precisely when
\[
              \dot f\,j+f\dot j\in fJ\quad\hbox{for all }j\in J.
\]
It is necessary and sufficient that \(f\mid\dot f\,j\) for every
\(j\in J\); the class of \(\dot j\) modulo \(J\) is then forced.
Taking the minimum valuation over a basis of \(J\), for each
irreducible polynomial, makes this equivalent to \(f\mid\dot f H\).
Writing \(f=cf_1\), \(H=cH_1\), with \(\gcd(f_1,H_1)=1\), gives
\(\dot f=f_1w\) with \(w\in S_k\).  The scalar tangent direction
\(\dot f\in\C f\) is exactly \(w\in\C c\).  This proves
\eqref{eq:overlap-kernel-v151} and its dimension assertion.
\end{proof}

\begin{theorem}[Normalization and the exact normal locus]
\label{thm:boundary-normalization-v151}
The finite map
\[
 \nu_g:\Pj(S_g)\times\Gr(r,S_h)\longrightarrow Y_g
\]
is the normalization of the natural degree-\(g\) divisor locus.
Its dimension is
\begin{equation}\label{eq:boundary-dimension-v151}
                    \dim Y_g=s_g-1+r(s_h-r).
\end{equation}
For a geometric point \(K\), the scheme \(Y_g\) is normal at \(K\)
if and only if \(F_K\) has exactly one degree-\(g\) divisor \([f]\)
and
\begin{equation}\label{eq:normal-criterion-v151}
                         \gcd(f,F_K/f)=1.
\end{equation}
Normality at a point is also equivalent to smoothness there.
Thus the two precise obstructions are more than one divisor choice,
or overlap between the selected divisor and the residual common
factor.  The latter gives the fibre tangent dimension in
\eqref{eq:overlap-kernel-v151}, and
\eqref{eq:fibre-schur-v151} determines its higher nilpotents.
\end{theorem}
\begin{proof}
The locus of gcd-free \(r\)-planes in \(S_h\) is nonempty: a space
containing \(x_1^h,x_2^h\) is one example, and it can be extended
to every allowed rank \(r\).  It is open, since the loci possessing
a positive-degree divisor are the finitely many closed images of
projective multiplication maps.  It is therefore dense in the
irreducible Grassmannian.  Its product with \(\Pj(S_g)\) maps to
the locus where the exact gcd has degree \(g\).

This last locus is open in \(Y_g\): remove the finitely many closed
images for divisor degrees \(a>g\).  Call it \(U_g\).  Every fibre
over a geometric point of \(U_g\) has one point and, by
Lemma~\ref{lem:fibre-equations-v151} with \(H=1\), zero tangent
space.  A finite local algebra over an algebraically closed field
with zero cotangent space is the field itself, by Nakayama's lemma.
The fibre consequently has length one, not merely one point.

Here is the local argument which also settles the scheme structure.
At a closed point of \(U_g\), write \(A\) for the local ring of
\(Y_g\), and \(B\) for its finite inverse-image algebra.  The map
\(A\to B\) is injective, and \(B/\mathfrak m_A B=\C\), generated
by \(1\).  Nakayama's lemma, applied to the finite module \(B/A\),
gives \(B=A\).  The non-isomorphism locus of a finite morphism is
closed, so the same conclusion holds throughout \(U_g\).  Hence
\(\nu_g\) is birational.  The source is smooth and integral, thus
normal.  On affine opens, a finite birational extension \(A\subset B\)
with \(B\) integrally closed is the integral closure of \(A\) in
the common fraction field: every element integral over \(A\) is
integral over \(B\) and hence belongs to \(B\).  This proves the
normalization assertion; compare \cite[Tag~035E]{StacksNormalV151}.
The dimension follows from the product description.

If \(Y_g\) is normal at \(K\), its normalization is an isomorphism
there, so its fibre is one reduced point.  Formula
\eqref{eq:divisor-fibre-set-v151} gives uniqueness of \([f]\), and
\eqref{eq:overlap-kernel-v151} gives \(s_k-1=0\), equivalent to
\(k=0\) because \(e\geq2\).  Conversely uniqueness and \(k=0\)
again make the entire fibre a single reduced point.  The same finite
local Nakayama argument shows that the normalization is an isomorphism
near \(K\).  Its source is normal, so \(Y_g\) is normal at \(K\).
The source is smooth, so every normal point of the target is smooth;
the converse holds for every complex variety.  This proves both
directions without replacing a nonreduced fibre by its set of points.
\end{proof}

\begin{corollary}[Universality of the exact stratum]
\label{cor:canonical-exact-v151}
The open subscheme \(U_g\subset Y_g\) represents the unmarked
relation families in the divisor incidence whose residual space is
gcd-free on every geometric fibre.  Over it the divisor line and
residual subbundle are unique and commute with arbitrary base change.
Moreover \(U_g\), as a locally closed subscheme of \(\mathscr G\),
is the reduced exact-gcd stratum of Lemma~\ref{lem:universal-gcd}.
\end{corollary}
\begin{proof}
The inverse image of \(U_g\) is
\(\Pj(S_g)\times X_{h,r}(E)\), and the proof above shows that its
finite map to \(U_g\) is an isomorphism.  Its source is reduced,
so \(U_g\) has the reduced induced structure on its locally closed
support.  The incidence functor supplies the asserted universal
property.  A family in the ambient Grassmannian which only has the
right geometric gcd need not factor through \(Y_g\).  In particular,
the transverse dual-number family of
Remark~\ref{rem:gcd-scheme-condition} is not admitted by this
corollary.  The corollary is a scheme-theoretic identification, not a
claim that fibrewise gcd tests detect nilpotent equations.
\end{proof}

\begin{example}[A nonreduced fibre at a divisor collision]
\label{ex:fat-divisor-fibre-v151}
Take \(E=\langle x,y,z\rangle\), \(g=1,h=2,r=2\), and
\(K=\langle x^2y,x^2z\rangle\).  Its gcd is \(x^2\), so the only
point of the normalization fibre has divisor \(x\).  In the chart
\(f=x+ay+bz\), division by this monic linear form amounts to
substitution \(x=-ay-bz\).  The two remainders are
\[
             (ay+bz)^2y,\qquad (ay+bz)^2z.
\]
Their coefficient ideal is \((a^2,ab,b^2)\).  Thus the full fibre is
\begin{equation}\label{eq:fat-fibre-v151}
                     \Spec\C[a,b]/(a,b)^2,
\end{equation}
of length three and tangent dimension two.  It is not a reduced
point, despite uniqueness of the geometric divisor.  The family
\[
 K_t=x\langle xy+t y^2,\ xz+t z^2\rangle
\]
has independent generators for every \(t\), exact gcd \(x\) when
\(t\ne0\), and the displayed collision at \(t=0\).  Consequently
this nonreduced fibre occurs on the boundary of the exact stratum,
not in an unrelated example.
\end{example}

\begin{example}[Several branches and a normal gcd jump]
\label{ex:branches-normal-jump-v151}
For the same \(g,h,r,E\), let
\(K=xy\langle x,z\rangle\).  The divisors \(x\) and \(y\)
give two points of the normalization fibre.  At both points the
overlap is one, so the fibre consists of two reduced points.
The point of \(Y_1\) is nonnormal.  By contrast, put
\(q=x^2+yz\) and \(K=xq\langle y,z\rangle\), now with \(h=3\).
The polynomial \(q\) is irreducible and is not divisible by \(x\).
There is exactly one linear divisor of \(\gcd K=xq\), and its
overlap with the residual gcd \(q\) is one.  This point of \(Y_1\)
is normal although its gcd degree is three.  These two examples
show why neither gcd degree nor the number of geometric lifts alone
is a sufficient normality criterion.
\end{example}

\begin{corollary}[An intrinsic boundary for first-relation algebras]
\label{cor:algebra-boundary-v151}
On the maximal-lower-Hilbert-function stratum of
Proposition~\ref{prop:relative-first-relation}, the algebras whose
intrinsic degree-\(D\) kernel has classifying map in \(Y_g\)
have, after removing only tangent-identity substitutions, stack
\([Y_g/\operatorname{GL}(E)]\).  Its normalization is the finite
representable morphism
\begin{equation}\label{eq:quotient-normalization-v151}
 [\mathscr X_g/\operatorname{GL}(E)]
                    \longrightarrow[Y_g/\operatorname{GL}(E)].
\end{equation}
All members of the universal relation family give finite locally
free algebras of rank \(\binom{e+D}{D}-r\).  This construction imposes
no pencil image equations and permits the divisor and residual
common factor to collide.
\end{corollary}
\begin{proof}
For every subbundle \(\mathcal K\) classified by the Grassmannian,
\(\Sym E/((\mathcal K)+\mathfrak m^{D+1})\) is, as a module,
the sum of \(S_j\) for \(j<D\) and \(S_D/\mathcal K\).
This proves finite local freeness and the rank.  The first-relation
proposition identifies the quotient after the specified nonlinear
inertia is removed.  Incidence and multiplication are
\(\operatorname{GL}(E)\)-equivariant.  Pulling
\eqref{eq:quotient-normalization-v151} back along the smooth frame
atlas \(Y_g\to[Y_g/\operatorname{GL}(E)]\) gives precisely
\(\nu_g\).  Finiteness, representability, and the normalization
identification follow by descent; equivalently normality and integral
closure may be checked on these smooth atlases.  This quotient is
not claimed to classify smoothings outside the specified Hilbert
stratum.
\end{proof}

For the pencil first relations one can take
\(e=n^2\), \(g=n(n-1)\), \(h=3n-4\), and
\(r=\binom{\binom{n+1}{2}}2\).  Every pencil then lies in the
exact-gcd open of this canonical boundary.  The boundary theorem
concerns the much larger space of all degree-\(g\) divisorial first
relations, rather than a closure defined by recovering a pencil and
reimposing its Pluecker equations.  It determines its normalization
and normal locus, but does not assert a classification of all its
linear orbits or of all finite-flat algebra deformations.

% BEGIN NEW MATHEMATICS v153

\section{Formal images and marked divisor germs}
\label{sec:formal-images-v153}
All rings and schemes in the following sections are over $\C$. Write
$S_j=\Sym^j E$, and let $Y_g\subset\Gr(r,S_D)$ be the
scheme-theoretic image of
\[
 \nu_g:\Pj(S_g)\times\Gr(r,S_{D-g})\longrightarrow\Gr(r,S_D),
 \qquad (f,J)\longmapsto fJ,
\]
where $\dim E\geq2$, $g\geq1$, and $2\leq r\leq\dim S_{D-g}$.
The source is smooth, and $\nu_g$ is its finite normalization: a fibre
records the degree-$g$ divisors of the common divisor, and on the
exact-degree-$g$ locus there is one reduced lift. We shall retain the
scheme structures of these fibres and of the image.

\begin{lemma}[Completed finite images]\label{lem:formal-image-v153}
Let $X\to M$ be a finite morphism of complex finite-type schemes, with
$X$ reduced and $M$ smooth, and let $Y$ be its scheme-theoretic image.
Fix a closed point $y\in Y$, put $R=\widehat{\OO}_{M,y}$, and write
$B_i=\widehat{\OO}_{X,x_i}$ for the points over $y$. Then
\[
 \widehat{\OO}_{Y,y}=R/\bigcap_i\ker(R\to B_i).
\]
This completed image is reduced. If the cotangent map from $R$ to
$B_i$ is surjective, then $R\to B_i$ is surjective. If there are two
such points, writing $I_i=\ker(R\to B_i)$ gives
\[
 \widehat{\OO}_{Y,y}=(R/I_1)\mathbin{\times}_{R/(I_1+I_2)}(R/I_2).
\]
The formal schemes at the different $x_i$ are disjoint. A local Artin
point of the finite inverse image belongs to exactly one of them.
\end{lemma}
\begin{proof}
On an affine neighbourhood the image ring is the image of the
ambient ring in a finite algebra. Exactness of completion on finite
modules preserves this kernel. The completed finite algebra is the
product of the $B_i$: the orthogonal idempotents of its Artin closed
fibre lift uniquely through the powers of the maximal ideal. Reduced
finite-type complex rings are excellent and analytically unramified,
so their completions are reduced; alternatively this applies directly
to the reduced image ring. For cotangent surjectivity, lift a basis
of the maximal ideal of $B_i$ modulo its square to $R$. The resulting
lifts generate every successive associated-graded piece. Subtracting
successive homogeneous lifts and taking their convergent sum expresses
any element of $B_i$ as the image of an element of $R$. The finite map
factors through the completed image, so this is also a closed immersion
into that image, not only into $M$. Finally the elementary exact sequence
\[
0\longrightarrow R/(I_1\cap I_2)\longrightarrow R/I_1\oplus R/I_2
 \longrightarrow R/(I_1+I_2)\longrightarrow0
\]
proves the fibre-product assertion. The idempotent decomposition gives
the last assertion, including over nonreduced Artin bases.
\end{proof}

\begin{lemma}[Coprime lifts and the split codimension]
\label{lem:coprime-lifts-v153}
If homogeneous forms $a,b$ have coprime reductions over a local Artin
base $A$, then $(a)\cap(b)=(ab)$ in $A[E]$. The relevant graded
quotients and Koszul kernels are $A$-flat. In the split-divisor model,
with $s_j=\binom{e+j-1}{j}$, $e\geq2$, $h>g\geq1$, and $r\geq2$,
\[
 c=r(s_h-s_{h-g})-(s_g-1)\geq(r-1)(s_g-1)\geq1.
\]
\end{lemma}
\begin{proof}
Over the residue field the two forms form a regular sequence. In
each degree their Koszul complex consists of finite free modules and
has its prescribed ranks. The acyclic-complex criterion over an Artin
local ring, proved by lifting splittings successively along its maximal
ideal, gives exactness and flatness of the lifted cokernels. Thus $b$
is a nonzerodivisor modulo $a$, proving the intersection identity.
For the inequality, express $s_h-s_{h-g}$ as the sum of $g$ consecutive
first differences of the sequence $s_j$. These differences are
nondecreasing for $e\geq2$. Their sum is at least $s_g-s_0=s_g-1$.
\end{proof}

In particular the two formal marked lifts in a coprime split model
are separated before their images are compared. Their intersection
functor is precisely $K_A=a_A b_A L_A$. Division gives a subbundle:
multiplication by a fibrewise nonzero form is locally split, and the
quotient of the ambient bundle by $K_A$ is locally free. Combining
this Artin-functor identification with Lemma~\ref{lem:formal-image-v153}
justifies the completed-image fibre product without any assumption
that a tangent-space comparison determines an intersection scheme.


\section{A global factorization theorem for binary normalization fibres}
\label{sec:binary-fibres-v153}
For this section $E$ has dimension two. Set $N_0=r(D+1-r)$ and
$m=r-1$. Let
\[
 \mathcal T_s=Y_s\setminus Y_{s+1},\qquad 1\leq s\leq D+1-r,
\]
with the reduced locally closed structure; an inadmissible $Y_j$ is
empty. Binary forms split into linear factors, so the $Y_j$ are
nested. The multiplication map identifies
\[
 \mathcal T_s\simeq\Pj(S_s)\times X_{D-s,r},
\]
where $X_{D-s,r}$ is the open Grassmannian of gcd-free residual spaces.
The first projection is the algebraic common-divisor morphism, not a
choice made independently on individual fibres.

For nonnegative $a,b$ define $H_{a,b}=\C$ when $ab=0$. Otherwise,
put $k=\min(a,b)$ and $l=\max(a,b)$, and set
\begin{equation}\label{eq:H-ab-v153}
 H_{a,b}=\C[q_1,\ldots,q_k]/(c_{l+1},\ldots,c_{l+k}),\quad
 \sum_{j\geq0}c_jz^j=(1+q_1z+\cdots+q_kz^k)^{-1}.
\end{equation}
The weight of $q_i$ is $i$. This grading is not the maximal-ideal
adic grading when $k>1$.

\begin{theorem}[Stratified finite flatness and all binary fibres]
\label{thm:binary-fibres-v153}
For $g\leq s$, the restriction of $\nu_g$ over $\mathcal T_s$ is the
base change, along the common-divisor morphism, of
\[
 \Pj(S_g)\times\Pj(S_{s-g})\longrightarrow\Pj(S_s),\qquad (f,H)\mapsto fH.
\]
It is finite locally free of degree $\binom{s}{g}$. If
$\gcd(K)=\prod_{i=1}^{b}\ell_i^{s_i}$, with distinct $\ell_i$ and
$\sum_i s_i=s$, its entire normalization fibre is
\begin{equation}\label{eq:binary-fibre-product-v153}
 \prod_{\substack{0\leq a_i\leq s_i\\\sum_i a_i=g}}
                   \bigotimes_{i=1}^{b} H_{a_i,s_i-a_i}.
\end{equation}
Every local factor is an Artin complete intersection and is Gorenstein.
Its length, weighted Hilbert series, and socle weight are, respectively,
\begin{gather*}
 \prod_i\binom{s_i}{a_i},\qquad
 \prod_i\prod_{j=1}^{\min(a_i,s_i-a_i)}
       \frac{1-z^{\max(a_i,s_i-a_i)+j}}{1-z^j},\qquad
 \sum_i a_i(s_i-a_i).
\end{gather*}
The total length is $\binom{s}{g}$, including at arbitrary simultaneous
root collisions. Over fixed $s$, the closure order of common-root
partition strata is the order generated by merging parts.
\end{theorem}
\begin{proof}
Over $\mathcal T_s$ the common divisor $G$ and residual subbundle $L$
are defined scheme-theoretically by the inverse to the exact-divisor
incidence. Locally on any base change, choose a member of $L$ coprime
to $G$ on each relevant geometric fibre. There is such a member over
$\C$, since only finitely many roots must be avoided. Regular-sequence
cancellation, or monic division and its flat quotient, gives
$f\mid GL$ if and only if $f\mid G$. This equivalence is an equality
of incidence functors over nonreduced bases. The residual subspace is
then forced. This proves the base-change description.

The displayed multiplication map is finite by properness and the
finiteness of the divisors of a fixed binary form. Its source and
target are smooth of dimension $s$. A finite map between these smooth
spaces is flat: a regular system of parameters in the target is a
system of parameters in each source local ring, hence a regular
sequence; the local flatness criterion applies. Its degree can be
computed at a squarefree form, where the $\binom{s}{g}$ divisors are
reduced and distinct.

Coprime root clusters split uniquely over Artin bases by Hensel
factorization. On the cluster with polynomial $x^{s_i}$, choosing a
degree-$a_i$ factor is the equation $f_i H_i=x^{s_i}$ with both factors
monic. Eliminating one factor by reciprocal expansion gives
\eqref{eq:H-ab-v153}. This proves the full product decomposition,
not only its closed points. Its equations have their sole common zero
at the origin: a monic factor of $x^{a+b}$ over $\C$ is a power of
$x$. Thus its $k$ equations in $k$ variables are a homogeneous system
of parameters and a regular sequence. The Hilbert series is the
quotient of the corresponding weighted polynomial Hilbert series.
Taking its value at one gives $\binom{a+b}{a}$; its top weight is
$ab$. A zero-dimensional complete intersection has a one-dimensional
socle. Equivalently $q_k^l$, up to a nonzero scalar, represents the
rectangular Schur class. The identification with the usual
Grassmannian cohomology presentation is classical; a purely algebraic
Schur-basis statement is \cite[Theorem~2.7]{Grinberg2019}.
Tensoring the local factors proves all three formulas. The total
length identity is the coefficient of $z^g$ in
$\prod_i(1+z)^{s_i}$. Finally the exact-gcd stratum is the product
with $X_{D-s,r}$ above. In the symmetric power of the projective line,
root clusters specialize by coalescence, and every merging is realized
by moving the corresponding roots together. This gives precisely the
stated closure order within fixed gcd degree.
\end{proof}

\begin{remark}
The finite flat statement is on the exact-gcd stratum with its indicated
scheme structure. It does not assert that normalization or formation of
scheme-theoretic images commutes with an arbitrary nonflat base change
of the ambient Grassmannian. It identifies a specific base change of
the already constructed incidence morphism.
\end{remark}


\section{Conductors and all branch intersections on the squarefree locus}
\label{sec:squarefree-v153}
Let $\mathcal U\subset\Gr(r,S_D)$ be the open set on which the common
divisor is squarefree. Its complement is the proper image of the
incidence $K=\ell^2L$. This condition concerns the common divisor,
not every member of the subspace.

\begin{theorem}[Squarefree conductor stratification]
\label{thm:squarefree-v153}
Let $K\in Y_g\cap\mathcal U$, let $s=\deg\gcd(K)$, and set
\[
 t=s-g,\qquad m=r-1,\qquad q=N_0-ms.
\]
Introduce $s$ disjoint blocks $u_i=(u_{i1},\ldots,u_{im})$, and let
$J_j$ be the ideal generated by products of one variable from each of
$j$ distinct blocks. Put $J_0=(1)$ and $J_{s+1}=0$. Then
\begin{equation}\label{eq:squarefree-ring-v153}
 \widehat{\OO}_{Y_g,K}\simeq
 \C[[z_1,\ldots,z_q,u_{11},\ldots,u_{sm}]]/J_{t+1}.
\end{equation}
The branches are indexed by the subsets $A\subset\{1,\ldots,s\}$
of cardinality $g$. The branch $A$ has ideal
$I_A=(u_{ij}:i\in A)$ in the ambient power-series ring. For any
collection of branches, their full scheme-theoretic intersection has
ideal $I_{\cup A}$; it is smooth of dimension
$q+m(s-|\cup A|)$.

The conductor in \eqref{eq:squarefree-ring-v153} is $J_t/J_{t+1}$.
Its restriction to branch $A$ is
$\prod_{i\notin A}(u_{i1},\ldots,u_{im})$. Consequently the global
conductor sheaf on $Y_g\cap\mathcal U$ is exactly
\begin{equation}\label{eq:global-conductor-v153}
 \mathfrak c_{Y_g}|_{\mathcal U}
        =\II_{Y_{g+1}/Y_g}|_{\mathcal U}.
\end{equation}
In particular it is radical there. These local rings are seminormal.
Their dimension, depth and multiplicity are
\[
 q+mt,\qquad q+t,\qquad \binom{s}{g}.
\]
They are Cohen--Macaulay exactly when $s=g$ or $r=2$.
The Hilbert series of the nonsmooth factor in its standard grading is
\begin{equation}\label{eq:block-hilbert-v153}
 \sum_{j=0}^{t}\binom{s}{j}\bigl((1-z)^{-m}-1\bigr)^j.
\end{equation}
The whole collection of branch intersections and conductor strata
is invariant under permutation of the roots and descends from any
ordered-root etale chart.
\end{theorem}
\begin{proof}
Choose $F_0\in K$ with a simple zero at each of the $s$ common roots,
and complete it to a frame $F_0,\ldots,F_m$. Hensel factorization gives
moving roots $x_1,\ldots,x_s$ in the completion. Use as coordinates
the $s$ root parameters and the $ms$ evaluations
$u_{ij}=F_j(x_i)$, with harmless invertible normalizing factors.
These are part of a regular parameter system: the evaluation map
$S_D/K\to\C^s$ is surjective because $K$ vanishes on the common-root
scheme, and $D\geq s-1$. Thus
$\Hom(K,S_D/K)$ allows independent variation of all $rs$ jets.
The passage to the displayed coordinates is triangular with invertible
diagonal, since $F_0$ has simple zeros. The other $N_0-rs$ parameters
and the $s$ root parameters give exactly $q=N_0-ms$ smooth parameters.

A marked degree-$g$ divisor near this fibre chooses a subset $A$ of
these roots. Its incidence equations are exactly $u_{ij}=0$ for
$i\in A$, with no nilpotents suppressed. There are no other divisor
lifts. Lemma~\ref{lem:formal-image-v153} therefore identifies the
completed image ideal with $\bigcap_{|A|=g}I_A$. A monomial belongs
to this intersection exactly when its block support meets every
$g$-element subset, or equivalently has at least $s-g+1=t+1$ blocks.
This proves \eqref{eq:squarefree-ring-v153}. Sums of the coordinate
ideals give all intersections, including intersections of three or
more branches.

For completeness compute the conductor inside the product of the
branch rings. Its component on branch $A$ consists of functions whose
extension by zero on every other branch belongs to the image ring.
The image of $\bigcap_{B\ne A}I_B$ modulo $I_A$ is precisely
$\prod_{i\notin A}(u_i)$. Indeed a surviving monomial has support
outside $A$, and must meet every $g$-subset other than $A$. Replacing
one element of $A$ by each $i\notin A$ forces all those blocks to
occur. Conversely their product meets every such subset. Hence a
conductor monomial has exactly $t$ surviving blocks; monomials on more
blocks are zero. This is $J_t/J_{t+1}$. The same coordinates identify
$Y_{g+1}$ with $J_t$, proving the sheaf equality by faithful flatness
of completion. The case $t=0$ means that the conductor is the unit
ideal and the next stratum is absent locally.

The ring is the ring of tuples on its coordinate branches that agree
on every pairwise intersection. To verify this directly, compare the
coefficient of each monomial on all branches where it survives; the
pairwise agreements give a unique common coefficient. If an element
$b$ of the total fraction ring has $b^2,b^3$ in the image, it first
belongs to the normalization, by integrality. Its two restrictions
on any intersection have equal squares and cubes, hence are equal
in that reduced ring. Thus $b$ lies in the image. This proves
seminormality by the square-and-cube criterion.

Formula~\eqref{eq:block-hilbert-v153} follows by sorting nonzero
monomials by their block support. The top-dimensional contributions
come from the $\binom{s}{t}$ supports of size $t$, giving dimension
$mt$ and multiplicity $\binom{s}{t}$. We supply the homological step
for the depth, since branch counting alone does not give it. The
simplicial complex of the squarefree monomial ideal $J_{t+1}$ consists
of sets using at most $t$ blocks. Any induced subcomplex meeting $b$
blocks is homotopy equivalent to the $(t-1)$-skeleton of the simplex
on those $b$ blocks: collapse each block to a representative vertex.
The two composites are contiguous to the identity, because adjoining
the representatives never increases the number of blocks. If $b\leq t$
the complex is a simplex; otherwise its only reduced homology is in
degree $t-1$, with rank $\binom{b-1}{t}$. This last assertion follows
also by the elementary induction that adds the last vertex to a
simplex skeleton.

The squarefree multidegree-$W$ part of the Koszul complex on all $sm$
variables is the augmented simplicial cochain complex of this induced
subcomplex, with degree shifted by $|W|-1$. Non-squarefree multidegrees
are contractible in the relevant summands, by pairing faces using a
variable whose exponent exceeds one. Consequently a nonzero positive
Betti number has homological degree $|W|-t$. The full vertex set gives
the nonzero degree $sm-t$, since $s>t$ when $g>0$. The projective
dimension is therefore $sm-t$, and the Auslander--Buchsbaum formula
gives depth $t$. When $t=0$ the quotient is simply $\C$. Adjoining
$q$ power-series parameters gives depth $q+t$ in every case, and
comparison with $q+mt$ proves the Cohen--Macaulay criterion.
All the constructions commute with relabelling roots, proving descent.
\end{proof}

\begin{corollary}[Binary divisor-image intersections]
\label{cor:binary-intersections-v153}
For admissible binary divisor degrees $a_1,\ldots,a_v$, their
scheme-theoretic intersection is $Y_{\max a_i}$. On the squarefree
locus all multiple intersections of normalization branches, rather
than only the degree images, are given by the union-of-subsets rule
in Theorem~\ref{thm:squarefree-v153}.
\end{corollary}
\begin{proof}
The reduced binary images are nested as closed subschemes: a binary
form of degree $b$ has a divisor of each degree $a\leq b$, and
reduced closed subschemes are determined by their closed points.
Their ideals are thus nested, so their sum is the largest ideal.
The assertion about branches has already been proved as an identity
of coordinate ideals, without taking radicals.
\end{proof}


\section{An incidence atlas for all binary root collisions}
\label{sec:collision-atlas-v153}
The squarefree theorem has a uniform replacement when any number of
root clusters collide. This replacement specifies the full completed
image, not merely a transverse family in it.

\begin{theorem}[Simultaneous collision atlas]\label{thm:atlas-v153}
Let $K\in Y_g$ be binary with exact common divisor
$\prod_{i=1}^b\ell_i^{s_i}$ and $s=\sum_i s_i$. There are formal
ambient coordinates consisting of the coefficients of monic
polynomials $P_i$ of degree $s_i$, polynomials $R_{ij}$ of degree
less than $s_i$ ($1\leq j\leq m$), and $N_0-rs$ smooth parameters
$z$. They are centred at $P_i=x_i^{s_i}$ and $R_{ij}=0$.
For every $\alpha=(a_i)$ with $0\leq a_i\leq s_i$ and $\sum a_i=g$,
let $B_\alpha$ be the regular power-series ring in the coefficients
of monic $A_i,B_i'$ of degrees $a_i,s_i-a_i$, polynomials $U_{ij}$
of degree less than $s_i-a_i$, and the same $z$. The completed
normalization and image are exactly
\begin{equation}\label{eq:collision-image-v153}
 T\longrightarrow\prod_\alpha B_\alpha,\qquad
 P_i\mapsto A_iB_i',\quad R_{ij}\mapsto A_iU_{ij},\qquad
 \widehat{\OO}_{Y_g,K}=T/\bigcap_\alpha\ker(T\to B_\alpha).
\end{equation}
Here $T=\C[[\operatorname{coeff}(P_i),\operatorname{coeff}(R_{ij}),z]]$.
Degree-zero monic factors are one, and negative-degree remainder
spaces are zero. The formula is independent, up to isomorphism of
completed incidence diagrams, of the frame, the selected $F_0$, and
the local root coordinates. It covers every higher and simultaneous
binary collision.
\end{theorem}
\begin{proof}
Choose $F_0\in K$ having exact order $s_i$ at every common root.
The required nonvanishing of its residual value excludes finitely
many hyperplanes in $K$, so a choice exists. Complete it to a frame.
Weierstrass preparation at the distinct roots produces $P_i$; division
of the remaining frame members produces $R_{ij}$. The map
\[
 S_D/K\longrightarrow\bigoplus_i\C[x_i]/(x_i^{s_i})
\]
is surjective: its kernel contains $K$, and the usual interpolation
map from $S_D$ is surjective. Thus the Grassmannian tangent space
permits independent variation of all $rs$ jets. Division and preparation
change these jets by an invertible triangular matrix, whose diagonal
contains the nonzero residual values of $F_0$. Their differentials
are independent. Formal inverse coordinates supply the remaining
$N_0-rs$ parameters; this integer is nonnegative since
$r\leq D-s+1$.

A marked divisor near a chosen normalization point splits uniquely
among the coprime root clusters. In cluster $i$ its monic representative
has degree $a_i$ and divides both $P_i$ and all $R_{ij}$. Monic division
is equivalent, over every local Artin base, to the unique equations
$P_i=A_iB_i'$ and $R_{ij}=A_iU_{ij}$. Conversely these equations give
the marked divisor and the residual subbundle. Thus the normalized
incidence functor is represented by the displayed regular ring.
It is finite over $T$: monic factor coefficients are integral over the
coefficients of $P_i$, while the remaining factors and quotients are
uniquely determined by division. Lemma~\ref{lem:formal-image-v153}
now gives the exact completed image and the intersection of kernels.

Changing a frame or a coordinate changes only a presentation of the
same functor of a subbundle with a marked divisor. The mutually inverse
changes of incidence data therefore induce isomorphisms of its complete
representing rings and of the resulting finite-image diagrams. This
also proves that the calculation describes the entire completed image.
A point of that image over an Artin ring need not lift to its
normalization; no such lifting assertion is used or implied.
\end{proof}


\section{The singular Fitting ideal of every double-root model}
\label{sec:fitting-v153}
The atlas specializes for $s=2$, $g=1$ to the following ring, after
completion of the square and suppression of smooth parameters:
\[
 A_m=R\oplus It\subset B=\C[[u_1,\ldots,u_m,t]],\qquad
 R=\C[[u_1,\ldots,u_m,\Delta]],\ I=(u_1,\ldots,u_m),\ \Delta=t^2.
\]
Put $w_i=u_it$. The defining equations are
$u_iw_j-u_jw_i$ and $w_iw_j-\Delta u_iu_j$.

\begin{theorem}[Closed singular-ideal formula]\label{thm:fitting-v153}
For every $m\geq1$, the intrinsic singular Fitting ideal is
\begin{equation}\label{eq:fitting-v153}
 \Fitt_{m+1}\Omega^{\mathrm{cont}}_{A_m/\C}
  =(I^{m+1}+\Delta I^m)\oplus I^m t
  = I^{m-1}(I^2,(w_1,\ldots,w_m),\Delta I).
\end{equation}
The first equality is a decomposition as an $R$-submodule of $B$;
the right side is an ideal of $A_m$. It is not merely an ideal of
Jacobian minors left unevaluated. For $m=1$ it is
$(u^2,w,\Delta u)$, and its radical is the conductor $(u,w)$ for
all $m$.
\end{theorem}
\begin{proof}
There are $2m+1$ ambient variables and dimension $m+1$, so the Fitting
ideal consists of the $m$-minors of the Jacobian of the displayed
relations. Evaluate this matrix in $B$ using $w_i=u_it$ and
$\Delta=t^2$. A minor not using the $\Delta$ column has $u$-degree
$m$. It vanishes at $t=0$: the only remaining rows are the Koszul
rows in the $w$ columns, of rank at most $m-1$. Hence its terms have
a positive power of $t$, and belong to $\Delta I^m\oplus I^m t$.
A minor using the $\Delta$ column has $u$-degree $m+1$, and hence
lies in $I^{m+1}\oplus I^{m+1}t$. This proves the first inclusion.

Use the $m$ rows
\[
 u_1w_j-u_jw_1\ (2\leq j\leq m),\qquad w_1^2-\Delta u_1^2.
\]
The columns $(w_1,w_2,\ldots,w_m)$, $(u_1,w_2,\ldots,w_m)$,
and $(\Delta,w_2,\ldots,w_m)$ give, up to sign and nonzero scalar,
$u_1^m t$, $\Delta u_1^m$, and $u_1^{m+1}$. The ideal is invariant
under simultaneous linear change of $u$ and $w$. In characteristic
zero the powers of linear forms span each symmetric power, so these
three minors generate the three indicated spaces of all $u$-monomials.
They generate the claimed $R$-module; it is an $A_m$-ideal since
multiplication by $w_i=u_it$ preserves it. This proves equality.
The formula also immediately gives its radical, without discarding
the nonreduced Fitting structure in the assertion itself.
\end{proof}


\section{A structural dichotomy for pure-power divisor fibres}
\label{sec:reciprocal-v153}
Let $E=\C a\oplus W$, $d=\dim W\geq1$, and $g,k\geq1$.
The factorization fibre of $a^{g+k}$ records monic factors of degrees
$g$ and $k$. Interchanging these factors identifies its two descriptions;
we may and shall take $g\geq k$. Write
\[
 Q_i\in\Sym^iW\quad(1\leq i\leq k),\qquad
 C_0=1,\quad C_j=-\sum_{i=1}^kQ_i C_{j-i},\quad C_j=0\ (j<0).
\]
Let $P=\Sym(\bigoplus_{i=1}^k(\Sym^iW)^*)$, with its generators
of weight $i$, and let $\mathfrak n$ be its variable ideal. Put
\begin{equation}\label{eq:reciprocal-algebra-v153}
 F_{g,k}(W)=P/\mathcal I,\qquad
 \mathcal I=(\operatorname{coeff}C_{g+1},\ldots,
                                    \operatorname{coeff}C_{g+k}).
\end{equation}
It is an Artin local algebra. Indeed its complex points are exactly
the factorizations of $a^{g+k}$, hence the single origin by unique
factorization. The polynomial and completed presentations coincide.
If $K=a^{g+k}L$ with $\gcd(L)=1$, cancellation by a member of $L$
not divisible by $a$, as in Lemma~\ref{lem:coprime-lifts-v153},
identifies this with the full degree-$g$ normalization fibre at $K$.

\begin{theorem}[Minimal equations and complete-intersection classification]
\label{thm:reciprocal-structure-v153}
The embedding dimension and minimal number of equations of
$F_{g,k}(W)$ in a minimal regular local presentation are
\begin{align}
 \operatorname{edim}F_{g,k}(W)&=\binom{d+k}{k}-1,\label{eq:edim-v153}\\
 \mu(\mathcal I)&=\sum_{j=g+1}^{g+k}\binom{d+j-1}{j}
                  =\binom{d+g+k}{d}-\binom{d+g}{d}.
                  \label{eq:minrels-v153}
\end{align}
More precisely the minimal relation module has the equivariant
weighted decomposition
\begin{equation}\label{eq:relation-module-v153}
 \mathcal I/\mathfrak n\mathcal I
                 \simeq\bigoplus_{j=g+1}^{g+k}(\Sym^jW)^*.
\end{equation}
Consequently this fibre is a complete intersection if and only if
$d=1$. For $d=1$, it has length $\binom{g+k}{k}$, weighted Hilbert
series $\prod_{i=1}^k(1-z^{g+i})/(1-z^i)$, and one-dimensional
socle generated by $Q_k^g$, of weight $gk$.
For $k=1$ and arbitrary $d$,
\begin{equation}\label{eq:k-one-v153}
 F_{g,1}(W)=\Sym(W^*)/\mathfrak n^{g+1},\quad
 \length F_{g,1}(W)=\binom{d+g}{g},\quad
 \operatorname{Soc}F_{g,1}(W)=\Sym^g(W^*).
\end{equation}
In particular the last algebra is Gorenstein precisely when $d=1$.
No Gorenstein classification for $d>1,k>1$ is needed for the theorem.
\end{theorem}
\begin{proof}
Since $g\geq k$, each defining form has no linear term in the
coefficient variables. Thus the presentation is minimal and
\eqref{eq:edim-v153} follows by summing $\dim\Sym^iW$.
For $d=1$, the regular-sequence argument in the proof of
Theorem~\ref{thm:binary-fibres-v153} applies to the $k$ reciprocal
relations. In particular none of them belongs to $\mathfrak n\mathcal I$.
It gives the stated Hilbert series, length, and the classical rectangular
socle class.

For general $W$, taking coefficients of $C_j$ gives an equivariant
map $(\Sym^jW)^*\to\mathcal I/\mathfrak n\mathcal I$.
Its source is irreducible. This map is nonzero: choose one coordinate
$x$ of $W$, keep only the coefficients of $x^i$ in every $Q_i$, and
set every other coefficient to zero. The coefficient of $x^j$ maps
to the $j$th binary reciprocal relation. Were it in
$\mathfrak n\mathcal I$ before specialization, its specialization
would be in the binary $\mathfrak n\mathcal I$, contradicting the
preceding minimality. Irreducibility now makes the entire coefficient
map injective. The different $j$ have different weights, so their
images are independent. They generate the relation module by the
chosen presentation, proving \eqref{eq:relation-module-v153} and
\eqref{eq:minrels-v153}.

If $d\geq2$, the sequence $\binom{d+j-1}{j}$ strictly increases in
$j$. Comparing its terms at $j=g+i$ and $j=i$ shows that the minimal
number of relations is strictly greater than the embedding dimension.
A zero-dimensional complete intersection in a minimal regular local
presentation has exactly that many minimal relations as variables;
hence these fibres are not complete intersections. This proves both
directions of the classification. Finally when $k=1$ the sole form
is $C_{g+1}=(-Q_1)^{g+1}$. Its coefficients are, with nonzero multinomial
factors, every monomial of degree $g+1$ in the $d$ variables. This gives
\eqref{eq:k-one-v153}, including the socle and length.
\end{proof}

\begin{corollary}[The ternary-pencil ramification fibre]
\label{cor:ternary-fibre-v153}
At the common limiting relation algebra with $e=9$, $g=6$ and $k=2$,
the normalization fibre has embedding dimension $44$ and exactly
\[
 \binom{14}{7}+\binom{15}{8}=9867
\]
minimal scalar equations. Their two weighted pieces have dimensions
$3432$ and $6435$, in weights seven and eight. In particular this
actual pencil-orbit boundary fibre is not a complete intersection.
\end{corollary}
\begin{proof}
Here $d=e-1=8$. Apply \eqref{eq:edim-v153} and
\eqref{eq:relation-module-v153}. These are minimal equation counts,
not the number of coefficient variables or the length of the fibre.
\end{proof}

The binary factorization rings in this section are classical. The
additional point of the multivariate argument is that polarization
preserves every full irreducible coefficient space in the minimal
relation module. Neither the Grassmannian identification nor a finite
calculation in a few variables is being used as a claim of priority
or as a substitute for that argument.


\section{Homological invariants of the divisor fibres}
\label{sec:homological-v154}
We first extract all graded Betti numbers from the squarefree local
model and then identify a determinant-apolar quotient of the
multivariate fibres. A separate exact calculation treats an entire
multivariate fibre, rather than only a quotient or its minimal
relations. The grading conventions will be specified in each case.

\begin{theorem}[Betti numbers of the squarefree block model]
\label{thm:block-betti-v154}
Let $S=\C[u_{ij}:1\leq i\leq s,1\leq j\leq m]$ with its standard
multigrading, let $0\leq t<s$, and put $A=S/J_{t+1}$, with the
block-support ideal of Theorem~\ref{thm:squarefree-v153}.
For a nonempty squarefree multidegree $F$ let $b(F)$ be the number
of blocks it meets. Besides $\beta_{0,0}=1$, the only nonzero
multigraded Betti numbers are
\begin{equation}\label{eq:block-multibetti-v154}
 \beta_{i,F}^S(A)=\binom{b(F)-1}{t}
       \quad\text{if }b(F)>t\text{ and }i=|F|-t.
\end{equation}
For $t=0$ this is the Koszul resolution of $\C$.
For all $t$, the singly graded numbers are
\begin{equation}\label{eq:block-betti-v154}
 \beta_{i,i+t}^S(A)=
 \sum_{b=t+1}^{\min(s,i+t)}\binom{s}{b}\binom{b-1}{t}
       [z^{i+t}]\big((1+z)^m-1\big)^b\quad(i\geq1).
\end{equation}
The ideal $J_{t+1}$ has a $(t+1)$-linear resolution. The projective
dimension of $A$ is $sm-t$. In the Cohen--Macaulay case $m=1$,
its type is $\binom{s-1}{t}$; it is Gorenstein exactly when
$t=0$ or $t=s-1$. The same Betti numbers over the completed ambient
regular ring apply to the local model after adjoining its smooth
parameters.
\end{theorem}
\begin{proof}
The Stanley--Reisner complex consists of sets meeting at most $t$
blocks. If $t\geq1$, its induced subcomplex on $F$ has, in each
nonempty block, a simplex all of whose vertices have identical
external links. Contract that simplex to one of its vertices: the
map and the inclusion are contiguous, because adjoining the chosen
vertex to any face does not change its block support. Applying this
one block at a time gives a homotopy equivalence with the
$(t-1)$-skeleton of the simplex on $b(F)$ vertices. If $b(F)\leq t$
this simplex is contractible. Otherwise its reduced homology is
concentrated in degree $t-1$, with dimension $\binom{b(F)-1}{t}$.
The latter assertion follows by attaching its top faces successively,
or from the augmented simplicial chain complex and the binomial
identity for its Euler characteristic. Hochster's multigraded
formula now gives \eqref{eq:block-multibetti-v154}; see
\cite[Chapter~5]{MillerSturmfels2005}. If $t=0$, every variable is
in the ideal and the Koszul formula gives the same result directly.

For a support of size $b$, the coefficient in
\eqref{eq:block-betti-v154} counts choices of $i+t$ vertices meeting
each selected block. Summing proves the formula. All ideal syzygies
therefore have the indicated linear degrees. Taking $F$ to be all
$sm$ vertices gives the nonzero last Betti number
$\binom{s-1}{t}$ in homological degree $sm-t$; no larger degree
occurs. When $m=1$, Auslander--Buchsbaum gives depth $t$, equal to
dimension, and the last Betti number is the Cohen--Macaulay type.
Completion and adjoining independent regular parameters are flat
operations preserving the minimal free resolution. These are
classical squarefree-monomial methods. The assertion here is their
full specialization to the normalization image model, not a claim
of priority for Hochster's formula or skeleton resolutions.
\end{proof}

For the remaining results use the weighted reciprocal algebra
$F_{g,k}(W)$ of \eqref{eq:reciprocal-algebra-v153}, where
$d=\dim W$ and $g\geq k\geq2$. Set all $Q_i$ except $Q_2$ equal
to zero and, in addition, impose every coefficient of $Q_2^2$.
Denote the resulting algebra by
\begin{equation}\label{eq:apolar-quotient-v154}
 B_d=\Sym((\Sym^2W)^*)/(\operatorname{coeff}Q_2^2).
\end{equation}
Its coefficient variables have reciprocal weight two. When referring
to its ordinary polynomial grading, we instead give them degree one.

\begin{theorem}[A determinant-apolar quotient of every higher fibre]
\label{thm:apolar-quotient-v154}
There is a $\operatorname{GL}(W)$-equivariant graded surjection
$F_{g,k}(W)\twoheadrightarrow B_d$ for every $g\geq k\geq2$.
After fixing the standard coefficient/differential pairing, $B_d$
is the apolar algebra of the determinant of a generic symmetric
$d\times d$ matrix. Its polynomial-degree Hilbert function is
\begin{equation}\label{eq:Narayana-v154}
 \dim(B_d)_p=\frac{1}{d+1}\binom{d+1}{p}\binom{d+1}{p+1}
                 \quad(0\leq p\leq d).
\end{equation}
It is Gorenstein, its reciprocal-weight socle is $2d$, and
\begin{equation}\label{eq:Catalan-v154}
 \length F_{g,k}(W)\ \geq\ \length B_d
       =\frac{1}{d+2}\binom{2d+2}{d+1}.
\end{equation}
For $k=2$ and $g=2$ or $3$, the quotient by the coefficients of
$Q_1$ alone is already $B_d$.
In particular the actual ternary-pencil fibre $F_{6,2}(\C^8)$ has
length at least $4862$, independently of its embedding dimension
$44$ and its $9867$ minimal equations.
\end{theorem}
\begin{proof}
Under the stated substitution, the reciprocal series becomes
$(1+Q_2z^2)^{-1}$. Its odd coefficients vanish and its even
coefficients are powers of $Q_2$. Since $g+1\geq3$, imposing
$Q_2^2=0$ annihilates all defining coefficients. This proves the
surjection, functorially for linear changes of $W$. When $k=2$ and
$g=2$ or $3$, the first even defining coefficient is $Q_2^2$ and
there are no further independent equations after $Q_1=0$.

Choose coordinates $x_1,\ldots,x_d$ on $W$ and write
$Q_2=\sum_i y_{ii}x_i^2+\sum_{i<j}y_{ij}x_ix_j$.
Let $X=(X_{ij})$ be symmetric, with independent entries for $i\leq j$,
and interpret $y_{ij}$ as $\partial/\partial X_{ij}$. Then
$Q_2(x)$ acts as directional differentiation in the rank-one
symmetric direction $xx^t$. The determinant of $X+zxx^t$ is linear
in $z$, so every coefficient of $Q_2(x)^2$ annihilates $\det X$.
These coefficients are, up to nonzero scalars, $y_{ii}^2$,
$y_{ii}y_{ij}$, $y_{ij}^2+2y_{ii}y_{jj}$,
$y_{ii}y_{jl}+y_{ij}y_{il}$ for distinct $i,j,l$, and
$y_{ij}y_{kl}+y_{ik}y_{jl}+y_{il}y_{jk}$ for four distinct indices.
They are exactly the quadratic generators in
\cite[Proposition~3.1 and Theorem~3.10]{ShafieiSymmetric}.
That theorem proves equality with the full apolar ideal, not merely
containment. Its Corollary~2.7 gives \eqref{eq:Narayana-v154} and
the Catalan length. Alternatively, that Hilbert function counts the
independent minors obtained by differentiating $\det X$.
Apolarity to one homogeneous form gives a perfect pairing between
complementary degrees and a one-dimensional socle in degree $d$.
Replacing polynomial degree by twice that degree proves the weight
statement. The length inequality follows from the surjection.
For $d=8$ the Catalan number is $4862$.

The apolar Hilbert function and its Gorenstein property are classical
results of the cited determinant theory. The new identification in
this setting is the functorial quotient of the entire family of
normalization fibres. A quotient's Gorenstein property is not being
attributed to the original fibre.
\end{proof}

\begin{proposition}[The complete first multivariate two-factor fibre]
\label{prop:F22-v154}
Let $A=F_{2,2}(\C^2)$ and write
$Q_1=ax+by$, $Q_2=cx^2+dxy+ey^2$.
Then
\begin{gather}\label{eq:F22-hilbert-v154}
 H_A^{\mathrm{wt}}(z)=1+2z+6z^2+6z^3+7z^4,\qquad
 \length A=22,\qquad\dim\operatorname{Soc}A=7,\\
 \label{eq:F22-local-hilbert-v154}
 H_{\gr_{\mathfrak m}A}(z)=1+5z+6z^2+5z^3+5z^4.
\end{gather}
The weight grading gives $a,b$ weight one and $c,d,e$ weight two.
In the minimal weighted free resolution over
$S=\C[a,b,c,d,e]$, the shift polynomials
$B_i(z)=\sum_j\beta_{i,j}^S(A)z^j$ are
\begin{align*}
 B_0&=1,& B_1&=4z^3+5z^4,\\
 B_2&=4z^5+12z^6+12z^7,&
 B_3&=21z^8+20z^9+z^{10},\\
 B_4&=21z^{10}+8z^{11},&B_5&=7z^{12}.
\end{align*}
In particular the total Betti numbers are $(1,9,28,42,29,7)$.
This fibre is not Gorenstein. The assertion concerns this whole
five-variable fibre, not the determinant-apolar quotient.
\end{proposition}
\begin{proof}
The nine initial equations are the coefficients of
$-Q_1^3+2Q_1Q_2$ and $Q_1^4-3Q_1^2Q_2+Q_2^2$.
For graded reverse lexicographic order $e>d>c>b>a$, a Groebner basis
is the following list (a scalar multiple does not change a member):
\begin{gather*}
 (c,d,e)^3,\quad b(c,d,e)^2,\quad a(c,d,e)^2,\\
 b^2e-e^2,\quad b^2d-de,\quad b^2c-a^2e,\quad b^3-2be,\\
 2abe-de,\quad 2a^2e+2abd-2ce-d^2,\quad 2abc-cd,\\
 3ab^2-2ae-2bd,\quad a^2d-cd,\quad a^2c-c^2,\\
 3a^2b-2ad-2bc,\quad a^3-2ac.
\end{gather*}
Here a monomial ideal in the list means all its minimal monomial
generators. There are $10+6+6+12=34$ polynomials. Polynomial
division gives both equality with the ideal of the nine equations
and zero remainders for all their basis S-pairs. This is a finite
exact arithmetic verification; the explicit input and all reduction
and rank tests are supplied in the accompanying calculation.
The standard monomials, grouped by weight, are
\begin{equation}\label{eq:F22-basis-v154}
\begin{array}{c|l}
0&1\\
1&a,b\\
2&a^2,ab,b^2,c,d,e\\
3&ac,bc,ad,bd,ae,be\\
4&c^2,cd,d^2,a^2e,ce,de,e^2.
\end{array}
\end{equation}
Reduction of products of these monomials gives
$\dim A/\mathfrak m=1$ and
$\dim(\mathfrak m^i)=(21,16,10,5,0)$ for $1\leq i\leq5$.
The common kernel of multiplication by $a,b,c,d,e$ is exactly the
last row of \eqref{eq:F22-basis-v154}. This proves both distinct
Hilbert series and the type.

For precision, the resolution calculation uses no numerical rank
thresholds. In weight $w$ and homological degree $i$, take basis
pairs $(J,m)$ with $J\subset\{a,b,c,d,e\}$ of size $i$, $m$ a
standard monomial, and $\operatorname{wt}(J)+\operatorname{wt}(m)=w$.
The differential is
\begin{equation}\label{eq:F22-koszul-v154}
 d(J,m)=\sum_{v\in J}(-1)^{\operatorname{pos}(v,J)-1}
                  (J\setminus\{v\},\operatorname{NF}(vm)).
\end{equation}
All entries are rational, explicitly determined by the displayed
34 polynomials. Exact row reduction and
$\beta_{i,w}=\dim C_{i,w}-\rank d_{i,w}-\rank d_{i+1,w}$ give the
six shift polynomials above. The calculation is fully specified by
\eqref{eq:F22-basis-v154}--\eqref{eq:F22-koszul-v154}; the archived
rank table makes each finite elimination independently repeatable.
As additional checks, their alternating sum is
$H_A^{\mathrm{wt}}(z)(1-z)^2(1-z^2)^3$ and the last Betti space is
the socle shifted by the total variable weight eight. These checks
do not replace the exact rank computations.
\end{proof}


\section{Spectral Fitting schemes as power-zero loci in one reciprocal section}
\label{sec:power-spectral-v155}
The algebra just calculated supplies a single geometric object in
which all spectral Fitting schemes can be read simultaneously.
This avoids constructing a separate binary form after each Fitting
ideal has already been recovered.

Let $V$ have dimension $n$ and $R\subset\Sym^2V$ be any pencil.
Fix $h\geq1$ and put $\widetilde E=\C a\oplus V^*$.
Consider the divisor-incidence point
\begin{equation}\label{eq:canonical-reciprocal-point-v155}
 K_h=a^{2h+2}\widetilde E
         \subset\Sym^{2h+3}\widetilde E.
\end{equation}
The entire degree-$2h$ normalization fibre over $K_h$ is
$F_{2h,2}(V^*)$, because $\widetilde E$ is gcd-free.  Its section
by the linear-factor coefficients is $B_{n,h}$ and
$(B_{n,h})_1=\Sym^2V$.  Thus the given pencil defines a line
$\Lambda=\Pj(R)$ in the projectivized degree-one space of this
actual reciprocal section.  This construction uses a new formal
line $\C a$ and is functorial for isomorphisms of $V$.

For $0\leq p\leq hn$ let $\mathcal J_p$ be the ideal sheaf on
$\Lambda$ locally generated by the coordinates of
\begin{equation}\label{eq:power-map-v155}
       \ell^p\in(B_{n,h})_p\otimes\OO_\Lambda(p),
\end{equation}
where $\ell$ is the universal pencil member.  Multiplying a local
representative of $\ell$ by a unit does not change this ideal.
We set $\mathcal J_0=\OO_\Lambda$.
Let $\mathcal T_R$ be the cokernel of the universal symmetric
matrix $V^*\otimes\OO_\Lambda(-1)\to V\otimes\OO_\Lambda$.
This definition includes singular pencils.

\begin{lemma}[The local power-ideal formula]
\label{lem:power-valuation-v155}
Let $A(t)$ be a symmetric matrix over $\C[[t]]$ of rank $r$ over
$\C((t))$.  Write its finite elementary exponents as
$0\leq e_1\leq\cdots\leq e_r$ and put $e_i=+\infty$ for $i>r$.
For $1\leq p\leq hn$, write $p=hq+s$, $0\leq s<h$.
The ideal generated by the coordinates of
$\ell_{A(t)}^p$ in $(B_{n,h})_p\otimes\C[[t]]$ is
\begin{equation}\label{eq:valuation-power-v155}
 \left(t^{\,h\sum_{i=1}^q e_i+s e_{q+1}}\right).
\end{equation}
An empty sum is zero, the last term is omitted if $s=0$, and
$t^{+\infty}$ means the zero ideal.
\end{lemma}
\begin{proof}
A symmetric matrix over this discrete valuation ring is congruent
to
\[\diag(t^{e_1},\ldots,t^{e_r},0,\ldots,0).\]
Here is the required elementary reduction.  Divide by the least
valuation of an entry.  A diagonal unit is a pivot; if only an
off-diagonal entry is a unit, the value of the form on a sum of two
basis vectors supplies a unit because $2$ is invertible.  Split
off the resulting rank-one summand by elimination and continue.
Every remaining unit has a square root in $\C[[t]]$, so the units
can be removed from the diagonal.  The finite exponents are the
Smith exponents, as congruence is in particular equivalence.

Congruence acts on every graded piece of $B_{n,h}$ by an invertible
matrix over $\C[[t]]$.  It therefore preserves the ideal in question.
The apolar embedding
\[
 (B_{n,h})_p\longrightarrow
       \C[X_{ij}]_{hn-p},\qquad
       D\longmapsto D(\det X)^h
\]
is injective and splits as a map of complex vector spaces.  After
tensoring with $\C[[t]]$, the coordinates of a vector and of its
image generate the same ideal.  We may therefore calculate the
coefficient ideal of $[z^p]\det(X+zA(t))^h$.

For diagonal $A(t)$, expand each of the $h$ determinant factors
in its chosen diagonal entries.  Every term is divisible by
$t^{\sum c_i e_i}$ with $0\leq c_i\leq h$ and $\sum c_i=p$.
The least such exponent is obtained by filling the smallest $e_i$
to capacity $h$, giving \eqref{eq:valuation-power-v155}.
There is no cancellation at this least exponent: specialize $X$
to $\diag(x_1,\ldots,x_n)$.  Then
\[
 \det(X+zA(t))^h=\prod_i(x_i+zt^{e_i})^h
\]
(with $x_i^h$ for a zero diagonal entry of $A$).
The monomial $\prod x_i^{h-c_i}$ has coefficient
$\prod\binom h{c_i}\,t^{\sum c_i e_i}$, which is nonzero and has
exactly the required valuation.  If $p>hr$, every term is zero.
This proves all the conventions, including the zero-ideal case.
\end{proof}

\begin{theorem}[A reciprocal multiplication model for every spectral Fitting ideal]
\label{thm:power-Fitting-v155}
For every complex quadratic pencil, regular or singular, and for
$0\leq j\leq n$, there is an equality of coherent ideal sheaves
on its intrinsic line
\begin{equation}\label{eq:power-Fitting-v155}
          \mathcal J_{h(n-j)}=\Fitt_j(\mathcal T_R)^h.
\end{equation}
In particular, when $h=1$, the complete spectral Fitting schemes,
including their multiplicities and any whole-line components, are
exactly the power-zero schemes inside the one quadratic section
of the normalization fibre at $K_1$.

For a regular pencil and a spectral point $x$, set
$b_{x,j}=\length\OO_{\Lambda,x}/\Fitt_j(\mathcal T_R)_x$.
The local length of the power-zero scheme is $h b_{x,j}$, and,
for $0\leq j<n$,
\begin{equation}\label{eq:power-Segre-v155}
 \lambda_{x,j+1}
   =\frac1h\left(\length(\OO_{\Lambda,x}/\mathcal J_{h(n-j),x})
           -\length(\OO_{\Lambda,x}/\mathcal J_{h(n-j-1),x})\right).
\end{equation}
Thus its multiplication, together with the embedded pencil line,
recovers the full Segre data and their common projective support.
For $n\geq3$, the unmarked inverse makes this an invariant of the
finite failure algebra.  For a singular pencil, the construction recovers its
normal rank and all its spectral Fitting ideals; it does not assert
that these ideals alone recover the singular minimal indices.
\end{theorem}
\begin{proof}
At a point of the smooth curve $\Lambda$, pass to the completed
discrete valuation ring.  The ideal of $(n-j)$-minors of a matrix
with the stated Smith exponents is
$(t^{e_1+\cdots+e_{n-j}})$, with the zero-ideal convention.
Apply Lemma~\ref{lem:power-valuation-v155} with $p=h(n-j)$.
The equality of completed ideals descends by faithful flatness.
This proves \eqref{eq:power-Fitting-v155} at every stalk and hence
globally.  This proof does not replace an ideal by its radical.

For a regular pencil the positive Smith exponents at $x$, in
reverse order, are $\lambda_{x,1}\geq\lambda_{x,2}\geq\cdots$.
Consequently $b_{x,j}=\sum_{i>j}\lambda_{x,i}$.
Equation~\eqref{eq:power-Fitting-v155} multiplies these lengths by
$h$, so taking successive differences proves
\eqref{eq:power-Segre-v155}.  The powers that vanish at the generic
point similarly determine the normal rank in the singular case.
Frame changes act on $\Lambda$ by projectivities and congruences
act equivariantly on the reciprocal construction.  Theorem~
\ref{thm:artin-local-inverse-v146} recovers precisely the needed
pencil up to these changes, proving the intrinsic assertion.
\end{proof}

\begin{remark}[Families and the scope of the ideal identity]
\label{rem:power-basechange-v155}
For a vector bundle and pencil subbundle over any complex base,
the quadratic-section algebra is obtained by applying the fixed
$\operatorname{GL}(V)$-equivariant construction to local frames.
Its graded pieces are vector bundles, and the coefficient ideal of
each power map commutes with arbitrary base change: it is the image
of the dual of a map between vector bundles.  No flatness of its
zero scheme is needed for that assertion.  Formula
\eqref{eq:power-Fitting-v155} has been proved on each complex
pencil line (and on a regular one-dimensional base by the same
local argument).  We do not infer an equality of those relative
ideals on an arbitrary nonreduced parameter space from equality
on its geometric fibres.  Nor are the power-zero schemes asserted
to be flat across a change in their length.
\end{remark}

\begin{corollary}[The native specialization order in reciprocal coordinates]
\label{cor:power-order-v155}
On the fixed-discriminant, fixed-corank locus
$X_{\boldsymbol m,\boldsymbol c}(A)$ of
\eqref{eq:fixed-spectral-locus}, a specialization from the labelled
partition tuple $\boldsymbol\lambda$ to $\boldsymbol\mu$ exists
if and only if, at every labelled root and every $j$, the local
length of the corresponding power-zero scheme does not decrease.
Equivalently the special power-zero divisor contains the generic
one, using the common labelled spectral line.
The reduced closure of a full $O(A)$-orbit within this locus is
therefore specified by these length inequalities.  Each inequality
can be written as the vanishing of the requisite jets of the
coordinates of the power map \eqref{eq:power-map-v155}.
Every strict transition is realized by a flat family of the original
finite failure neighbourhoods and is first distinguished at order
$n^2+2n-4$.
\end{corollary}
\begin{proof}
For fixed multiplicity $m_x$, dominance
$\mu_x\preceq\lambda_x$ is equivalent to
$\sum_{i>j}\mu_{x,i}\geq\sum_{i>j}\lambda_{x,i}$ for every $j$.
Use \eqref{eq:power-Segre-v155} and
Theorem~\ref{thm:fixed-spectral-classification}.
In a local coordinate $t$ at a labelled root, a coefficient ideal
has colength at least $b$ exactly when each of its generators is
divisible by $t^b$; this is the asserted finite jet condition.
The closure order and the flat realization are the inherited
full-orthogonal results, ultimately using Ohta's theorem.  The
new assertion is their expression inside one reciprocal
multiplication algebra.  No reducedness assertion is made for the
scheme cut out by the jet equations before taking its reduction.
\end{proof}

\begin{example}[A rank profile with a genuinely different power-zero scheme]
\label{ex:power-profile-v155}
At a root of multiplicity four the partitions $(3,1)$ and $(2,2)$
have equal determinant and corank.  In $B_{4,h}$ their sixth powers
when $h=2$ (in general their $3h$th powers) cut out local schemes
of lengths $h$ and $2h$, respectively.  Their $4h$th powers have
length $4h$ in both cases.  Thus the nonreduced power-zero scheme,
not just the set of non-Lefschetz elements, distinguishes them.
The specialization $(3,1)\rightsquigarrow(2,2)$ increases the first
of these lengths.  Additional simple spectral blocks give the
same distinction in all larger dimensions, with power indices
$h(n-1)$ and $hn$.
\end{example}

The representation theory and strong Lefschetz property used above
are classical inputs, explicitly credited to Nagaoka--Wachi and,
for $h=1$, Shafiei.  The geometric statements here are the exact
section of the whole reciprocal family, its size and Loewy-depth
consequences for actual first-relation boundary fibres, and the
power-ideal identity on pencil lines.  The last construction is an
intrinsically induced reciprocal model; it is not a claim that the
common extreme $\operatorname{GL}(E)$ boundary point remembers the
starting pencil.  Its embedded pencil line is essential.

The preceding relative caution concerns the valuation-ring proof. The following polynomial identity establishes the equality over arbitrary complex base algebras.


\section{Descent and complete-local details}
\label{sec:formal-details-v154}
The following statements supply the local-to-global steps used in
Sections~\ref{sec:formal-images-v153}--\ref{sec:reciprocal-v153}.
They also specify which assertions concern a completed image, a
normalization fibre, or a family on an exact-divisor stratum.

\begin{lemma}[Complete-local surjectivity and reduced completion]
\label{lem:complete-local-v154}
In the cotangent-surjectivity assertion of
Lemma~\ref{lem:formal-image-v153}, the rings are Noetherian complete
local $\C$-algebras with residue field $\C$, and the homomorphism
is local, continuous, and the identity on that residue field. Under
these hypotheses cotangent surjectivity implies ring surjectivity.
A reduced local ring essentially of finite type over $\C$ has
reduced completion.
\end{lemma}
\begin{proof}
Choose preimages in the source maximal ideal of generators of the
target cotangent space. Their degree-$q$ products span every quotient
$\mathfrak m_B^q/\mathfrak m_B^{q+1}$. Given an element of $B$, first
lift its residue, then its error successively in degrees $q=1,2,\ldots$.
The correction in degree $q$ lies in $\mathfrak m_R^q$. Their sum
converges in $R$, and continuity and separatedness imply that its
image is the desired element of $B$. This is not a statement about
arbitrary noncomplete local homomorphisms.

Finite-type complex rings are excellent. Quasi-excellent local
rings are Nagata and have geometrically reduced formal fibres;
see \cite[Tags 07QV and 0BJ0]{StacksLocalV154}. For a reduced local
ring $A$, the injection into the product of its finitely many
minimal-prime quotients remains injective after flat completion.
Each quotient is a Nagata domain and its completion is reduced.
Thus $\widehat A$ is a subring of a product of reduced rings.
This justifies the reduced completed-image step of the earlier lemma.
\end{proof}

\begin{lemma}[Coprime cluster products over Artin bases]
\label{lem:hensel-products-v154}
Let $A$ be a local Artin $\C$-algebra. If the reduction of a monic
polynomial has a factorization into pairwise coprime monic factors,
there is a unique factorization into monic lifts of the same degrees.
In a factorization fibre, the functor of divisors with prescribed
degrees in the distinct residual clusters is the product of the
cluster divisor functors. Its coordinate algebra is therefore the
tensor product of their Artin coordinate algebras.
\end{lemma}
\begin{proof}
For two coprime factors $p,q$, the linear correction map
\[
 (\dot p,\dot q)\longmapsto q\dot p+p\dot q,
 \quad\deg\dot p<\deg p,\quad\deg\dot q<\deg q,
\]
is an isomorphism on coefficient spaces over the residue field.
Its kernel is zero by coprimeness, and source and target have equal
dimension. It is therefore an isomorphism over $A$ as well. Across
a square-zero extension of Artin rings it uniquely corrects a
lifted product to the given polynomial. Induction on the nilpotence
order and then on the number of clusters proves existence and
uniqueness. Apply the same argument to both the marked divisor and
its complement. The cluster product and decomposition constructions
are mutually inverse on every Artin base; Yoneda gives the asserted
product of finite schemes. Passing to inverse limits gives the
corresponding assertion on complete local bases.
\end{proof}

\begin{lemma}[The local choice and descent in the binary stratum]
\label{lem:residual-selection-v154}
The base-change identification of
Theorem~\ref{thm:binary-fibres-v153} holds as an identification of
incidence functors, not only on closed fibres. The required choice
of a member of the residual space coprime to the common divisor can
be made on a Zariski open cover of the exact-gcd stratum; the choices
are auxiliary and the resulting identifications glue uniquely.
\end{lemma}
\begin{proof}
Write a point of the stratum as $(G,L)$ in
$\Pj(S_s)\times X_{D-s,r}$. The vanishing of a member of $L$ at
one of the finitely many roots of $G$ is a proper linear condition
on that member. Over $\C$ a finite union of these proper hyperplanes
does not cover $L$. Choose a linear combination avoiding them, lift
its coefficients on a bundle-trivializing neighbourhood, and invert
its resultant with $G$. This gives the required Zariski neighbourhood.
These neighbourhoods cover the stratum and may be pulled back by
any base change.

On the marked-divisor incidence, pass locally to a monic coordinate
for the proposed divisor $f$. The quotient by $f$ is finite free.
On every geometric fibre $f$ divides $G$, so the chosen residual
member is invertible in that quotient, by coprimeness. Its
multiplication matrix has invertible determinant locally, hence it
is invertible also over a nonreduced base. Thus $f\mid G L$ implies
$f\mid G$ as an equation on that base, and the converse is immediate.
Monic division forces the remaining subspace. The identification is
independent of which coprime residual member established cancellation,
so it glues on overlaps. No lifting of arbitrary image points to the
normalization is asserted.

The finite factorization map used there is flat by miracle flatness:
its target local rings are regular, its source local rings are
Cohen--Macaulay, it is finite, and dimensions agree. These are exactly
the hypotheses of \cite[Tag 00R4]{StacksLocalV154}. This names the
flatness input separately from the incidence-functor argument.
\end{proof}

\begin{lemma}[The coordinate equalizer and coherent conductor]
\label{lem:equalizer-v154}
For the block ideals $I_A$ of Theorem~\ref{thm:squarefree-v153},
the completed image ring is the equalizer of the branch rings with
their restrictions to all pairwise intersections. This assertion
uses the coordinate monomial structure. The completed conductor
calculation there implies equality of the corresponding coherent
ideal sheaves on the squarefree open.
\end{lemma}
\begin{proof}
A monomial has a coefficient on each branch where its support
survives. If it survives on two branches, it also survives on their
coordinate intersection. Pairwise agreement therefore assigns it
one well-defined coefficient. Collecting these coefficients gives
one formal series modulo the intersection of the $I_A$; conversely
such a series has compatible restrictions. Formal convergence causes
no issue, since there are finitely many monomials in each degree.
The claim would be false for general configurations of ideals without
this coordinate argument.

Write $A\subset B$ for the local finite normalization algebra.
The conductor is $\Ann_A(B/A)$, a coherent ideal. Flat completion
preserves the exact finite-module sequence and its annihilator:
compute the latter as the intersection of the kernels of multiplication
on a finite generating set. The completed finite algebra is the
product of the normal branch rings. Thus the conductor and the ideal
of $Y_{g+1}$ have equal completed stalks at every squarefree closed
point. Faithful flatness gives equal stalks, and coherence on a
finite-type complex scheme gives equality of the ideal sheaves.
\end{proof}

\begin{proposition}[Finiteness and uniqueness of the collision diagram]
\label{prop:atlas-functor-v154}
Every branch map in Theorem~\ref{thm:atlas-v153} is finite, and the
completed incidence diagram is determined, up to the unique
isomorphism respecting its incidence data, by its marked-divisor
functor. It is independent of the chosen frame and root coordinates.
\end{proposition}
\begin{proof}
For each cluster $P_i=A_iB_i'$, the coefficients of the monic $A_i$
are integral over the coefficients of $P_i$. This follows from the
finite factorization morphism, or by expressing them as elementary
symmetric functions of a selected subset of the roots in a finite
splitting algebra. Division by a monic polynomial expresses every
coefficient of $B_i'$ polynomially in those of $A_i$ and $P_i$.
Likewise $R_{ij}=A_iU_{ij}$ expresses all quotient coefficients in
terms of $A_i$ and $R_{ij}$. Consequently the branch algebra is
generated over the ambient coefficient algebra by finitely many
integral coefficients of the $A_i$, so it is finite. Completing
preserves this finite-module assertion.

The equations represent the functor of a subbundle with a marked
divisor and its uniquely determined quotient. A change of frame or
local coordinate gives inverse natural transformations of this same
functor and of its forgetful map to the ambient Grassmannian.
The complete representing rings and their maps are uniquely
isomorphic by Yoneda. Their scheme-theoretic images consequently
agree as completed diagrams. This proves coordinate independence
without identifying arbitrary Artin points of an image with lifted
Artin points of its normalization.
\end{proof}

\paragraph{Partition closures.}
For a fixed partition $(s_1,\ldots,s_b)$, the proper map
$(\Pj^1)^b\to\Pj(S_s)$ sending $(\ell_i)$ to
$\prod\ell_i^{s_i}$ has closed image. Its open set of distinct
$\ell_i$ is dense. At the boundary precisely some of these points
coincide, and the associated parts add. Every such coincidence can
be approached from distinct points. This proves exactly the merging
closure order in the fixed-degree symmetric product; compare the
coincident-root constructions of \cite{Chipalkatti2003,Kurmann2012}.
It is not a closure classification for all analytic singularities
of the completed collision atlas.

\paragraph{Equivariance of the singular Fitting calculation.}
The group used in Theorem~\ref{thm:fitting-v153} is
$\operatorname{GL}_m(\C)$, acting simultaneously on the columns
$u$ and $w$ and fixing $\Delta$. It preserves the presentation and
therefore its intrinsic Fitting ideal. In the normalization the
three displayed minors span, under this group, respectively
$\Sym^m(\C^m)t$, $\Delta\Sym^m(\C^m)$, and
$\Sym^{m+1}(\C^m)$, with the coordinate/dual convention determined
by the action on $u$. To see spanning directly, their translates
are nonzero scalar multiples of the corresponding powers of all
linear forms; polarization in characteristic zero spans each
symmetric power. Multiplication by $R$ gives exactly the modules
asserted in \eqref{eq:fitting-v153}. This is stronger than invoking
invariance without identifying the representations.

\paragraph{Specialization of minimal relations.}
For the surjective coefficient specialization $\pi:P\to P'$ to
one coordinate in Theorem~\ref{thm:reciprocal-structure-v153}, one
has $\pi(\mathfrak n)=\mathfrak n'$ and
$\pi(\mathcal I)=\mathcal I'$. Consequently
$\pi(\mathfrak n\mathcal I)=\mathfrak n'\mathcal I'$.
A coefficient whose image is a minimal binary relation cannot lie
in $\mathfrak n\mathcal I$. This explicitly justifies the nonzero
irreducible coefficient map used in that proof. All reciprocal
series and socle weights in these arguments use the weight of a
$Q_i$ coefficient equal to $i$. They do not claim ordinary
maximal-ideal Hilbert functions; \eqref{eq:F22-local-hilbert-v154}
gives one such function separately.

% BEGIN PRESERVED PART 44-conductor-and-collisions-v152.tex

\section{Conductors and local models of divisor collisions}
\label{sec:conductor-v152}

We retain the notation of Section~\ref{sec:canonical-boundary-v151}.
All completed local rings below are completed at closed complex points.
For a reduced local ring $A$ with finite normalization $B$, its
\emph{conductor} is $\mathfrak c=\Ann_A(B/A)$, regarded also as an
ideal of $B$.  We use the intrinsic singular subscheme defined by
$\Fitt_{\dim A}\Omega_{A/\C}$; on completions, differentials mean
continuous differentials.  The conductor and this Fitting ideal need
not agree.  Our first result determines both on a multivariate
stratum.  The second identifies the way two branches merge when a
common binary root becomes double.

The formal-image and separated-germ steps below are isolated in
Lemma~\ref{lem:formal-image-v153}. Lemma~\ref{lem:coprime-lifts-v153}
gives the Artin cancellation and a direct positive lower bound for the
split codimension. Theorem~\ref{thm:atlas-v153} supplies the full
Weierstrass incidence functor in every collision order;
Theorem~\ref{thm:fitting-v153} evaluates the singular ideal for all ranks.

\subsection{Two coprime divisor choices}
\begin{theorem}[The split-divisor stratum]\label{thm:split-conductor-v152}
Suppose $h>g$, $2\leq r\leq s_{h-g}$, and
\[
 K=abL,\qquad [a],[b]\in\Pj(S_g),\qquad
 L\in\Gr(r,S_{h-g}),\qquad \gcd(L)=1.
\]
Assume that $a,b$ are coprime and that the only degree-$g$ divisors of
$ab$, up to scalar, are $a$ and $b$.  Put
\begin{align*}
 t&=2(s_g-1)+r(s_{h-g}-r),\\
 c&=r(s_h-s_{h-g})-(s_g-1),\qquad d_g=t+c.
\end{align*}
Then $c\geq1$, and there are formal coordinates in which
\begin{equation}\label{eq:split-ring-v152}
 \widehat{\OO}_{Y_g,K}=
 \C[[z_1,\ldots,z_t,x_1,\ldots,x_c,y_1,\ldots,y_c]]/
                       (x_i y_j:1\leq i,j\leq c).
\end{equation}
Its normalization is
$\C[[z,x]]\times\C[[z,y]]$.  The conductor is
\begin{equation}\label{eq:split-conductor-v152}
 \mathfrak c=(x_1,\ldots,x_c,y_1,\ldots,y_c),\qquad
 \widehat{\OO}_{Y_g,K}/\mathfrak c=\C[[z]],
\end{equation}
and the intrinsic singular subscheme has ideal $\mathfrak c^{\,c}$.
The two smooth branches meet scheme-theoretically in a smooth
$t$-dimensional subscheme.  The local ring has multiplicity two and
depth $t+1$; it is Cohen--Macaulay exactly when $c=1$.
\end{theorem}

\begin{proof}
We first determine the intersection as a scheme, not by comparing
closed points or tangent dimensions.  The normalization has precisely
two points above $K$, namely $(a,bL)$ and $(b,aL)$.  At either point,
the overlap in Lemma~\ref{lem:fibre-equations-v151} is zero.  The
map from the completed local ring of the ambient Grassmannian to
each completed source ring is consequently surjective: its cotangent
map is surjective, and completeness and Nakayama's lemma lift
successive homogeneous terms.  Thus each source completion is a
closed smooth formal branch of the image.

Here is the required assertion over infinitesimal bases.  Let $A_0$
be a local Artin $\C$-algebra.  A point on both formal branches carries
unique marked lifts $a_{A_0},b_{A_0}$ near $a,b$.  If homogeneous
polynomials have coprime reductions, those reductions form a regular
sequence in the polynomial ring over $\C$.  Its graded Koszul
complex is exact.  In each degree the polynomial modules are finite
free; lifting a splitting of the complex, or applying the local
criterion for flatness successively, shows that its lift over $A_0$
is exact and that the quotient is $A_0$-flat.  In particular
\begin{equation}\label{eq:coprime-artin-v152}
 (a_{A_0})\cap(b_{A_0})=(a_{A_0}b_{A_0}).
\end{equation}
This also follows by noting that $b_{A_0}$ is a nonzerodivisor
modulo $a_{A_0}$ and solving a relation between them.  Applying
\eqref{eq:coprime-artin-v152} to the columns of the relation
subbundle gives, uniquely,
\[
                  K_{A_0}=a_{A_0}b_{A_0}L_{A_0}.
\]
The residual module is a subbundle: multiplication by
$a_{A_0}b_{A_0}$ is a split injection in degree $D$, and the
quotient of the ambient free module by $K_{A_0}$ is free.  The
kernel of the induced surjection onto the complementary quotient
is therefore free.  This proves that the formal intersection is
represented by the product of the two divisor charts and the
Grassmannian chart for $L$, with no nilpotents omitted.

The differential of this product parametrization is injective.
Indeed, an infinitesimally constant product gives
\[
       (\dot a b+a\dot b)l+ab\dot l\in abL\quad(l\in L).
\]
Taking valuations at the factors of $a$, and using
$\gcd(a,b)=\gcd(L)=1$, forces $a\mid\dot a$; similarly
$b\mid\dot b$.  These are the two projective scalar directions,
and then $\dot L=0$.  The intersection is thus a smooth formal
closed subscheme of dimension $t$.

Let $R_0$ be the ambient completed local ring and $I_a,I_b$ the
ideals of these branches.  Finiteness of the normalization, exactness
of completion on finite modules, and the definition of the
scheme-theoretic image give
\[
 \widehat{\OO}_{Y_g,K}=R_0/(I_a\cap I_b)
   = (R_0/I_a)\mathbin{\times}_{R_0/(I_a+I_b)}(R_0/I_b).
\]
This is the elementary conductor-square construction; see also
\cite[Lemma~1.3 and \S1.3.2]{Ferrand2003}.  All three quotient rings
are regular.  Choosing compatible regular parameters for the two
surjections gives \eqref{eq:split-ring-v152}.  The branch dimension
is $d_g=s_g-1+r(s_h-r)$, so its codimension along the intersection
is the asserted $c$.  It is positive: if $c=0$, both branch ideals
would coincide locally with the intersection ideal, contradicting
the two distinct normalization branches of a reduced local ring.

In the fibre-product description, an element belongs to the
conductor exactly when both its restrictions to the intersection
vanish.  This gives \eqref{eq:split-conductor-v152}.  The Jacobian
of the relations $x_i y_j$ has nonzero entries only among the $x_i$
and $y_j$.  The Fitting ideal in question is its ideal of $c$-minors.
Every such minor belongs to $\mathfrak c^{\,c}$.  Conversely,
choosing the $c$ columns indexed by the $y_j$ and one row $(i_j,j)$
for each $j$ gives every monomial $x_{i_1}\cdots x_{i_c}$.
The symmetric choice gives every degree-$c$ monomial in the $y$'s.
Mixed monomials vanish in the ring.  Thus the ideal is exactly
$\mathfrak c^{\,c}$, including its nonreduced structure.

Finally use
\[
 0\longrightarrow \widehat{\OO}_{Y_g,K}
 \longrightarrow\C[[z,x]]\oplus\C[[z,y]]
 \longrightarrow\C[[z]]\longrightarrow0,
\]
where the last arrow is the difference of restrictions.  The depth
lemma gives depth $t+1$; for $c=1$ this equals the dimension.
For $c>1$ it is strictly smaller.  Additivity of leading Hilbert
coefficients gives multiplicity two, since the last module has
smaller dimension than the two regular branches.
\end{proof}

\begin{remark}\label{rem:split-scope-v152}
The hypotheses are nonempty in every number of variables: for $e=2$
one may take $g=1$, and for $e\geq3$ one may take two distinct
irreducible forms of degree $g$.  The theorem treats an explicitly
specified stratum, not all pairs of overlapping degree-$g$ divisors.
In particular, it does not substitute a binary factorization rule
for multivariate divisibility.  Its high-codimension branch
intersections need not be normal-crossing divisors.
\end{remark}

\subsection{A double common root}
For $m\geq1$, put $R_m=\C[[u_1,\ldots,u_m,\Delta]]$ and
$I=(u_1,\ldots,u_m)\subset R_m$.  Define
\begin{equation}\label{eq:pinch-algebra-v152}
 A_m=R_m\oplus I t\ \subset\ B_m=\C[[u_1,\ldots,u_m,t]],
                    \qquad \Delta=t^2.
\end{equation}
The right side of \eqref{eq:pinch-algebra-v152} is a subring,
not a square-zero extension: $(It)^2=\Delta I^2$.  For $m=1$
it is the ordinary pinch-point ring.

\begin{theorem}[Collision of two binary divisor choices]
\label{thm:binary-pinch-v152}
Let $e=2$, $g=1$, $D=h+1$, and $2\leq r\leq D-1$.
At a point $K_0=\ell^2 L$ with
$L\in\Gr(r,S_{D-2})$ and $\gcd(L)=1$, set
\[
             m=r-1,\qquad q=r(D-1-r)+1.
\]
Then
\begin{equation}\label{eq:binary-completion-v152}
              \widehat{\OO}_{Y_1,K_0}\simeq A_m[[z_1,\ldots,z_q]].
\end{equation}
Writing $w_i=u_i t$, a complete presentation of $A_m$ is
\begin{equation}\label{eq:pinch-equations-v152}
 \C[[u_1,\ldots,u_m,\Delta,w_1,\ldots,w_m]]/
 \bigl(u_iw_j-u_jw_i,\ w_iw_j-\Delta u_i u_j\bigr),
\end{equation}
where the first relations have $i<j$ and the second have $i\leq j$.
Its normalization is $B_m$, and its conductor is
\[
 \mathfrak c=(u_1,\ldots,u_m,w_1,\ldots,w_m),\qquad
 \mathfrak c B_m=I B_m.
\]
The induced map on conductor quotients is
\begin{equation}\label{eq:conductor-cover-v152}
              \C[[\Delta]]\longrightarrow\C[[t]],\qquad\Delta\longmapsto t^2.
\end{equation}
Thus the two reduced divisor lifts merge through a ramified double
cover of the conductor.  The completed local ring has dimension
$q+m+1$ and depth $q+2$.  It is Cohen--Macaulay if and only if $r=2$.
For $r=2$, its intrinsic singular ideal, with the smooth variables
suppressed, is $(w,\Delta u,u^2)$, whereas its conductor is $(u,w)$.
\end{theorem}

\begin{proof}
Choose an affine coordinate $x$ with the root of $\ell$ at zero.
Some member $F_0\in K_0$ has multiplicity exactly two there, since
$\gcd(L)=1$.  Extend it to a frame $F_0,\ldots,F_{r-1}$.
In the completed local ring of the ambient Grassmannian the formal
Weierstrass theorem writes the universal $F_0$ as a unit times
\[
                         P(x)=x^2+a x+b.
\]
Divide each $F_i$, $i>0$, by this monic polynomial.  Its remainder
is $u_i x+v_i$.  The $2r$ functions
$a,b,u_1,v_1,\ldots,u_m,v_m$ are part of a regular parameter
system.  To see this rather than assume versality, note that
$K_0$ has zero two-jet at the root, while
$S_D/K_0\to\C[x]/(x^2)$ is surjective.  The Grassmannian tangent
space $\Hom(K_0,S_D/K_0)$ therefore allows independent variations
of both jets in every frame column.  Passing from these jets to
the displayed remainders is an invertible triangular operation,
with leading coefficient $F_0/x^2(0)\ne0$.  Their differentials
are consequently independent.

Set
\[
       t=x+a/2,\quad \Delta=a^2/4-b,\quad
       w_i=a u_i/2-v_i.
\]
The incidence of a common root near zero is now exactly
\[
                       \Delta=t^2,\qquad w_i=u_i t.
\]
There are no other normalization points above $K_0$, because its
exact common divisor is $\ell^2$.  Consequently the image of this
completed incidence is the completed image defining $Y_1$.
The ambient dimension is $N_0=r(D+1-r)$.  Of the $2r$ parameters
above, $a$ is free, as are the other $N_0-2r$ parameters.  This
accounts for $q=N_0-2r+1$ smooth variables in
\eqref{eq:binary-completion-v152}.

It remains to justify elimination scheme-theoretically.  The
relations in \eqref{eq:pinch-equations-v152} reduce any series to
$p+\sum_i w_i p_i$ with $p,p_i\in R_m$.  If its image in $B_m$
is zero, the decomposition $B_m=R_m\oplus R_m t$ gives
$p=0$ and $\sum_i u_i p_i=0$.  The first syzygies of the regular
sequence $u_1,\ldots,u_m$ are generated by the Koszul syzygies.
They are precisely the relations $u_iw_j-u_jw_i$.  This proves
the presentation without taking a radical.  The same argument
works first over polynomial rings and then by exact completion.

The extension $A_m\subset B_m$ is finite, since $t^2=\Delta$,
and birational, since $t=w_1/u_1$ in the fraction field.  The
regular ring $B_m$ is thus the normalization.  For
$p+v t\in R_m\oplus I t$, multiplication by $t$ stays in
$A_m$ if and only if $p\in I$.  As $1,t$ generate $B_m$,
the conductor is $I\oplus I t$, proving the quotient map
\eqref{eq:conductor-cover-v152}.  Equivalently,
\begin{equation}\label{eq:pinch-square-v152}
 A_m=B_m\mathbin{\times}_{\C[[t]]}\C[[\Delta]],
\end{equation}
where the first map sets $u=0$ and the second sends $\Delta$ to
$t^2$.  This is a specific conductor square, not an assumption
that the normalization descends through its reduced fibres.

As an $R_m$-module, $A_m=R_m\oplus I$.  For $m=1$, the ideal
$I$ is free.  For $m>1$, the exact sequence
$0\to I\to R_m\to\C[[\Delta]]\to0$ gives
$\operatorname{depth}_{R_m}I=2$.  Hence $A_m$ has depth two
for every $m\geq1$.  The maximal ideal of $R_m$ generates an
ideal with maximal radical in $A_m$, so this is also its local
ring depth.  Adding the $q$ smooth parameters gives the assertion.
When $m=1$, the equation is $w^2-\Delta u^2=0$; its partial
derivatives generate $(w,\Delta u,u^2)$ in characteristic zero.
\end{proof}

\begin{corollary}[The split and ramified parts of the same model]
\label{cor:collision-strata-v152}
On the binary locus where the common divisor has degree exactly two,
there are two strata.  Over two distinct roots the formal local ring
has the split form \eqref{eq:split-ring-v152}, with
$t=q+1$ and $c=r-1$.  Over a double root it has
\eqref{eq:binary-completion-v152}.  The support of the singular
locus in the latter model is the conductor, and the normalization
fibre at a collision is $\Spec\C[t]/(t^2)$.
For arbitrary $m$, its singular subscheme is given without
radicalization by the $m$-minors of the Jacobian of the displayed
$m^2$ relations in \eqref{eq:pinch-equations-v152}.
\end{corollary}
\begin{proof}
The assertion at distinct roots is Theorem~\ref{thm:split-conductor-v152}.
In the collision model the normalization is an isomorphism wherever
some $u_i$ is invertible.  On $u=w=0$, a nonzero value of $\Delta$
gives two distinct points after an \'etale square-root extension,
while $\Delta=0$ gives the stated length-two fibre.  The exact
normality criterion proves the assertion about singular support.
There are $2m+1$ nonsmooth coordinates and dimension $m+1$;
the cotangent presentation therefore gives the $m$-minors as
claimed.  This formula specifies the whole singular scheme, even
where it is thicker than its conductor support.
\end{proof}

\begin{remark}[Descent to intrinsic first relations]
\label{rem:conductor-stack-v152}
On the maximal-lower-Hilbert-function stratum of
Corollary~\ref{cor:algebra-boundary-v151}, these statements compute
the completed smooth-atlas singularities of the intrinsic
first-relation algebra stack after tangent-identity rigidification.
The conductor is equivariant, and formation of its finite-module
annihilator and of the Fitting ideals commutes with flat base change.
Thus the calculations descend through the frame atlas.  This does
not identify these rings with completed local rings of a coarse
orbit space, where stabilizers can introduce further singularities.
\end{remark}

% BEGIN PRESERVED PART 46-divisor-intersections-v152.tex

\section{Intersections of divisor-degree images}
\label{sec:degree-intersections-v152}

Here the total relation degree $D$ is fixed.  Write $Y_a\subset
\Gr(r,S_D)$ for the image admitting a divisor of degree $a$;
empty Grassmannians are omitted.  In several variables these images
are not generally nested.  Their reduced intersections nevertheless
have a finite description by common parts of two marked divisors.

\begin{proposition}[The common-part formula]
\label{prop:degree-intersections-v152}
For $1\leq a,b<D$ and $0\leq j\leq\min(a,b)$, let $M_j$ be the
scheme-theoretic image of
\begin{equation}\label{eq:intersection-map-v152}
\begin{split}
 \Pj(S_j)\times\Pj(S_{a-j})\times\Pj(S_{b-j})
       \times\Gr(r,S_{D-a-b+j})&\longrightarrow\Gr(r,S_D),\\
                   (c,u,v,L)&\longmapsto cuvL.
\end{split}
\end{equation}
A negative last degree or insufficient dimension gives the empty
image, and $\Pj(S_0)$ is a point.  Then
\begin{equation}\label{eq:intersection-reduced-v152}
             (Y_a\cap Y_b)_{\mathrm{red}}
                         =\left(\bigcup_j M_j\right)_{\mathrm{red}}.
\end{equation}
On the locus of two marked divisors $f_a,f_b$ with exact common
part of degree $j$, the corresponding factorization is
$c=\gcd(f_a,f_b)$, $f_a=cu$, $f_b=cv$, with $\gcd(u,v)=1$.
In particular, at a relation $K=a_0^mL$ with $a_0$ a nonzero
linear form and $\gcd(L)=1$, the
set of positive divisor degrees is exactly $\{1,\ldots,m\}$.
For the boundary point in Theorem~\ref{thm:universal-boundary-v152},
$m=D-q_n$.
\end{proposition}
\begin{proof}
Each map in \eqref{eq:intersection-map-v152} is defined on a
projective scheme: multiplication by a nonzero form has constant
rank, so the subspace construction is a morphism.  Its image is
closed and admits the two divisors $cu$ and $cv$, proving one
inclusion.  Conversely, if $K$ has divisors $f_a,f_b$, unique
factorization gives their common part $c$ of some degree $j$.
Their least common multiple is $cuv$ and divides every member of
$K$.  Division yields $L$ of degree $D-a-b+j$, of the same rank.
This proves equality on geometric points.  Over $\C$, a reduced
closed subscheme of a finite-type scheme is determined by its closed
points, giving \eqref{eq:intersection-reduced-v152}.  The last
assertion follows because every homogeneous divisor of a power
of one linear form is a power of that form.
\end{proof}

The reduction in \eqref{eq:intersection-reduced-v152} is essential
to the statement of this global formula.  It is not a prescription
to reduce the individual $Y_a$, their normalization fibres, or their
singular subschemes.  The full intersection of the two formal
branches on the coprime stratum was computed over Artin bases in
Theorem~\ref{thm:split-conductor-v152};
Theorem~\ref{thm:binary-pinch-v152} computes the thickening when
those branches collide.  Thus the degree-poset description and
the local scheme calculations have distinct, explicit roles.

% BEGIN PRESERVED PART 45-pencil-orbit-boundary-v152.tex

\section{A common ramified boundary of pencil failure orbits}
\label{sec:orbit-boundary-v152}

The fixed tensor presentation of a pencil has exact common divisor
$\delta^{n-1}$.  Varying the pencil alone therefore stays in the
exact-divisor locus.  The intrinsic finite algebra permits a larger,
linear coordinate-change action: changing the basis of its entire cotangent
space.  We now describe a boundary point produced by such changes.
The generic members are genuine pencil failure algebras.  The central
member has a larger common divisor, and its normalization fibre is
nonreduced.  Thus the boundary of Section~\ref{sec:canonical-boundary-v151}
occurs inside the orbit closures of the motivating algebras themselves.

Let $n=\dim V=\dim U\geq3$, $E=V^*\otimes U$, and put
\[
 e=n^2,\quad N=\binom{n+1}{2},\quad r=\binom N2,\quad
 g=n(n-1),\quad h=3n-4,\quad D=g+h.
\]
For a pencil $R$, write $K_R=\delta^{n-1}C_R\subset\Sym^D E$
as in Lemma~\ref{lem:primitive-pencil}.  Define the projective orbit
closure
\[
 Z_R=\overline{\operatorname{GL}(E)\cdot K_R}
                  \ \subset\ \Gr(r,\Sym^D E).
\]
It is independent of a chosen basis of the cotangent space.  Since
$Y_g$ is closed and invariant, $Z_R\subset Y_g$.

The group in $Z_R$ is the full $\operatorname{GL}(E)$, not the
$\PGL(V)$ congruence group of the pencil. Fixing the standard flag,
scalar direction and complement specifies the displayed common point;
no choice-independent distinguished moduli point is asserted.

\begin{theorem}[A common point in full first-relation orbit closures]
\label{thm:universal-boundary-v152}
Set
\[
 q_n=\begin{cases}3,&n=3,\\4,&n\geq4,\end{cases}
 \qquad k_n=h-q_n.
\]
There is an explicitly constructed point $K_\infty$ belonging to
$Z_R$ for every quadratic pencil $R$, with
\begin{equation}\label{eq:boundary-gcd-v152}
 K_\infty=a^{g+k_n}L_n,\qquad
 L_n\subset\Sym^{q_n}E,\quad \rank L_n=r,\quad\gcd(L_n)=1,
\end{equation}
where $a\in E$ is a nonzero linear form.  The flag pencil
$R_*=\langle x_1^2,x_1x_2\rangle$ admits a finite-flat family of
first-relation algebras over $\A^1$, with special relation space
$K_\infty$ and every nonzero fibre isomorphic to its failure algebra.
Together with Theorem~\ref{thm:closed-pencil-v151}, this gives a
chain of two explicit finite-flat specializations from the failure
algebra of every $R$ to the same boundary algebra.

The fibre of $\nu_g$ over $K_\infty$ has a single geometric point,
marked by $a^g$, and is nonreduced.  Its tangent dimension is
\begin{equation}\label{eq:boundary-tangent-v152}
                 \binom{e+k_n-1}{k_n}-1.
\end{equation}
In particular $K_\infty$ lies in the singular locus of $Y_g$ and
outside the exact-degree-$g$ locus.  For $n=3$ this tangent dimension
is $44$.
\end{theorem}

The construction requires the exact osculating order of one elementary
flag embedding.  We include the calculation to distinguish the theorem
from a computationally suggested degeneration.

\begin{lemma}[The flag coefficient filtration]
\label{lem:flag-jets-v152}
Write a matrix of independent linear coordinates as
\[
                    T=aI_n+B,\qquad B_{11}=0.
\]
For $C_*=C_{R_*}$, the filtration by vanishing order at $T=I_n$
on the chart $T_{11}=1$ has last nonzero graded piece in degree $q_n$.
It has a nonzero degree-zero piece.  Hence an adapted basis
$c_1,\ldots,c_r$ has orders $0\leq d_j\leq q_n$, attaining both
endpoints, and
\begin{equation}\label{eq:flag-initial-v152}
 \lim_{\tau\to0}\tau^{-d_j}c_j(aI_n+\tau B)
                           =a^{h-d_j}p_j(B),
\end{equation}
where the displayed initial forms are linearly independent.
\end{lemma}
\begin{proof}
The complementary-minor identity for $\Sym^2T$ identifies the vector
of primitive coefficients, up to a fixed invertible coordinate
change, with
\begin{equation}\label{eq:flag-coefficients-v152}
 \delta(T)^3\bigl(v^2\wedge vw\bigr),\qquad
                 v=T^{-1}e_1,\quad w=T^{-1}e_2.
\end{equation}
Indeed $\det(\Sym^2T)=\delta^{n+1}$, the complementary minors have
order two, and division of $\delta I_{N-2}$ by $\delta^{n-1}$ leaves
$\delta^3$.  This identity is initially on $\delta\ne0$; both sides
are the same polynomial vector of degree $h$.

A standard affine flag chart has representatives
\[
       v=e_1+\sum_{i=2}^n z_i e_i,
       \qquad w=e_2+\sum_{i=3}^n y_i e_i.
\]
Every coordinate of $v^2\wedge vw$ is a polynomial of degree at most
four.  The map from $\operatorname{GL}_n$ to this flag chart is
smooth, and restricting to $T_{11}=1$ remains smooth: the scalar
subgroup is contained in each flag stabilizer and is transverse to
this chart.  Normalizing $v,w$ only multiplies the whole wedge by
a common unit.  Smooth pullback preserves the order of a nonzero
function, as does multiplication by a unit.  The coordinates in
\eqref{eq:flag-coefficients-v152} are linearly independent by Lemma~\ref{lem:pencil-irreducibility} and the nonzero
matrix-coefficient inclusion in \eqref{eq:general-coefficient-map}.  Thus no nonzero linear
combination can have order exceeding four.  For $n\geq4$, the
coordinate indexed by $e_3^2\wedge e_3e_4$ is
\[
                       z_3^2(z_3y_4-z_4y_3),
\]
which is nonzero and has order exactly four.

For $n=3$, put $z=z_2$, $w=z_3$, $t=y_3$.  In the ordered basis
$e_1^2,e_1e_2,e_1e_3,e_2^2,e_2e_3,e_3^2$, the two factors are
\[
 (1,2z,2w,z^2,2zw,w^2),\qquad
 (0,1,t,z,zt+w,wt).
\]
Their fifteen wedge coordinates have the same linear span as
\begin{equation}\label{eq:ternary-flag-basis-v152}
\begin{gathered}
 1,\ z,\ w,\ t,\ z^2,\ zw,\ w^2,\ zt,\ wt,\ z^2t,\ zwt,\ w^2t,\\
 z^2w-z^3t,\quad zw^2-z^2wt,\quad w^3-zw^2t.
\end{gathered}
\end{equation}
This is obtained by subtracting the two coordinates with a common
quadratic term and rescaling by $2$.  Their lowest homogeneous terms
are distinct monomials, all of degrees at most three, and some have
degree three.  Thus the third jet is injective on this coefficient
space and the second is not.  This proves $q_3=3$.  The coordinate
$e_1^2\wedge e_1e_2$ is $1$ in every dimension, proving the
order-zero assertion.  Finally, homogeneity of degree $h$ gives
\eqref{eq:flag-initial-v152}; an adapted basis is by definition a
basis on each associated-graded piece, so its initial forms are
independent.
\end{proof}

\begin{proof}[Proof of Theorem~\ref{thm:universal-boundary-v152}]
The linear substitution $T\mapsto aI_n+\tau B$ is invertible for
$\tau\ne0$.  Divide the adapted coefficient $c_j$ by its exact
power $\tau^{d_j}$ after substitution.  There are no negative powers:
all its terms have $B$-degree at least $d_j$.  The resulting columns
span a rank-$r$ subbundle $J_\tau\subset\Sym^h E\otimes\C[\tau]$.
Independence at zero follows from the lemma, and away from zero from
the invertible substitution.  The ideal of its maximal minors in $\C[\tau]$ has no common
closed zero, and is therefore the unit ideal. On each maximal-minor
open a corresponding square submatrix is invertible and splits the
image. These opens cover $\operatorname{Spec}\C[\tau]$, proving the
subbundle assertion over the whole line, including all scheme points.  Set
\[
                 K_\tau=\det(aI_n+\tau B)^{n-1}J_\tau.
\]
Multiplication by the nonzero homogeneous polynomial in every fibre
again has constant rank.  Consequently $K_\tau$ is a subbundle, and
\[
 \mathcal A_\tau=\Sym E\otimes\C[\tau]\big/
                           (K_\tau,\mathfrak m^{D+1})
\]
is finite locally free of rank $\binom{e+D}{D}-r$.
For $\tau\ne0$, its relation space is a linear transform of $K_{R_*}$.

At zero the marked factor becomes $a^g$, while
$J_0=\langle a^{h-d_j}p_j(B)\rangle$.  It contains a nonzero multiple
of $a^h$, so its gcd is a power of $a$.  It also contains a polynomial
with exact $a$-adic order $h-q_n$, and every column is divisible by
that power.  Therefore $\gcd(J_0)=a^{k_n}$, proving
\eqref{eq:boundary-gcd-v152}.  Notice that $0<k_n<g$ for every
$n\geq3$.

The closed-flag specialization of Theorem~\ref{thm:closed-pencil-v151}
and its failure-algebra lift imply
$K_{R_*}\in Z_R$.  A closed invariant subset containing a point
contains that point's orbit closure.  Hence the just-constructed
$K_\infty\in Z_{R_*}$ belongs to every $Z_R$.  This proves the
assertion about the two-stage specialization without replacing it
by an unproved assertion about one fixed one-parameter subgroup
acting on every initial pencil.

The sole degree-$g$ divisor of $a^{g+k_n}$ is $a^g$.  Its overlap
with the residual gcd is $a^{k_n}$.  Lemma~\ref{lem:fibre-equations-v151}
gives \eqref{eq:boundary-tangent-v152}; the fibre is supported at one
point and has positive tangent space.  The normalization criterion
then proves singularity.  With $n=3$, $e=9$ and $k_n=2$, the dimension
is $\binom{10}{2}-1=44$.
\end{proof}

\subsection{The complete normalization fibre at the limiting algebra}
The last theorem is not restricted to the tangent space of its fibre.
The full local algebra can be given uniformly.

\begin{proposition}[Pure-power divisor fibre]
\label{prop:pure-power-fibre-v152}
Let $E=\C a\oplus W$, $g\geq k\geq1$, and
$K=a^{g+k}L$ with $\gcd(L)=1$.  Introduce coefficient variables for
$q_i\in\Sym^i W$, $1\leq i\leq k$, and define forms $c_j$ recursively by
\begin{equation}\label{eq:reciprocal-v152}
 c_0=1,\qquad c_j=-\sum_{i=1}^k q_i c_{j-i},\qquad c_j=0\ (j<0).
\end{equation}
The entire fibre of $\nu_g$ over $K$, supported at its unique point,
is the Artin algebra
\begin{equation}\label{eq:pure-power-fibre-v152}
 \C[[\text{coefficients of }q_1,\ldots,q_k]]\big/
     (\text{coefficients of }c_{g+1},\ldots,c_{g+k}).
\end{equation}
Its tangent dimension is $s_k(E)-1$.
\end{proposition}
\begin{proof}
Choose $p\in L$ not divisible by $a$.  Over any local Artin base,
a divisor $f$ near $a^g$ is monic in $a$, and the quotient by $f$
is flat over the base.  The reductions of $f,p$ are coprime.
The regular-sequence argument used in
\eqref{eq:coprime-artin-v152} makes $p$ a nonzerodivisor modulo $f$.
Thus $f\mid a^{g+k}p$ is equivalent to $f\mid a^{g+k}$, functorially
on these bases.  Once this condition holds, it holds for every
member of $K$.

Polynomial division gives a unique monic residual factor
$H=a^k+q_1a^{k-1}+\cdots+q_k$.  The equation $fH=a^{g+k}$ determines
the coefficients of $f$ successively as $c_1,\ldots,c_g$ in
\eqref{eq:reciprocal-v152}.  The remaining $k$ product coefficients
vanish if and only if $c_{g+1},\ldots,c_{g+k}$ vanish.  This follows
successively by comparing the product with the reciprocal series
$(1+\sum q_i a^{-i})^{-1}$.  Hence
\eqref{eq:pure-power-fibre-v152} represents the full infinitesimal
fibre, without reduction.  Finiteness of $\nu_g$ makes it Artin;
there is no other geometric divisor, so completion loses no part
of the fibre.  Since $g\geq k$, these equations have no linear
terms in the coefficient variables.  The number of coefficient variables is
$\sum_{i=1}^k\dim\Sym^iW=\dim\Sym^kE-1$, as asserted.
\end{proof}

\begin{example}[The nine-generator limit]\label{ex:n3-boundary-v152}
For $n=3$, the construction can be written with four of the nine
linear coordinates, denoted $a,u,v,w$.  The other five coordinates
remain generators of the algebra.  Up to nonzero scalar changes
of the relation basis, $K_\infty=a^8L_3$, where
\begin{equation}\label{eq:n3-limit-v152}
\begin{split}
 L_3=\langle& a^3,\ a^2u,\ a^2v,\ a^2w,\ au^2,\ auv,\ av^2,
             \ auw,\ avw,\\
            &u^2v,\ u^2w,\ uv^2,\ uvw,\ v^3,\ v^2w\rangle.
\end{split}
\end{equation}
Here one may take $u=B_{21}$, $v=B_{31}$, $w=B_{32}$; the monomials
are the lowest terms in \eqref{eq:ternary-flag-basis-v152}.
The relation space has dimension $15$ and exact gcd $a^8$.
The selected divisor has degree $g=6$, and its fibre has $44$
coefficient variables: eight coefficients of a linear form $Q_1$
and thirty-six of a quadratic form $Q_2$ in $W$.
Its defining ideal consists of all coefficients of the two forms
\begin{align*}
 c_7&=-Q_1^7+6Q_1^5Q_2-10Q_1^3Q_2^2+4Q_1Q_2^3,\\
 c_8&=Q_1^8-7Q_1^6Q_2+15Q_1^4Q_2^2-10Q_1^2Q_2^3+Q_2^4.
\end{align*}
This is a full presentation of the normalization fibre.  Its embedding
dimension $44$ is not being identified with its length, nor is this
fibre algebra being identified with the completed local ring of $Y_g$.
\end{example}

\begin{remark}[Relation to pencil invariants]\label{rem:segre-boundary-v152}
The action defining $Z_R$ is the intrinsic $\operatorname{GL}(E)$
action on first relations, not merely the congruence action on $V$.
The flag pencil is singular: every member has rank at most two.
It is therefore outside the regular-pencil open set treated by
\cite[Theorem~1.1 and \S5]{FevolaMandelshtamSturmfels}, where Segre
symbols organize classification and degenerations.  The theorem
above does not purport to classify that poset or all orbits in $Z_R$.
Its conclusion is instead a uniform common boundary and an exact
ramification algebra for genuine finite failure algebras.  The
increase from $g$ to $g+k_n$ in the intrinsic gcd is precisely what
is invisible in a fixed-determinant pencil specialization.
\end{remark}



% BEGIN PRESERVED PART 00-principal-introduction-v152.tex

% BEGIN PRESERVED PART 17-assistance.tex
\section*{Acknowledgment of assistance}
AI assistance was used in developing and drafting the arguments and
finite verification scripts.  All structural proofs are supplied for
independent mathematical scrutiny; computational receipts are not
formal proof certificates.  No AI system is listed as an author.

% BEGIN PRESERVED PART 43-assistance-specific-v151.tex
For the present revision, GPT-6 Astra Pro assisted specifically with the
incidence normalization and overlap criterion in
Section~\ref{sec:canonical-boundary-v151}, the universal-property proof
in Section~\ref{sec:rigidification-v151}, the explicit specialization
in Section~\ref{sec:closed-pencil-v151}, and the associated exact tests
and exposition.  These arguments require independent human verification.
No statement here represents a completed journal review or a formal
proof-assistant certification.

% BEGIN PRESERVED PART 47-assistance-specific-v152.tex
\paragraph{Additional disclosure for the present revision.}
GPT-6 Astra Pro assisted specifically with the conductor and collision
models, the scalar-matrix degeneration, the flag-jet calculation,
the reciprocal-polynomial fibre presentation, literature comparisons,
and exact computational checks.  The proofs above state the complete
mathematical arguments; finite computations are consistency checks,
not universal proof certificates.  No journal decision or historical
priority determination is inferred from the build and source audits.

\section*{Assistance in the present revision}
An AI assistant helped formulate and write the binary conductor stratification,
the simultaneous incidence atlas, the evaluated singular Fitting ideal, and
the minimal-relation-module argument for pure-power fibres. It also prepared
the source-preserving reorganization and finite exact consistency checks.
The written proofs, rather than those finite checks, are the mathematical
basis of the statements and are submitted to independent scrutiny.

\paragraph{Revision 154.}
AI assistance was used for the polarized moduli compatibility,
spectral divisor-fibre construction, homological calculations and
proof expansions. The fixed five-variable Betti table uses exact
computer-assisted rational elimination with its complete finite
inputs recorded. The general theorems rely on the written arguments
and the explicitly attributed classical results, not on regression
counts or source-preservation receipts.

\paragraph{Revision 155.}
AI assistance was used to formulate and draft the exact quadratic
coefficient-section theorem and the reciprocal power-ideal description
of spectral Fitting schemes, to check their finite examples, and to
prepare the companion-paper organization.  Classical apolar and
orthogonal-orbit results are explicitly attributed.  The general
claims rely on the written proofs; finite exact tests and source
preservation are not formal proof certificates.

\paragraph{Revision 157.}
AI assistance was used in deriving and drafting the universal power-ideal
formula, its complete-quadric graph interpretation, the contact-divisor
and resolved-pencil statements, and their exact finite checks. Classical
determinantal, apolar and complete-quadric results are attributed at their
points of use. The proofs, rather than the computations, support the
general statements. No formal proof certification or journal endorsement
is claimed.

\paragraph{Scope of the new intrinsic construction.}
The source-normalized envelope is an algebra bundle with an intrinsic
projective power diagram. It is not a claimed preferred subquotient
of the sharp failure algebra. Singular-pencil classification is read
through the classical complex symmetric Kronecker form. The explicit
Hilbert tail is a boundary calculation for the stated transverse
family, not a classification of every fibre of the incidence
compactification. Exact computations accompanying the source check
finite instances; they do not certify the general proofs.
\begin{thebibliography}{99}

\bibitem{Ballico93}
E.~Ballico,
On the failure locus of higher order properties of embeddings in
projective spaces,
\emph{Math. Nachr.} \textbf{163} (1993), 5--13.
doi:10.1002/mana.19931630102.

\bibitem{Ballico96}
E.~Ballico,
On the failure cycles for the quadratic normality of a projective
variety,
\emph{Pacific J. Math.} \textbf{172} (1996), 307--313.

\bibitem{ChengWang}
S.-J. Cheng and W. Wang, Howe duality for Lie superalgebras,
\emph{Compositio Math.} \textbf{128} (2001), 55--94;
arXiv:math/0008093v1, \S2 and Theorem 3.5.

\bibitem{FevolaMandelshtamSturmfels}
C.~Fevola, Y.~Mandelshtam, and B.~Sturmfels,
Pencils of quadrics: old and new,
\emph{Le Matematiche} \textbf{76} (2021), no.~2, 319--335.
doi:10.4418/2021.76.2.2; arXiv:2009.04334v2,
Theorem~1.1, Corollary~2.1, Remark~2.3, Conjecture~4.5, and Theorem~5.1.

\bibitem{Ohta1986}
T.~Ohta, The singularities of the closures of nilpotent orbits in certain
symmetric pairs, \emph{Tohoku Math. J. (2)} \textbf{38} (1986), 441--468.
doi:10.2748/tmj/1178228456.

\bibitem{SmithQuadrics}
I.~Smith, Floer cohomology and pencils of quadrics,
\emph{Invent. Math.} \textbf{189} (2012), 149--250.
doi:10.1007/s00222-011-0364-1; arXiv:1006.1099, \S4.2.

\bibitem{EliasRossiShort}
J.~Elias and M.~E.~Rossi,
Isomorphism classes of short Gorenstein local rings via Macaulay's
inverse system,
\emph{Trans. Amer. Math. Soc.} \textbf{364} (2012), no.~9, 4589--4604.
doi:10.1090/S0002-9947-2012-05430-4;
arXiv:0911.3565v1, Theorem~3.3 and Corollaries~3.4--3.5.

\bibitem{EliasRossiCompressed}
J.~Elias and M.~E.~Rossi,
Analytic isomorphisms of compressed local algebras,
\emph{Proc. Amer. Math. Soc.} \textbf{143} (2015), no.~3, 973--987.
doi:10.1090/S0002-9939-2014-12313-6;
arXiv:1207.6919v1, Theorem~3.1 and Examples~3.3--3.5.

\bibitem{MarcusMoyls}
M.~Marcus and B.~N.~Moyls,
Transformations on tensor product spaces,
\emph{Pacific J. Math.} \textbf{9} (1959), no.~4, 1215--1221.
doi:10.2140/pjm.1959.9.1215, Theorem~1.

\bibitem{StacksGroups}
The Stacks Project Authors,
\emph{The Stacks Project}, Section~39.8, Lemma~39.8.2,
\url{https://stacks.math.columbia.edu/tag/0BF6}, accessed September 24, 2026.
\bibitem{Milne2017}
J.~S. Milne, \emph{Algebraic Groups: The Theory of Group Schemes of
Finite Type over a Field}, Cambridge Studies in Advanced Mathematics
\textbf{170}, Cambridge University Press, Cambridge, 2017.

\bibitem{StacksClosedImmersion}
The Stacks Project Authors, \emph{The Stacks Project},
Lemma~41.7.2 (Tag 04XV),
\url{https://stacks.math.columbia.edu/tag/04XV}, accessed September 24, 2026.

\bibitem{StacksQuotient}
The Stacks Project Authors, \emph{The Stacks Project},
Theorem~97.17.2 and Section~97.18,
\url{https://stacks.math.columbia.edu/tag/06PI}, accessed September 24, 2026.

\bibitem{StacksSmoothQuotient}
The Stacks Project Authors, \emph{The Stacks Project},
Lemma~101.33.7 (Tag 0DLS),
\url{https://stacks.math.columbia.edu/tag/0DLS}, accessed September 24, 2026.

\bibitem{StacksFiniteV151}
The Stacks Project Authors, \emph{The Stacks Project},
Lemma~37.44.1 (Tag 02LS),
\url{https://stacks.math.columbia.edu/tag/02LS}, accessed September 24, 2026.

\bibitem{StacksImageV151}
The Stacks Project Authors, \emph{The Stacks Project},
Section~29.6, Scheme theoretic image (Tag 01R5),
\url{https://stacks.math.columbia.edu/tag/01R5}, accessed September 24, 2026.

\bibitem{StacksNormalV151}
The Stacks Project Authors, \emph{The Stacks Project},
Section~29.55, Normalization (Tag 035E),
\url{https://stacks.math.columbia.edu/tag/035E}, accessed September 24, 2026.

\bibitem{StacksRigidV151}
The Stacks Project Authors, \emph{The Stacks Project},
Subsection~112.5.7, Rigidification (Tag 04V2), including the discussion
of noncentral normal inertia subgroups,
\url{https://stacks.math.columbia.edu/tag/04V2}, accessed September 24, 2026.


\bibitem{Chipalkatti2003}
J.~V. Chipalkatti, On equations defining coincident root loci,
\emph{J. Algebra} \textbf{267} (2003), 246--271.
doi:10.1016/S0021-8693(03)00336-3;
arXiv:math/0110224v1, Proposition~5.1, Theorem~5.4,
and Corollary~5.8.

\bibitem{Kurmann2012}
S.~Kurmann, Some remarks on equations defining coincident root loci,
\emph{J. Algebra} \textbf{352} (2012), 223--231.
doi:10.1016/j.jalgebra.2011.10.045;
arXiv:1108.4532v1, Proposition~1.3, Theorem~1.5,
and Proposition~2.1.

\bibitem{Ferrand2003}
D.~Ferrand, Conducteur, descente et pincement,
\emph{Bull. Soc. Math. France} \textbf{131} (2003), no.~4, 553--585.
doi:10.24033/bsmf.2455.

\bibitem{AOV2008}
D.~Abramovich, M.~Olsson, and A.~Vistoli, Tame stacks in positive
characteristic, \emph{Ann. Inst. Fourier (Grenoble)}
\textbf{58} (2008), no.~4, 1057--1091.
doi:10.5802/aif.2378; Appendix~A, Theorem~A.1.

\bibitem{HuLinShao2007}
Y.~Hu, J.~Lin, and Y.~Shao, A compactification of the space of
algebraic maps from $\Pj^1$ to $\Pj^n$,
arXiv:math/0701255v1 (2007), Theorem~1.1,
\S2.2, and Corollary~3.6.


\bibitem{Grinberg2019}
D.~Grinberg, \emph{A basis for a quotient of symmetric polynomials},
arXiv:1910.00207, version of September 24, 2021, Theorem~2.7.

\bibitem{MFK1994}
D.~Mumford, J.~Fogarty, and F.~Kirwan,
\emph{Geometric Invariant Theory}, third edition,
Ergebnisse der Mathematik und ihrer Grenzgebiete (2), vol.~34,
Springer-Verlag, Berlin, 1994.

\bibitem{MillerSturmfels2005}
E.~Miller and B.~Sturmfels, \emph{Combinatorial Commutative Algebra},
Graduate Texts in Mathematics, vol.~227, Springer, New York, 2005.
Chapter~5 (Stanley--Reisner rings and Hochster's formula).

\bibitem{ShafieiSymmetric}
S.~Shafiei, \emph{Apolarity for determinants and permanents of generic
symmetric matrices}, arXiv:1303.1860v3 (2021), Proposition~3.1,
Theorem~3.10, and Corollary~2.7.

\bibitem{Krishna2026}
A.~Krishna, \emph{A universal discriminant formula for pencils of
quadrics}, arXiv:2607.05340v2 (2026), Theorem~1.1 and Lemma~3.1.

\bibitem{StacksLocalV154}
The Stacks Project Authors, \emph{The Stacks Project},
Tags \href{https://stacks.math.columbia.edu/tag/00R4}{00R4}
(miracle flatness),
\href{https://stacks.math.columbia.edu/tag/07QV}{07QV}
(quasi-excellent rings are Nagata), and
\href{https://stacks.math.columbia.edu/tag/0BJ0}{0BJ0}
(geometrically reduced formal fibres), accessed September 25, 2026.

\bibitem{NagaokaWachi2024}
T.~Nagaoka and A.~Wachi,
\emph{The strong Lefschetz property of Gorenstein algebras generated
by relative invariants}, arXiv:2403.05492v1 (2024),
Theorem~4, Example~5, Example~9, Proposition~10, Remark~11,
Lemma~13 and Proposition~16.
Published version: doi:10.1090/proc/16870.

\bibitem{BrunsConca2003}
W.~Bruns and A.~Conca,
\emph{Gr\"obner bases and determinantal ideals},
arXiv:math/0302058v3 (2003), Lemma~2.5 and Theorem~2.4.

\bibitem{HenriquesVarbaro2016}
I.~Henriques and M.~Varbaro,
\emph{Test, multiplier and invariant ideals},
Adv. Math. \textbf{287} (2016), 704--732,
doi:10.1016/j.aim.2015.09.028; arXiv:1407.4324.

\bibitem{Vainsencher1982}
I.~Vainsencher,
\emph{Schubert calculus for complete quadrics},
in Enumerative Geometry and Classical Algebraic Geometry (Nice, 1981),
Progress in Mathematics \textbf{24}, Birkh\"auser (1982), 199--235.

\bibitem{Massarenti2020}
A.~Massarenti,
\emph{On the birational geometry of spaces of complete forms I:
collineations and quadrics},
Proc. Lond. Math. Soc. \textbf{121} (2020), 1579--1618,
doi:10.1112/plms.12377; arXiv:1803.09161v2, \S2.

\bibitem{CasarottiCornianiMassarenti2023}
A. Casarotti, E. Corniani and A. Massarenti,
\emph{Complete singular collineations and quadrics},
International Mathematics Research Notices \textbf{2023} (2023),
no.~18, 15370--15407. DOI: 10.1093/imrn/rnac271.
The cited construction is also in arXiv:2111.02940,
Definition~2.5, Remark~2.6 and Theorem~2.14.

\bibitem{Thompson1991}
R. C. Thompson, \emph{Pencils of complex and real symmetric and skew
matrices}, Linear Algebra and its Applications \textbf{147} (1991),
323--371. DOI: 10.1016/0024-3795(91)90238-R.


\bibitem[Stacks Project]{StacksProject160}
The Stacks Project Authors, \emph{The Stacks Project},
Tags 0806 (universal property of blowing up), 02LS (proper quasi-finite
morphisms), and 0H1G (smooth centres and their blow-ups),
\url{https://stacks.math.columbia.edu}, accessed September 25, 2026.

\end{thebibliography}

\end{document}
